% ============================================================================
%  ASSUMPTIONS -- MAIN TEXT.  project `panel-npiv-discrete`
%
%  WHAT THIS FILE HOLDS.  The assumption STATEMENTS, and -- as of dispatch 25 --
%  one declaration block that has to precede them: the
%  numbered conditions the core results of the article cite, each copied
%  verbatim from assumptions.tex.  The interpretation paragraphs, the retirement
%  records, the worked feasibility tables and the "what this condition rules
%  out" discussion are in assumptions_supp.tex.
%
%  THE ONE EXCEPTION TO "STATEMENTS AND NOTHING ELSE", ADDED IN DISPATCH 25 AND
%  RECORDED HERE RATHER THAN LEFT AS A SILENT BREACH OF THIS FILE'S OWN
%  CONTRACT.  A \paragraph immediately under \section*{Assumptions} declares the
%  instrument's support, the propensity vector, the contrast vector, Omega, the
%  types of L / ell / b / ell^Pi, and the third factor of the range reference
%  measure.  IT COULD NOT GO IN assumptions_supp.tex.  That file compiles FIFTH
%  of the eight inputs, and A0 (first), B1 (second), A2' (third), A4 and
%  def:hatT all USE those symbols; a declaration landing after its first
%  consumer is the defect def:spaces records at its own 2026-08-03 note, and it
%  is invisible to the compiler because every reference still resolves.  The
%  block is a DECLARATION and not a condition -- the conditions it feeds are A0
%  (sans-L finite, fixed, known), A4(i) (the overlap floor) and B1(iv.d) (the
%  masses of Pi_w) -- so the "conditions only" rule of
%  .claude/references/theorem-environment-style.md is not breached either: no
%  restriction is smuggled into it.
%
%  assumptions_main.tex + assumptions_supp.tex ARE the old assumptions.tex.  The
%  split is a PARTITION of the printing text: every line that printed before
%  prints in exactly one of the two, no statement was rewritten, retitled or
%  renumbered, and no clause number moved.  Three whole conditions moved --
%  A3 and A6 (both RETIRED) and C4' (inert: no consumer anywhere in the corpus)
%  -- and they moved entire.
%
%%  DO NOT \input THIS FILE ALONGSIDE results.tex OR assumptions.tex.
%%  results.tex and assumptions.tex SURVIVE UNDIVIDED and remain the source of
%%  record: scripts/build-theory/build.sh compiles them, scripts/render-sync/
%%  check_rendering_sync.py reads them, and ~100 files under scripts/symbolic/
%%  specs/ cite them BY LINE NUMBER.  Neither was edited.  Every \label in this
%%  file is therefore declared TWICE on disk -- here and in its parent -- and the
%%  two must never reach the same document.
%
%  SUPPLEMENT POINTERS.  A statement kept here may cite a statement that lives
%  in the supplement.  Those calls are written \suppref{Lemma}{name} ->
%  "Lemma S:name" and \suppeq{name} -> "(S:name)", both defined in
%  paper_header.tex.  They are not \cref calls and cannot render "??".
% ============================================================================

% ============================================================================
%  CONSOLIDATED ASSUMPTION BLOCK -- project `panel-npiv-discrete`
%
%  Stage 3, dispatch 1 (proof_strategy_memo.md Sec.10, Steps 1-2).
%  Produced by the theorist, 2026-07-28.  ROUND 2: revised against
%  proofs_review.md (56/100).  Change log at the foot of this file.
%
%  CONTRACT.  Tags are STABLE: every later theorem cites by tag
%  (A0, A1, A1b, A1c, B1, A2, A2', H1, A3, A4, N1, A5, A8, A7, C1-C5),
%  never by restating.  The environment numbers may move; the tags may not.
%  FIVE THINGS ARE RETIRED, AND THEIR TAGS AND CLAUSE NUMBERS ARE KEPT RATHER
%  THAN DELETED, for the same reason the contract exists: several files cite
%  them, and a reader who follows one of those citations must land on the reason
%  and not on a gap.  A retired tag or clause is NOT AVAILABLE AS A HYPOTHESIS,
%  and no hypothesis list -- including a BLOCK citation of the parent tag -- may
%  carry one.  THE FIVE RETIREMENTS REST ON THREE DISTINCT GROUNDS AND THE
%  GROUND IS RECORDED AT EACH SITE:
%  [COUNT CORRECTED 2026-08-27, AND THE COMMENT WAS THE LAST HALF OF THE 08-05
%  SWEEP LEFT UNDONE.  A1d was retired 2026-08-05; the PRINTED text at the A1d
%  clause has said "five retirements and not four, on three grounds" since that
%  date, and assumptions_supp.tex's census was corrected with it.  THIS BLOCK --
%  the file's own register of the rule -- still said FOUR and enumerated FOUR.
%  No reader was misled, the comment being uncompiled; that is exactly the
%  condition under which this project's standing lesson says a rule rots, and it
%  is the third time this header's count has been repaired (it said TWO through
%  dispatch 17).  A COUNT IN A COMMENT MUST BE SWEPT WITH THE PRINT THAT CARRIES
%  IT.]
%    A6 (ass:tcc, dispatch 6)  RETIRED BECAUSE FALSE -- prp:tccfails proves its
%       event is null.  Nothing replaces it as a hypothesis; def:a6b is a
%       property no theorem consumes.
%    A3 (ass:rank, dispatch 13)  RETIRED BECAUSE PROVED -- lem:momzero(iv)
%       derives it from A1 and A1b(iv) plus the finiteness tags already on every
%       list that carried it.  Striking it WEAKENS four hypothesis lists and
%       changes no conclusion.
%    A2'(ii) (ass:sourcecondp, dispatch 14)  RETIRED BECAUSE PROVED, on the A3
%       ground -- it follows from A2'(i) in full with B1(i), (ii), (iv.a), (iv.d)
%       via cor:elldiff and (P1b) -- (P1b) and NOT (P1), the constant 2 being
%       c_Wtilde^2 and the bound Euclidean; (P1) is the PER-CELL |Wtilde_l| <= 1
%       and does not bound the vector norm.  Ground at assumptions_supp.tex's
%       retirement note and at lem:momzero's closing paragraph, which prints "the
%       bound on ||Wtilde|| is the Euclidean one and NOT (P1)"; cor:rangegap's
%       proof records that reading it off (P1) produced 2 bar-c_pi^{1/2}, FALSE at
%       sans-L >= 3.  [Tag corrected 2026-08-24.]  The clause number is kept as a struck
%       placeholder.  [(iv.d) ADDED 2026-08-24 with cor:elldiff's own line; the
%       RETIREMENT IS UNAFFECTED -- (iv.d) is in force wherever A2'(i) is -- but
%       a retirement ground is a derivation and must name every clause it uses.]
%    A1d (dispatch 25, 2026-08-05)  RETIRED BECAUSE FALSE, on A6's ground and not
%       on a new one -- its hypothesis set is EMPTY.  A1 makes Wtilde
%       conditionally centred given (X_1, X_2), so the clause's own displayed
%       event fails at Lambda == 1, which lies in the class the clause quantifies
%       over.  The clause carries a printed struck banner at its own site and the
%       printed sentence there says "five retirements and not four, on three
%       grounds"; this entry is the comment half of that sweep.
%    B1(vi) (ass:bierens, dispatch 14)  RETIRED BECAUSE IT IS ABOUT A DIFFERENT
%       ESTIMATOR -- the THIRD ground, and it is neither of the other two.  The
%       growing finite family T_n occurs nowhere outside its own sentence, while
%       def:hatT integrates against Pi_t over ALL of T.  What its heading claimed
%       -- identification of varphi_dagger -- is A1b(iii) here, assumed outright.
%  REGISTER RULE, ADDED IN DISPATCH 18 BECAUSE B1(vi)'s RETIREMENT WAS HALF DONE
%  FOR FOUR DISPATCHES: a retirement is complete only when the clause is struck,
%  when no hypothesis list carries it -- INCLUDING VIA A BLOCK CITATION OF THE
%  PARENT TAG -- and when nothing anywhere describes it in the present tense as
%  a live condition.  B1(vi) failed the second test at prp:phirate and the third
%  at three compiled sites; both are closed in that dispatch.
%  THE RULE HAS A SECOND HALF, RULED BY THE AUTHOR 2026-08-24 AND STATED HERE SO
%  THAT NO LATER PASS RE-DERIVES IT: *** A DEFINITION'S BLOCK CITATION IS A
%  LOCATING POINTER AND CARRIES NO HYPOTHESIS -- def:spaces's bare C1 and B1(iv)
%  LOCATE the objects they name and impose nothing -- WHILE A BLOCK CITATION ON A
%  RESULT'S HYPOTHESIS LINE CARRIES EVERY CLAUSE UNDER THE TAG. ***
%    AMENDED 2026-08-26: THE POINTER HALF WAS STATED AS A UNIVERSAL AND IS FALSE
%    AT TWO SITES.  A2b IS DEFINED AS "admits a version satisfying A2'(i.a)-(i.b)"
%    -- the citation is its ENTIRE CONTENT -- and def:hatT builds eq:trunset out of
%    A4(i)'s constant.  A CONSTITUTIVE citation is not a pointer.  The printed rule
%    now scopes the pointer half ("a pointer WHEN IT LOCATES") and gives the
%    discriminant: ask what survives deleting the citation.
%  *** AND A THIRD REGISTER, RULED BY THE AUTHOR 2026-08-24 (SECOND PASS).  A
%  CITATION OF AN ASSUMPTION'S *DECLARATIVE HALF* -- the part that INTRODUCES a
%  symbol, a type, a carrier set or a construction -- IS A **DECLARATION** AND
%  CARRIES NO HYPOTHESIS, because a declarative half is STANDING: in force at every
%  statement in this corpus, named or not, since without it the statement has no
%  objects to be about.
%  H1's H^s scale, B1(iii)'s T subset R^p and B1(iv)'s STEM
%  Pi = Pi_u (x) Pi_t (x) Pi_w are declarations; B1(iv.a)-(iv.d), B1(iii)'s
%  Leb(T) > 0, H1's rank/decay/eq:h1match and A8's basis clause are clauses. ***
%  *** THE DISCRIMINANT WAS THE NEGATION TEST -- "statements stop parsing ->
%  DECLARATION" -- AND IT IS SUPERSEDED 2026-08-26, BECAUSE THE CORPUS REFUTES IT.
%  A4(i) was EXISTENTIAL, so negating it left eps_varrho with no referent and
%  stopped eq:trunset, hat-Wtilde, eq:hatT and the seven results built on them from
%  PARSING -- so the test ruled OVERLAP a declaration carrying no hypothesis, while
%  this block carries it as a condition at seven sites, correctly.  The test also
%  turned on an unstated reading convention: ungenerous at A0 it returned
%  "declaration", generous at A8's basis clause -- which ALSO introduces the b_j
%  that def:lfstep needs -- it returned "carrier", and the two readings stood three
%  lines apart in the printed text.
%  THE REPLACEMENT, PRINTED: *** SPLIT FIRST, THEN TEST THE REMAINDER.  Cut the
%  fragment where it stops INTRODUCING and starts RESTRICTING; what introduces is
%  the declaration, and the remainder is a CLAUSE iff there is a spelling of this
%  block, with the declaration in force, under which every symbol used anywhere in
%  the corpus still has a referent and the remainder is FALSE. ***  B1(iii) was
%  already written this way and is the model.  A0, A4(i) and A8 were NOT split and
%  now are; the CONSEQUENCE is that A0's compact interval, A4(i)'s overlap
%  inequality and A8's basis requirement are all CLAUSES and are carried by the
%  results that read them.  A0's interval accordingly moves ONTO lem:intker's
%  hypothesis line, and lem:levels -- which does not read it at all, see the
%  intersection-form repair there -- carries A0's COMPACTNESS and not its interval.
%  *** ALL THREE REGISTERS ARE NOW **PRINTED**, in the \paragraph "How a citation
%  of this block is read" under the notation block below.  THAT MATTERS AND IS WHY
%  THE THIRD WAS NEEDED AT ALL: through this date BOTH halves above lived only in
%  comments -- this one, and for the carrier half the ARCHAEOLOGY comment under
%  lem:opnoise in results_supp.tex ("a block citation of the parent tag is a
%  citation of every clause under it"), cited by ANCHOR TEXT and not by line
%  because this same pass moves that file's line numbering --
%  which no compiled document renders, while THREE PRINTED instance-defences
%  disagreed in rhetoric about what a B1(iv) block citation meant.  THIS FILE'S OWN
%  PRECEDENT FOR AN UNPRINTED CONVENTION IS SIXTY LINES BELOW: the \ref-to-an-
%  unnumbered-assumption rule accumulated 82 BAD CALLS ACROSS FOUR FILES while it
%  was true and unwritten.  A RULING THAT LIVES ONLY IN A COMMENT IS NOT A RULING.
%  This comment is retained as the record of the ruling, NOT as its home. ***
%  THE CORPUS ALREADY ACTED ON THIS EDITORIALLY AND ONLY THE SENTENCE WAS MISSING:
%  prp:hatAprop lists (iv.a)-(iv.d) explicitly, which would be redundant if
%  def:spaces's pointer had already imposed them, and lem:aerr prints a defence of
%  why B1(iv.a) sits on its own line.
%  TWO CONSEQUENCES, BOTH LIVE, AND THEY RUN IN OPPOSITE DIRECTIONS.
%   (1) Eight results cited B1(iv.d) at a named proof step while their hypothesis
%       lines excluded it.  Those are GENUINE DEFECTS and were repaired 2026-08-24:
%       prp:wellposed, prp:frechet, prp:compact, lem:adjfubini, lem:opnoise,
%       lem:unifop and lem:aerr were promoted to the block B1(iv) -- six of the
%       seven printed the IDENTICAL line "B1(i), (ii) and (iv.c)", so it was one
%       line copied and not six oversights -- and cor:elldiff, whose proof reads
%       neither (iv.b) nor (iv.c), was promoted BY CLAUSE to "(iv.a) and (iv.d)".
%       THAT REPAIR OVERSHOT AT SIX OF THE SEVEN AND IS ITSELF REPAIRED 2026-08-24
%       (SECOND PASS), on the THREE REGISTERS rule now PRINTED under the notation
%       block above.  The block carries (iv.b) at all six and (iv.a) at five, and
%       NO PROOF OF ANY OF THE SIX OPENS (iv.b).  The six now cite BY CLAUSE:
%         prp:wellposed, prp:frechet, prp:compact, lem:opnoise, lem:unifop
%                                  -> B1(i), (ii), (iv.c), (iv.d)
%         lem:aerr                 -> B1(i), (ii), (iv.a), (iv.c), (iv.d)
%       lem:aerr's (iv.c) is NOT a fourth-clause flourish: proofs.tex Step 1 and
%       Step 2 both integrate bar-w^2 against Pi_t and the stated O((n h_x h_u)^-1)
%       needs that integral finite.  The clause was CONSUMED at two steps and NAMED
%       at neither through 2026-08-23.  THE 2026-08-24 PASS REPAIRED STEP 1 ONLY,
%       and both this entry and lem:aerr's own statement then said the repair was
%       done -- a closure claim short by one, at the very site whose subject is a
%       clause named at no step.  STEP 2 NAMES (iv.c) FROM 2026-08-26; both now do.
%       WHAT NO LINE CARRIES: the FACTORISATION Pi = Pi_u (x) Pi_t (x) Pi_w, which
%       all six need to write their iterated integrals.  It is B1(iv)'s
%       DECLARATION and is STANDING -- in force at every statement, cited by none.
%       lem:adjfubini KEEPS the block: its proof opens all four clauses, naming
%       (iv.c) and (iv.d) BY CLAUSE at the c_Pi equality and the probability-mass
%       factor, and (iv.a)/(iv.b) THROUGH THE BLOCK at the "product of THREE finite
%       measures" sentence, the clause attribution being printed at the STATEMENT.
%       IT IS NOT THE ONLY SUCH LINE.  This entry said it was through 2026-08-25,
%       and so did the printed sentence at lem:adjfubini; both are now retyped.
%       cor:elldiff's by-clause line is UNCHANGED and is now an instance of the
%       printed rule rather than its own local defence.
%       prp:stopfeas inherits lem:opnoise's line by reference and is repaired by
%       that one edit, in this pass as in the first.
%       THE TEN FURTHER B1(iv) BLOCK-FORM LINES ARE NOW AUDITED PROOF BY PROOF AND
%       ALL TEN STAND; NO HYPOTHESIS LINE MOVES.  An eleventh was found in the same
%       sweep -- thm:rootn's group (T-3) -- and stands on the same ground.  By
%       ground:
%         A. OPENS ALL FOUR AT NAMED STEPS OF ITS OWN PROOF
%              lem:adjfubini, lem:adjoint, prp:tccfails
%         B. OPENS SIX OF THE SEVEN FRAGMENTS AT NAMED STEPS; (iv.c) IS CONSUMED BY
%            THE STATEMENT AND NOT BY A STEP.  lem:bierens clause (d) and
%            prp:kerchar clause (a) are statements about the L^2(Pi)-VALUED maps of
%            eq:Tdef and eq:Adef, whose typing is prp:wellposed / prp:compact under
%            (iv.c).  THIS IS THE lem:aerr SPECIES -- consumed at the statement,
%            named nowhere -- and a tag census undercounts both by exactly one.
%              lem:bierens, prp:kerchar
%         C. CARRIES IT BY INHERITANCE, THE PROOF NAMING NO B1 CLAUSE WHATEVER.  A
%            census over B1 tags returns ZERO at each of these and would report the
%            block unused; every edge is a \cref edge.  This is the cor:centarity
%            failure mode read in reverse.
%              cor:opsol    re-runs lem:bierens's proof at arbitrary (varphi,theta)
%              prp:localise cor:opsol; its own proof says "the only place in the
%                           proof at which B1 is used"
%              lem:solset   lem:adjoint(ii), prp:orthog, cor:elldiff
%              lem:biasexp  lem:adjoint(iii), prp:orthog, lem:momzero, cor:elldiff
%              prp:orthog   cor:elldiff for (iv.a),(iv.d); lem:adjoint(iii) for the
%                           block, cited in terms at its Part (i)
%       NO LINE PINS IN THIS ENTRY, DELIBERATELY.  The ten were recorded here at
%       results_supp.tex :609/:671/:1190/:2960/:5102 and
%       results_s2_identification.tex :562/:648/:794/:943/:2057, and EVERY ONE OF
%       THOSE TEN PINS HAD ROTTED before the audit ran -- by +9 in results_supp and
%       by up to +130 in results_s2.  Cite by tag, never by line.
%   (2) The claims that C1(i) is a hypothesis of ONE lemma (lem:aerr) and,
%       transitively, of NO THEOREM -- results_supp.tex under lem:aerr and
%       assumptions_supp.tex's five-clause reading -- STAND and are not weakened.
%       Nothing about a definition's pointer creates a C1(i) edge.
%  THIS HEADER READ "THE TWO RETIREMENTS" THROUGH DISPATCH 17, WHICH HAD BEEN
%  WRONG SINCE DISPATCH 14 ADDED THE OTHER TWO.
%  DISPATCH 7 also adds A5(vi') and moves def:intker and def:hatT into Part V,
%  because A7 is a statement ABOUT T-hat and a maintained assumption must not
%  name an undefined map.
%
%  HOW RESULTS CITE THESE, AND THE ONE THING YOU MUST NOT DO.
%  Each assumption's heading carries its tag, e.g. "Assumption (A2: Source
%  condition for theta)".  Theorems and lemmas cite the TAG in running text
%  ("Suppose B1(i)-(iv) hold").  The tag is the ONLY identifier.
%
%  *** NEVER WRITE \ref{ass:...} OR \cref{ass:...}. ***
%  The environment is declared \newtheorem*{assumption}{Assumption} -- with the
%  STAR, hence UNNUMBERED -- at preamble.tex:283 and in both rendered preambles.
%  A \ref to an unnumbered environment has no counter to read and silently
%  prints whatever counter last incremented: NOTHING AT ALL in the build.sh
%  corpus ("with s as in ."), and the SECTION NUMBER in the article
%  ("with s as in 2").  It is not an undefined reference -- the label exists --
%  so build.sh, render.sh and check_rendering_sync.py ALL report clean.  Only
%  reading the PDF finds it.
%
%  THROUGH 2026-08-01 THIS HEADER SAID THE OPPOSITE: that the environment was
%  \newtheorem{assumption}{Assumption}[section] and that "\cref to the
%  environment label returns the printed NUMBER".  Both halves were false, and
%  they are why 82 such calls accumulated across four files -- an author reading
%  this header would reasonably believe \ref worked.  build.sh now greps for the
%  pattern; see its "assumption refs" finding.
%
%  Requires the macros \IR, \un, \upmodels, \vdagger from preamble.tex.
%
%  Symbols follow quality_reports/theory/panel-npiv-discrete/notation_glossary.md.
%  In particular:  gamma  (not mu) is the Bierens signed measure;
%                  mu     is the source-condition index;
%                  p_K    is the kernel order;  p = dim(Z);
%                  d_Z    is the EFFECTIVE dimension of Delta Z, d_Z <= p-1;
%                  varsigma, s_0 are RESERVED LOCAL BOUND VARIABLES with no
%                           meaning across results -- and varsigma is NOT the
%                           Step-4 direction, which is vartheta (round-7 rename);
%                  vartheta in R^p is the coordinate-contraction direction,
%                           BINDING ON STEP 4;
%                  tau_DP is A7's discrepancy constant; tau = (tau_1,tau_2,tau_3)
%                           is Part F's perturbation parameter.
%
%  ADDED IN DISPATCH 25 (the multinomial instrument), and declared in the
%  notation block immediately below \section*{Assumptions}:
%                  sans-L  is the NUMBER OF INSTRUMENT CATEGORIES, an integer
%                           >= 2, fixed and known.  ITALIC L IS UNTOUCHED and
%                           remains the debiasing weight L(u,z); L_x, L_u remain
%                           the sieve dimensions and L^2(.) the function spaces.
%                           THE SANS/ITALIC PAIR IS THIS PROJECT'S OWN DEVICE,
%                           already run at sans-C_{h_u} against C_{h_u} and at
%                           sans-r against italic r;
%                  varrho  is now the VECTOR (varrho_1, ..., varrho_{sans-L - 1})'
%                           of conditional cell probabilities, with
%                           varrho_l(z) = P(W = l | Delta Z = z) and
%                           varrho_0 = 1 - sum_{l >= 1} varrho_l.  At sans-L = 2
%                           the vector has one entry and IS the scalar propensity
%                           of every earlier round;
%                  Wtilde  is the (sans-L - 1)-VECTOR of contrasts, eq:contrast;
%                  Omega   is the conditional covariance of Wtilde, eq:omegadef.
%                           IT IS NOT A WEIGHT.  It appears in exactly three
%                           roles: as that covariance; as the PREFACTOR of
%                           eq:omegaprefac, E[Wtilde 1{DU<=u}|DZ] = Omega d; and
%                           as the object whose invertibility the identification
%                           discussion turns on.  Omega^{-1} APPEARS IN NO
%                           ESTIMATING EQUATION -- neither eq:hatT nor eq:psidef
%                           carries it -- BUT IT DOES APPEAR IN DERIVATIONS: at
%                           eq:omegainv, in the ||Omega^{-1}|| divergence of the
%                           Khan-Tamer passage under A4(d), in ell_star = ell
%                           Omega^{-1}, and throughout rem:htread.  (Round 4 wrote
%                           "appears ONLY inside rem:htread", which is false; the
%                           true statement is about the ESTIMATOR, not about the
%                           file.)
%                  Wtilde^Omega := Omega^{-1} Wtilde, eq:omeganorm.  A READING,
%                           NOT AN OBJECT THE ESTIMATOR USES.  THE ESTIMATOR
%                           CARRIES THE RAW CONTRAST ON BOTH SIDES: the operator
%                           at eq:hatT and the moment at eq:psidef.  Normalising
%                           was considered in dispatch 25 and REJECTED; the
%                           ground is at rem:htread and in the change log.
%                           *** THE CASE PAIR IS CLOSED, 2026-08-09.  Omega WAS
%                           a declared case pair: Omega(varphi, varphi-tilde)
%                           also named the intertwining family of prp:norangeinv
%                           (9 sites, in results_supp.tex and proofs.tex),
%                           separated BY TYPE ALONE -- one an operator family on
%                           L^2(Pi), the other a real symmetric
%                           (sans-L - 1) x (sans-L - 1) matrix.  (A first draft
%                           of this line also claimed ARITY and that was false:
%                           the intertwining family was bound to BARE Omega and
%                           a bare Omega^{-1} was printed -- arity zero, the exact
%                           string the normalisation reading owns.)
%                           THE RENAME IS DONE: the family is cal-O, declared in
%                           prp:norangeinv; results_supp.tex was swept in Phase 4
%                           round 1 and the last five proofs.tex tokens in round
%                           2.  eq:omegadef's contrast covariance is now the ONLY
%                           Omega in the corpus and needs no disambiguation rule.
%                           The nine printed lines that recorded the pair and
%                           gave the reader a disambiguation rule ARE DELETED
%                           from the type-declaration block below; see the
%                           archaeology block at that site. ***
%                  Pi_w    is the third factor of the range reference measure, a
%                           PROBABILITY measure on {1, ..., sans-L - 1} with
%                           masses varpi_w(l) > 0.  B1(iv.d).  At sans-L = 2 it
%                           is the point mass at 1 and Pi = Pi_u (x) Pi_t
%                           verbatim, so c_Pi of B1(iv.c) does not move.
% ============================================================================

\section*{Assumptions}

% ===========================================================================
%  NOTATION BLOCK -- NEW IN DISPATCH 25 (2026-08-08).
%  Phase 1 of quality_reports/plans/2026-08-08_multinomial-instrument.md.
%
%  WHAT CHANGED AND WHY IT IS ONLY A REPRESENTATION.  Through dispatch 24 the
%  instrument was BINARY and the restriction A1 was carried by a SINGLE SCALAR
%  contrast Wtilde = W - varrho(Delta Z).  A1's STATEMENT never said "binary":
%  it says Delta U independent of W given Delta Z, which is a statement about
%  sigma-fields and is indifferent to the cardinality of W's support.  WHAT WAS
%  BINARY WAS THE REPRESENTATION -- one contrast exhausts the restriction only
%  when W has two cells.  For W in {0, ..., sans-L - 1} the restriction is
%  carried by the (sans-L - 1)-VECTOR of contrasts declared below.  THIS IS WHY
%  THE GENERALISATION TOUCHES REPRESENTATION ONLY: no assumption's content
%  changes, and A1's statement is reproduced verbatim.
%
%  NOTHING IS APPROXIMATED IN THE W DIRECTION, AND THIS IS THE REASON THE
%  GENERALISATION IS CLEAN.  The conditionally centred functions of (W, Delta Z)
%  form, at each Delta Z, a space of dimension EXACTLY sans-L - 1, and the
%  sans-L - 1 dummy contrasts are a basis for it (derivation in
%  assumptions_supp.tex, paragraph "A1").  So no Bierens-type device is needed
%  on the W side; the Bierens family stays exactly where it was, over Delta Z,
%  and B1 is untouched except for the range measure's third factor.
%
%  *** WHAT THIS BLOCK MUST NOT BE READ AS SAYING. ***
%  The cardinality of W does NOT act through the conditioning of a linear
%  ill-posed inverse problem, and no sentence anywhere in this dispatch may be
%  written as though it did.  varphi solves the NONLINEAR operator equation
%  eq:opeq under a CONDITIONAL INDEPENDENCE restriction; def:operator's T is not
%  a linear Fredholm operator in varphi and A_cond is its DERIVATIVE, not the
%  problem.  Importing "more instrument values improve the conditioning of the
%  inverse problem" would import a different problem.  What extra cells buy is
%  stated below and is a different thing entirely.
%
%  IDENTIFICATION IS NOT AT STAKE AND THE BLOCK MUST NOT SUGGEST IT IS.
%  Identification already holds at sans-L = 2 under A1b(iii).  Moreover A1b(iii)
%  at the FULL sans-L-category W is IMPLIED BY A1b(iii) at any coarsening of W:
%  a coarsening imposes fewer restrictions, so its V_CI is LARGER, and
%  V_CI = {varphi_dagger} at the coarsening gives the same at the full W.  A
%  finer instrument therefore makes the identification hypothesis WEAKER, never
%  stronger.  ~~The extra contrasts are SURPLUS relative to the binary design:
%  they buy efficiency and, at sans-L >= 3, OVERIDENTIFICATION -- sans-L - 2
%  testable restrictions, which is what converts A1 from an untestable
%  assumption into one with evidence for or against it.~~  STRUCK IN ROUND 5.
%  A1 IS TESTABLE AT sans-L = 2: the restriction is d_1(u|z) = 0 for every u at
%  a.e. z, a CONTINUUM of equations against one unknown function, because this is
%  a conditional-INDEPENDENCE restriction with a continuously distributed
%  conditioning variable and not the mean-independence case in which a binary
%  instrument supplies two moments.  Section 6's own probe rejects A1 at
%  sans-L = 2.  The "sans-L - 2 restrictions" count was a moment-counting reflex
%  from the wrong case.
%  WHAT THE EXTRA CELLS BUY IS **LOCALISATION**, and three concepts must be kept
%  apart: TESTABILITY (present at sans-L = 2, not what sans-L > 2 adds);
%  INSTRUMENT STRENGTH (a property of the data -- splitting creates no variation
%  and strength can be examined at sans-L = 2 by subsetting -- also not what
%  sans-L > 2 adds); and LOCALISATION, which is the claim.  At sans-L = 2 the
%  restriction is one object, so a violation in one channel and its absence in
%  another average into a single verdict; at sans-L >= 3 the restriction is
%  d_l(u|z) = 0 PER CELL and a violation can be ATTRIBUTED to a cell.  That
%  question cannot be posed at sans-L = 2.  The joint-consistency / J-type
%  statement is the APPARATUS; localisation is the MOTIVATION.  Set out in full
%  in the Discussion paragraph A1, with the Section 6 chronic-versus-transient
%  split it exists to test.
%  NO SENTENCE IN THIS CORPUS MAY SAY THAT MORE CATEGORIES ARE NEEDED FOR
%  IDENTIFICATION, OR THAT THEY MAKE A1 TESTABLE.
%
%  WHY THE COUNT IS sans-L AND NOT ITALIC L.  Italic L is the debiasing weight
%  L(u,z) (glossary row L, ~200 sites); L_x and L_u are the sieve dimensions;
%  L^2(.) are the function spaces.  "L - 1" in italic would read as the weight
%  minus one.  THE PLAN AT quality_reports/plans/2026-08-08_multinomial-
%  instrument.md WRITES ITALIC L AND THIS BLOCK DEPARTS FROM IT DELIBERATELY,
%  on the project's own established sans/italic device -- sans-C_{h_u} against
%  C_{h_u} (def:intker against C3), sans-r against italic r (the varrho
%  direction against A4(ii)'s Holder order).  Recorded rather than done
%  silently, because the plan is the contract and a departure from a contract
%  must be visible.
%  ONE COLLISION SURVIVES AND IS RECORDED: sans-L_Gamma is the Laplace transform
%  in lem:bierenscore's proof (proofs.tex, 8 sites, one proof).  It ALWAYS
%  carries the subscript Gamma and is a proof-local object; the category count
%  NEVER carries a subscript.  The cheap repair, flagged for the author and not
%  taken here because it is in a Phase-3 file, is to rename that transform.
% ===========================================================================

\paragraph{The instrument, its contrasts, and the range reference measure.}
The instrument $W$ takes finitely many values, $W\in\{0,1,\dots,\mathsf L-1\}$, with
$\mathsf L\ge2$ an integer that is \emph{fixed and known to the analyst}. \emph{Sans-serif
$\mathsf L$ is the number of instrument cells and nothing else}: italic $L$ remains the debiasing
weight of \cref{def:weight}, $L_x$ and $L_u$ remain the sieve dimensions, and $L^{2}(\cdot)$ remain the
function spaces. Write, for $l\in\{0,1,\dots,\mathsf L-1\}$,
$\varrho_l(z):=P(W=l\mid\Delta Z=z)$, and collect the \emph{non-base} cells into
\begin{equation}\label{eq:contrast}
\varrho\;:=\;\big(\varrho_1,\dots,\varrho_{\mathsf L-1}\big)',
\qquad
\widetilde W_l\;:=\;\un\{W=l\}-\varrho_l(\Delta Z),
\qquad
\widetilde W\;:=\;\big(\widetilde W_1,\dots,\widetilde W_{\mathsf L-1}\big)'\in\IR^{\mathsf L-1},
\end{equation}
so that $\varrho_0=1-\mathbf 1'\varrho$ with $\mathbf 1$ the vector of ones in $\IR^{\mathsf L-1}$.
Cell $0$ is the base cell. The \emph{conditional contrast covariance} is
\begin{equation}\label{eq:omegadef}
\Omega(z)\;:=\;\operatorname{Var}\big(\widetilde W\mid\Delta Z=z\big)\;=\;\operatorname{diag}\big(\varrho(z)\big)-\varrho(z)\varrho(z)' ,
\qquad z\in\IR^{p} .
\end{equation}

At $\mathsf L=2$ these are the scalar objects: $\varrho=\varrho_1=P(W=1\mid\Delta Z)$ is the propensity,
$\widetilde W=W-\varrho(\Delta Z)$, and $\Omega=\varrho(1-\varrho)$.

The contrasts span, and nothing in the $W$ direction is approximated: every measurable
$a(W,\Delta Z)$ with $E[|a|\mid\Delta Z]<\infty$ and $E[a\mid\Delta Z]=0$ almost surely is
$a=c(\Delta Z)'\widetilde W$ for a $\Delta Z$-measurable $c$, uniquely so at a $\Delta Z=z$ at which
every cell is non-null, where the space has dimension exactly $\mathsf L-1$ and \eqref{eq:contrast} is
a basis for it. \emph{The spanning statement needs no overlap condition}; at a $z$ with null cells the
dimension is one less than the number of non-null cells. Consequently the Bierens device is needed only
in the $\Delta Z$ direction, where it always was, and B1 is unchanged apart from B1(iv.d). The
derivation, the exact statement at a degenerate $z$, and the fact that neither depends on which cell is
taken as the base are in the Discussion paragraph \emph{A1}.

Type declaration, stated once and binding on every statement in this corpus. The debiasing
weight and the objects built from it carry one column per contrast: $L$, $\ell=\partial_uL$, the
anchor $b$ and their starred and hatted variants are $\IR^{p\times(\mathsf L-1)}$-valued, one
\emph{row} per coordinate of $\theta$ and one \emph{column} per contrast, and the range-space
representer $\ell^\Pi$ carries the contrast index as a third argument,
$\ell^\Pi(u,t,l)$. \emph{Norms carried by these objects are Frobenius}; matrix norms elsewhere in this
corpus are spectral unless the site says otherwise, and the two are named at every site at which both
occur. The moment function is
$\psi=\Delta Z\,\Delta U+L(\Delta U,\Delta Z)\,\widetilde W$ of \eqref{eq:psidef}, the contrast standing
to the \emph{right} of the $p\times(\mathsf L-1)$ weight. \emph{At $\mathsf L=2$ each of these is the
printed object}: a
$p\times1$ matrix is a vector in $\IR^{p}$, the Frobenius norm is the Euclidean norm,
$\ell^\Pi(u,t,1)=\ell^\Pi(u,t)$, and $L\widetilde W=\widetilde WL$ because $\widetilde W$ is then a
scalar.

The estimator carries the \textbf{raw} contrast \eqref{eq:contrast}, on both sides: the operator
equation is written at $\widetilde W$ and so is $\psi$, while $\Omega$ enters as the conditional
covariance \eqref{eq:omegadef} and as the prefactor \suppeq{omegaprefac}, and \emph{never as a weight}.
The normalised contrast $\widetilde W^{\Omega}:=\Omega(\Delta Z)^{-1}\widetilde W$, what it makes
explicit, and the reason the paper does \textbf{not} estimate at it are \suppref{Remark}{htread}.

% ===========================================================================
%  THE "DECLARED CASE PAIR" PARAGRAPH IS DELETED HERE, 2026-08-09, PHASE 4
%  CLOSE-OUT, BY AUTHOR RULING.  NINE PRINTED LINES STOOD HERE SAYING:
%    "Omega is a declared case pair.  The symbol Omega ALSO names the
%     intertwining family of prp:norangeinv ... the intertwining family is
%     RENAMED WHEN ITS OWN FILE IS NEXT OPENED, AND UNTIL THEN [a live
%     disambiguation rule for the reader]."
%
%  THAT RENAME IS DONE.  results_supp.tex:2954 now prints "The intertwining
%  family is written cal-O, and neither Omega nor cal-R", and the five
%  surviving proofs.tex tokens were swept in Phase 4 round 2.  So the shipped
%  document printed BOTH "closed" (in the supplement) AND "pending" (here),
%  and instructed the reader to disambiguate a collision that no longer exists.
%  That is this project's dominant defect species -- a claim about the artifact
%  that the artifact contradicts -- and it was in PRINTED text.
%
%  WHAT REPLACES IT: NOTHING.  eq:omegadef's conditional contrast covariance is
%  the only Omega left in the corpus and needs no disambiguation; cal-O stands
%  on its own declaration inside prp:norangeinv.  No mathematics moved and no
%  \label was removed.
%
%  DOWNSTREAM, AND IT HAD TO MOVE WITH THIS.  results_supp.tex:2961-2963 cited
%  this paragraph as the thing it discharges ("the case pair the type
%  declaration of eq:contrast recorded ... as that declaration said would
%  happen").  Deleting here without retyping there would have left a printed
%  citation of a retracted claim -- the same contradiction, relocated.  That
%  sentence is retyped in the same pass.
%
%  PRECEDENT FOR REOPENING A CLOSED FILE: Phase 4's own A2'/A2b/eq:lambdamu
%  reopening, same phase, same ground.
% ===========================================================================
The range reference measure gains a third factor: $\Pi=\Pi_u\otimes\Pi_t\otimes\Pi_{\mathsf w}$ with
$\Pi_{\mathsf w}$ a probability measure on $\{1,\dots,\mathsf L-1\}$, masses $\varpi_{\mathsf w}(l)$,
required strictly positive by B1(iv.d). \emph{At $\mathsf L=2$, $\Pi_{\mathsf w}$ is the point mass at
$1$ and $\Pi=\Pi_u\otimes\Pi_t$}, under the canonical identification of $\IR\times\mathcal T\times\{1\}$
with $\IR\times\mathcal T$; $c_\Pi$ of B1(iv.c) is unchanged at every $\mathsf L$, $\Pi_{\mathsf w}$
being a probability measure and contributing the factor $1$.

% ============================================================================
%  SHORTENED 2026-08-28.  ~1,700 PRINTED WORDS OF REGISTER TAXONOMY MOVED TO
%  assumptions_supp.tex, "The three registers: how a citation of the assumption
%  block is read".  What stays is the RULE; what left is the argument for it --
%  the split-first test, the three analyst-constant registers, the refutation of
%  the superseded negation test, and the two positions the taxonomy leaves out.
%  THE RULE HAD TO STAY: A0, A4(i), A8, B1(iii) and B1(iv) are written in the
%  split form and read unintelligibly without it, and prp:frechet, prp:compact
%  and five sites in results_supp.tex cite "the three registers" by name.
% ============================================================================
\paragraph{How a citation of this block is read.}
\emph{On a result's hypothesis line} --- inside a ``Suppose \dots\ hold'' --- a citation is a
\textbf{carrier}: naming a tag carries every condition under it, and naming a clause carries that clause
and no other. \emph{Inside a definition} a citation is a \textbf{pointer} when it merely locates the
object it names, as \cref{def:spaces}'s ``$\Pi=\Pi_u\otimes\Pi_t\otimes\Pi_{\mathsf w}$ as in
Assumption~B1(iv)'' does, and imposes nothing; it is \textbf{not} a pointer when it is constitutive, as
A2b's ``admits a version satisfying A2$'$(i.a)--(i.b)'' and \cref{def:hatT}'s use of the constant A4(i)
declares are. \emph{And several conditions below are written in two halves}: a \textbf{declarative}
half, which introduces a symbol, a type or a carrier set and is \emph{standing} --- in force at every
statement in this corpus, named or not, since without it the statement has no objects to be about, so
it is never listed as a hypothesis --- and a \textbf{clause}, which restricts, and which is carried on
the hypothesis line of every result whose proof reads it. A0, A4(i), A8, B1(iii) and B1(iv) are split
that way. Where the split falls, why it replaced an earlier test, and the register of the three
constants the analyst chooses are in the Discussion of the assumptions, under the heading \emph{The
three registers}.

% WHY Pi_w IS A PROBABILITY MEASURE AND NOT A GENERAL FINITE ONE.  Two reasons,
% and the first is arithmetic.  (1) c_Pi^2 = Pi_u(R) int bar-w^2 dPi_t is
% multiplied by Pi_w({1,...,sans-L-1}) once the third factor is present; fixing
% that mass at 1 leaves B1(iv.c) LITERALLY UNCHANGED and leaves every constant
% that carries c_Pi -- prp:frechet's operator-norm bound, prp:compact's HS bound,
% C_A -- unchanged as well.  A general finite measure would inflate all of them
% by a factor no result tracks.  (2) It makes the sans-L = 2 reduction exact
% rather than up to a constant: Pi_w = delta_1 is the only probability measure on
% a one-point set.
%
% *** AND Pi_w IS NOT THE NORMALISATION Omega^{-1}.  DO NOT CONFLATE THEM. ***
% Pi_w weights the contrasts INSIDE THE L^2(Pi) NORM of the operator equation --
% it is an analyst's choice on the range side, exactly as Pi_t is, and it is
% priced by B1(iv) with the other reference measures.  Omega^{-1} sits INSIDE the
% operator's integrand, pointwise in Delta Z, and is determined by the
% data-generating process through eq:omegadef.  They sit in different places, are
% chosen by different agents, and answer different questions.  A dispatch that
% sets Pi_w = "the efficient weights" has made a category error.
%
% ===========================================================================
%  NORMALISATION BY Omega^{-1}: CONSIDERED, AND REJECTED BY THE AUTHOR.
%  THE ESTIMATOR CARRIES THE RAW CONTRAST Wtilde ON BOTH SIDES.  DO NOT RE-OPEN.
%
%  THE HISTORY IS KEPT BECAUSE THE REJECTED ROUTE IS THE ONE A LATER READER WILL
%  RE-INVENT.  Dispatch 25 tried three answers to "where does Omega^{-1} go" and
%  the author rejected all three: round 1 declared it and placed it nowhere;
%  round 2 placed it on the operator and argued redundancy on psi; round 3
%  normalised both sides.  THE REVERSAL IS NOT A RETREAT TO ROUND 1: rounds 2
%  and 3 established real facts, and those facts are what justify the rejection.
%
%  *** THE AUTHOR'S GROUND, AND IT IS DECISIVE. ***  NORMALISING MAKES A_cond A
%  DIFFERENT OPERATOR.  A2's Picard/range condition, A2', A1c and H1 -- the last
%  through sigma_j(A_cond) asymp j^{-a} -- are ALL STATED AT A_cond IN THIS FILE.
%  Downstream of the decay exponent a sit mu = s/(2a), the frontier eq:frontier,
%  the whole feasibility table in the Discussion, and the CFF ceiling comparison.
%  Re-stating that chain at a normalised operator is not free, and THIS FILE
%  CONTAINS ITS OWN PROOF THAT IT IS NOT: A4's note (c) says in terms that the
%  two operators MUST NOT BE ASSUMED to share range and kernel, and the round-3
%  scoping of the range-invariance sentence (see (d) below) restricts range
%  invariance to CONSTANT varrho.  A z-varying Omega is exactly the case that
%  scoping excludes.  THE PRICE IS THE ENTIRE SMOOTHNESS-AND-ILL-POSEDNESS
%  BUDGET; THE GAIN WAS INTERPRETATIONAL.  The author declines the trade.
%
%  WHAT SURVIVES THE REJECTION, AND IT IS THE POINT OF THE DISPATCH: the
%  sans-L-category instrument, the contrast VECTOR, Omega as covariance and as
%  the prefactor of eq:omegaprefac, A4(i)'s per-cell overlap, eq:trunset,
%  eq:omegafloor, and every sans-L >= 3 correction of rounds 2 and 3.  NONE OF
%  THOSE DEPENDED ON THE NORMALISATION.
%
%  WHAT IS STRUCK WITH IT: the five-item cost list; "A4(i) becomes a hypothesis
%  of the operator" and the consequent reversal of M-10; the affine-to-ratio
%  second-order remainder at lem:opnoise; "the binary estimator changes"; and the
%  two-exception sans-L = 2 guarantee.  AT THE RAW CONTRAST THE sans-L = 2
%  REDUCTION HAS **ZERO** EXCEPTIONS AND THE BINARY ESTIMATOR IS UNCHANGED.
%
%  THE ROUND-2 BIJECTION ARGUMENT IS STILL FALSE AND ITS REFUTATION IS STILL
%  WORTH KEEPING, BECAUSE IT IS THE FIRST THING A RE-INVENTOR WILL REACH FOR.
%  It claimed psi could stay at Wtilde while the operator normalised, "because
%  L |-> L Omega is a bijection of def:weight's admissible class".
%     Round 2 checked FOUR BOUNDS -- (L^b_ell) Omega = L^{b Omega}_{ell Omega};
%     ||Omega|| <= 1 so ||ell Omega||_inf <= ||ell||_inf; E||b Omega||^2 <=
%     E||b||^2; and the inverse map with ||Omega^{-1}|| <= sans-L/eps_varrho --
%     and each of the four is CORRECT.  They are jointly INSUFFICIENT, because
%     def:weight DOES NOT DEFINE ITS ADMISSIBLE CLASS BY BOUNDS.  Its clause (a)
%     binds ell := cal-W ell^Pi, with cal-W the Bierens integral operator, so
%     membership requires the generator to lie in R(cal-W).  cal-W's z-sections
%     are Laplace transforms of finite measures against w = exp on a compact
%     index set (B1(i)-(ii), (iv.c)), hence REAL-ANALYTIC in z; A4(ii) makes
%     varrho, and therefore Omega, only Hoelder of order r.  Right-multiplication
%     by a non-analytic Omega(z) takes an element of R(cal-W) OUT of R(cal-W).
%     The counterexample runs already at sans-L = 2 with p = 1 and a varrho that
%     is Hoelder but non-analytic at a point.  SO THE TWO PARAMETERISATIONS DO
%     NOT HAVE THE SAME ADMISSIBLE WEIGHT CLASS.
%     THE FOUR BOUND-CHECKS ARE LEFT ON RECORD RATHER THAN DELETED: they are
%     true, and the lesson is that checking the bounds of a class is not checking
%     membership of it when the class is defined by a RANGE condition.  THE SAME
%     ANALYTIC-VERSUS-HOELDER MECHANISM IS WHAT MAKES THE AUTHOR'S GROUND BITE:
%     it is why a normalised A_cond cannot be assumed to inherit the range and
%     decay statements A2, A2', A1c and H1 are made at.
%
%  WHAT THE REJECTED ROUTE WOULD HAVE BOUGHT, RECORDED SO THAT THE TRADE IS
%  LEGIBLE AND NOT RE-LITIGATED FROM MEMORY.  Three interpretational facts, all
%  of them identities and NONE of them an efficiency or conditioning claim:
%    (a) by eq:omegaprefac the conditional forward residual under Wtilde^Omega
%        would be the cell difference d(u|z) itself, where under Wtilde it is
%        Omega(z) d(u|z), a matrix mixture of the differences;
%    (b) Wtilde^Omega is the vector of INVERSE-PROBABILITY (Horvitz-Thompson)
%        contrasts of each cell against the base;
%    (c) the quadratic form Wtilde' Omega^{-1} Wtilde = (1 - varrho_W)/varrho_W
%        is BASE-CELL-FREE, which the componentwise diag(varrho_l(1-varrho_l))
%        is not.
%  ALL THREE ARE RETAINED, AS INTERPRETATION, AT rem:htread.  They explain what
%  the contrast means.  THEY ARE NOT THE ESTIMATOR'S METRIC.
%
%  TWO CLAIMS MADE FOR THE REJECTED ROUTE WERE UNSUPPORTED AND STAY WITHDRAWN,
%  because they would otherwise be re-offered as reasons to revisit the decision.
%    (i) "EFFICIENCY."  NO EFFICIENCY CALCULATION EXISTS IN THIS CORPUS, and
%        Omega^{-1} is not the efficient conditional weight for this moment.  For
%        the indicator moment at (u, t) that weight is
%        Var(Wtilde 1{Delta U <= u} | Delta Z)^{-1}, and UNDER A1 the variance is
%        Omega(z) F(u|z): Wtilde is conditionally independent of Delta U and
%        1{.}^2 = 1{.}, so E[Wtilde Wtilde' 1{.} | Delta Z] = Omega F while
%        E[Wtilde 1{.} | Delta Z] = 0.  Omega^{-1} therefore agrees with it only
%        UP TO THE SCALAR F(u|z), which depends on both u and z.  A genuine
%        efficiency statement would concern the covariance OPERATOR across (u,t),
%        not a matrix.
%   (ii) "CONDITIONING."  Also unsupported, and it may run BACKWARDS: Omega^{-1}
%        has norm of order eps_varrho^{-1}, so normalising UP-WEIGHTS thin-cell
%        observations -- which is where Section 6's sans-L = 4 design sits, its
%        thinnest cell having n = 66.
%
%  AND ONE GROUND OFFERED IN ROUND 2 WAS A MISQUOTATION.  It cited
%  lf_panel.R:308-313 as saying the alternative "solves an equation whose
%  population zero is not varphi-dagger".  THAT PASSAGE IS ABOUT CENTRING --
%  pooled w - wbar against conditional w - rho -- AND NOT ABOUT NORMALISATION.
%  A global positive scalar changes no root set, no kernel and no range, so it
%  cannot move a population zero.  THE ENGINE'S EXISTING scal IS EXACTLY SUCH A
%  SCALAR (lf_panel.R:385, :401 set wbar <- mean(w), scal <- wbar*(1-wbar)), and
%  :320 records that it "leaves the discrepancy on the scale the stopping
%  constants were calibrated against".  THAT ARRANGEMENT IS UNDISTURBED BY THIS
%  DISPATCH: with the raw contrast adopted, eq:hatT is the pre-dispatch display
%  at sans-L = 2 and the engine needs no change on this account.
% ===========================================================================
%  WHAT THE REVERSAL COSTS: NOTHING IN THIS BLOCK.  The five-item cost list that
%  stood here in round 3 -- A4(i) becoming a hypothesis of the operator and the
%  reversal of M-10; the loss of the free bound |Wtilde| <= 1; lem:bierens (d)
%  acquiring A4(i); T-hat's dependence on varrho-hat turning from affine into a
%  ratio, with its eps^{-2} sensitivity and eps^{-3} curvature; and the change to
%  the binary estimator -- WAS ENTIRELY A CONSEQUENCE OF NORMALISING, AND IS
%  STRUCK WITH IT.  At the raw contrast every one of the five is moot:
%    - ||Wtilde||_inf <= 1 is FREE again, from W finitely supported and
%      varrho_l in [0,1], so A4(i) IS IDLE in prp:wellposed, prp:frechet and
%      prp:compact and M-10's removal STANDS.  Nothing is reversed there.
%    - Wtilde-hat_{i,l} = 1{W_i = l} - varrho-hat_l is AFFINE in varrho-hat, so
%      lem:opnoise's propensity branch is exactly linear and its second-order
%      term is identically zero, as it always was.  No new proof obligation.
%    - lem:bierens needs no overlap in ANY clause: M and T are written at
%      Wtilde, which is bounded and defined everywhere.  The C2 hand-off is
%      therefore unqualified again -- see the corrected note (b) below.
%    - eq:hatT at sans-L = 2 IS the pre-dispatch display, so the binary
%      estimator is unchanged and no engine constant needs re-deriving on
%      account of this dispatch.
%  THE sans-L = 2 REDUCTION THEREFORE HAS **ZERO** EXCEPTIONS.  Round 2 claimed
%  one and round 3 claimed two; both counts were artefacts of the normalisation.
%
%  ONE RELATIONSHIP IS RECORDED HERE, AND IT IS NOT A CONSEQUENCE OF THIS
%  DISPATCH -- IT PRE-DATES IT AND SURVIVES THE REVERSAL.  IT CONCERNS THE
%  DISCREPANCY CONSTANT, AND THE RECONCILIATION RUNS THROUGH c_14, NOT THROUGH
%  tau_DP.  *** TWO EARLIER FRAMINGS OF THIS NOTE WERE WRONG AND ARE STRUCK: that
%  the engine ran tau_DP = 0.7 "below the theory's floor", and that the code and
%  the theory carried "two noise-level conventions" reconciled by a ratio c with
%  tau^{code} c > 2.  Neither is the relationship. ***
%
%  THE THREE OBJECTS, KEPT APART.
%    (1) THE THEORY'S delta-hat_n IS A DETERMINISTIC RATE, defined at A5(i) as
%        max{n^{-1/2} sqrt(ln n), r_n} WITH CONSTANT 1.  It is a rate, not an
%        estimate of anything.
%    (2) THE ENGINE'S res$noise IS A DATA-DEPENDENT PLUG-IN of the L^2(Pi)
%        sampling standard deviation of T-hat.
%    (3) eq:lfsched, in results_s3_estimation.tex, defines the STAGEWISE
%        EFFECTIVE LEVEL
%            delta-hat^eff_{n,k} := c_14 ( delta-hat_n ln n + h_u^{-2} eps-hat^2_{k-1} ),
%        and says in terms that c_14 "must be CHOSEN, because it is not identified
%        by the data".
%  (1) AND (2) ARE DIFFERENT SPECIES -- a rate against a plug-in -- WHICH IS WHY
%  NO CONSTANT EVER RECONCILED THEM.  THE CORRECT STATEMENT IS THAT (2) IS AN
%  ESTIMATOR OF c_14 delta-hat_n ln n, i.e. of the FIRST SUMMAND OF (3), and NOT
%  a substitute for delta-hat_n.
%
%  CONSEQUENCES, AND THE FIRST IS THAT NOTHING DEVIATES FROM A7.
%    - tau_DP = 2.1 > 2.  THE THEORY'S FLOOR IS SATISFIED.  A7 asks tau_DP >= 1
%      (eq:a7amend below) and the Landweber rate asks tau_DP > 2
%      (results_s4_asymptotics.tex:214, proofs.tex:6165 and :6781); 2.1 clears
%      both.  NO SENTENCE MAY CALL THE IMPLEMENTATION'S CONSTANT A DEVIATION FROM
%      A7, AND NONE MAY SPEAK OF A DROPPED ln n.
%    - THE CALIBRATION LIVES IN c_14, at approximately 0.018 on the application's
%      numbers.  That is where a free constant was always meant to sit.
%
%  THE ARITHMETIC, PRINTED BECAUSE IT IS WHAT MAKES THE POINT CONCRETE.  At the
%  application's n = 983, d_Z = 5, r = 2: ln n = 6.891; n^{-1/2} sqrt(ln n) =
%  0.0837; r_n = n^{-r/(2r+d_Z)} = 983^{-2/9} = 0.216; so delta-hat_n = 0.216, the
%  propensity branch binding.  A threshold built directly on it at tau_DP = 2 is
%  2 * 0.216 * 6.891 = 2.98, against a MEASURED initial discrepancy of 0.081 --
%  A FACTOR OF 37 TOO LARGE, so the rule would fire at m = 1, before any Landweber
%  step.  With c_14 ~ 0.018 the effective threshold is
%  2.1 * 0.018 * 0.216 * 6.891 ~ 0.056, which sits BELOW the initial discrepancy
%  and lets the iteration run.
%
%  *** AND THE RESIDUAL GAP IS RECORDED HONESTLY RATHER THAN PRESENTED AS
%  CLOSED. ***  prp:phirate Step 2 needs c_14 >= c_16, and the paper says c_16 is
%  NOT KNOWN TO THE ANALYST.  A SMALL c_14 IS MORE LIKELY TO VIOLATE AN UNKNOWN
%  LOWER BOUND THAN A LARGE ONE, and 0.018 is small.  SO THE CALIBRATION IS
%  ADMISSIBLE BUT UNVERIFIABLE.  THAT IS THE THEORY'S OWN DECLARED GAP -- c_14 is
%  chosen, c_16 is unknown -- AND NOT ONE THE IMPLEMENTATION INTRODUCES.  Phase 6
%  should report the realised c_14 and the realised stopping index; it cannot
%  verify c_14 >= c_16 and must not claim to.
% ===========================================================================

% ---------------------------------------------------------------------------
% PART I.  SAMPLING, MODEL, IDENTIFICATION
% ---------------------------------------------------------------------------

\begin{assumption}[A0: Sampling]\label{ass:sampling}
$\{(Y_{i1},Y_{i2},X_{i1},X_{i2},Z_{i1},Z_{i2},W_i)\}_{i=1}^{n}$ are independent and identically
distributed draws from a distribution $P$ on
$\IR^{2}\times\mathcal X^{2}\times\IR^{2p}\times\{0,1,\dots,\mathsf L-1\}$, where
$\mathcal X\subseteq\IR$ is the carrier of the regressor, $\underline x:=\inf\mathcal X$ and
$\bar x:=\sup\mathcal X$ --- \emph{that much is A0's \textbf{declaration}} --- and $\mathsf L\ge2$ is a
finite integer that does not depend on $n$ and is known to the analyst. \textbf{The condition on
$\mathcal X$} is that it is a \emph{compact interval}, so that $\mathcal X=[\underline x,\bar x]$ with
$\underline x<\bar x$; \emph{that is a clause and is carried by the results that read it}.
\end{assumption}

% A0 is not in the plan's list; it was implicit throughout and is stated here
% because every cross-fitting, empirical-process and rate argument from Step 6
% onwards uses it.
%
% NOTE FOR THE ASSUMPTION-AUDITOR, CORRECTED TWICE.  Through Step 2, A0 was used
% by NOTHING: every result was a statement about the population distribution P
% and held for a single observation.  THAT IS NO LONGER TRUE.  What is consumed,
% and by which clause:
%   lem:levels    COMPACTNESS ONLY, AND THE CONNECTEDNESS HALF OF THIS ENTRY IS
%                 STRUCK 2026-08-26.  It read "COMPACTNESS AND CONNECTEDNESS OF X
%                 -- it integrates from the LEFT ENDPOINT of X, which presupposes
%                 an interval".  THAT OVER-ATTRIBUTES, and the lemma's OWN clause
%                 (ii) is the refutation: (ii) already writes the intersection
%                 form int_{{x >= v} cap X}, which parses on any measurable X.
%                 Read (i) as D_0^{-1}h(x) := int_{[x_lo, x] cap X} h dv -- now
%                 the printed form -- and on a fat Cantor X the lemma is TRUE:
%                 Cauchy-Schwarz is untouched, and injectivity survives because
%                 F(x) = int_{[x_lo,x] cap X} h is absolutely continuous with
%                 F' = h 1_X a.e., so at a.e. point of the level set {F = c} one
%                 has F' = 0, hence h = 0 a.e. on X.  WHAT IS STILL READ: X
%                 BOUNDED, for |X| < oo in the constant and for x_lo := inf X to
%                 be finite -- i.e. A0's COMPACTNESS clause, now on the lemma's
%                 hypothesis line.  What is NOT read is the INTERVAL;
%                 *** COMPLETED 2026-08-27.  THE 2026-08-26 ENTRY WAS SHORT BY
%                 ONE CLAUSE.  It swept (i) and stopped.  Clause (iii)'s proof
%                 ALSO read the interval, at its density step -- "R(D_0^{-1})
%                 contains every centred C^1 function on X" -- and that step is
%                 FALSE on a fat Cantor X (witness zeta(x) = x; see the change
%                 note at the proof).  So the 2026-08-26 narrowing opened the
%                 lemma to an X at which one of its own four proofs did not run.
%                 REPAIRED BY REROUTING, not by restoring the interval:
%                 injectivity of (D_0^{-1})^* on V is clause (ii) at
%                 nu_g = zeta pi, after a continuity extension of the tail
%                 function from X to R.  Neither step reads connectedness, so
%                 THE ENTRY ABOVE STANDS AS AMENDED -- lem:levels reads A0's
%                 COMPACTNESS at (i) AND at (iii), and its INTERVAL nowhere. ***
%   lem:intker    COMPACTNESS AND THE INTERVAL, at clause (iv), twice: for
%                 eq:ivalx to be an element of frak-I, and for the substitution
%                 x = x* + h_x varsigma to carry the X-integral onto a SINGLE
%                 interval.  Added to this list 2026-08-26 together with the
%                 register repair: the interval is a CLAUSE, so the site that
%                 reads it carries it, and lem:intker's hypothesis line now does;
%   lem:balance   the i.i.d. clause, through the sample-splitting convention
%                 stated in its hypotheses (dispatch-3 round 7);
%   A8            "quasi-uniform mesh over X" is undefined on a disconnected X;
%   H1            "the Sobolev scale of order s" likewise;
%   lem:balance Step 4  Schumaker's direct theorem is a statement about splines
%                 on an interval.
%
% ROUND-7 CORRECTION 1: "COMPACT" WAS NOT ENOUGH, AND THE GAP WAS REAL.  A0 said
% only "X compact".  Compactness admits a Cantor set, a convergent sequence with
% its limit, or a countable union of intervals, and C1(ii) plus int f = 1 forces
% only finite positive Lebesgue measure -- WHICH A FAT CANTOR SET HAS.  On such
% an X all four consumers above are undefined or false.  A0 now says COMPACT
% INTERVAL.  The strengthening is free: the application has a scalar continuous
% regressor.
%
% ROUND-7 CORRECTION 2: THE lem:adjfubini APPEAL IS STRUCK, NOT LISTED.  Round 6
% recorded lem:adjfubini as a second A0 consumer, for sigma-finiteness of P.
% THAT IS OVER-ATTRIBUTION: P is a probability measure by def:model, hence
% sigma-finite with no appeal to A0 at all.  The sentence at the foot of that
% proof is corrected rather than the hypothesis list extended -- the cheaper of
% the two repairs, and the correct one.
%
% 2026-08-26 RE-SPELLING OF THE TWO ENDPOINT SYMBOLS, AND IT CHANGES NO LAW.
% A0 read "X = [x_lo, x_hi] subset R a compact INTERVAL, x_lo < x_hi", which
% introduces x_lo and x_hi AS ENDPOINTS OF AN INTERVAL.  That is the same trap
% A4(i)'s existential fell into: negate the restriction and the symbols lose
% their referents, so the block's own (superseded) negation test returned
% "declaration" for a condition that four results read.  A0 now DECLARES
% X subset R with x_lo := inf X and x_hi := sup X -- available at any X -- and
% states the compact interval AS THE CONDITION.  Under the condition the two
% readings agree (inf = min = left endpoint), so no display, constant or proof
% moves; what moves is that the condition is now CITABLE AS A CLAUSE, which is
% what lem:intker's hypothesis line now does and what lem:levels' does not.
%
% DISPATCH-25 AMENDMENT: THE SAMPLE SPACE, AND WHAT IT DOES NOT DO.  The last
% factor read {0,1} through dispatch 24 and now reads {0,1,...,sans-L - 1}.
% Three things about the amendment.
%  (a) AT sans-L = 2 IT IS THE OLD CONDITION.  The support {0,1,...,sans-L - 1}
%      evaluates to {0,1}; the printed glyphs differ, the set does not.
%  (b) "sans-L FINITE, FIXED IN n, AND KNOWN" IS A CONDITION AND IS STATED
%      HERE RATHER THAN ASSUMED.  Finiteness is what makes the contrast vector
%      finite-dimensional and is therefore what every display in the notation
%      block presupposes.  FIXED IN n matters separately: a sans-L growing with
%      n would put a growing number of cells against a fixed overlap floor and
%      A4(i) would become infeasible, since varrho_l >= eps forces
%      eps <= 1/sans-L.  KNOWN matters because the multinomial first stage of
%      A4(iii) fits sans-L - 1 functions and the analyst must know how many.
%      None of the three was needed when the support was {0,1}, which is why
%      none was stated before.
%  (c) IT ADDS NO CONSUMER.  A0's consumers -- lem:levels, lem:intker,
%      lem:balance, A8, H1, and lem:balance Step 4, and NO CLAIM OF
%      EXHAUSTIVENESS IS MADE -- all read the X-side or the i.i.d. clause.
%      (THE WORD "four" STOOD HERE THROUGH 2026-08-25 IN FRONT OF A LIST OF
%      FIVE, and the list is six with lem:intker.  A closure claim short by one,
%      again, and in the same entry that was correcting an over-attribution.)
%      None reads the W factor at all, so the amendment changes no hypothesis
%      list and no proof.  The W factor is read by A4 and by def:hatT, which
%      cite A0 for the fold split and not for the support.

\begin{assumption}[A1: Conditional exogeneity]\label{ass:condind}
$\Delta U \upmodels W \mid \Delta Z$ and $E[\Delta U \mid \Delta Z] = 0$.
\end{assumption}

\noindent\emph{Note on A1.} A1 is a statement about $\sigma$-fields and carries no reference to the
cardinality of $W$'s support. Its representation is \eqref{eq:contrast}, and that representation is
exact: restriction (A) of \cref{def:admissible} written at the $\mathsf L-1$ contrasts is
\emph{equivalent} to $\Delta U\upmodels W\mid\Delta Z$, at every $\Delta Z$, at every $\mathsf L$, and
under \textbf{no} overlap condition --- \emph{A4(i) is not a hypothesis of that equivalence and must not
be listed as one}. The equivalence, the base-cell invariance, and the sense in which extra cells are
\emph{surplus} rather than necessary are in the Discussion of the assumptions, under the heading
\emph{A1 (conditional exogeneity)}.

% WHY THIS NOTE IS PRINTED PROSE AND NOT A \begin{remark}.  Environment numbers
% advance in \input order across the eight compiled files, and this file is
% position 1: a new numbered environment here would shift EVERY subsequent number
% in the article and in the supplement.  The note is therefore written on the
% \noindent\emph{Notes on ...} pattern this file already runs at def:hatT, C3 and
% C5', which prints and numbers nothing.  The fuller treatment is a \paragraph in
% assumptions_supp.tex, which is where this block puts interpretation anyway.
%
% AND THE POINTER IS PROSE, NOT \suppref.  \suppref{Kind}{name} renders
% "Kind S:name" and presumes a NUMBERED target; a \paragraph has no number.
% Writing \suppref at a paragraph would print a reference that resolves to
% nothing and that no build would flag, which is the same species of silent
% failure the header records for \ref at an unnumbered assumption.

\begin{assumption}[A1b: Identification]\label{ass:ident}~
\begin{enumerate}
\item[(i)] \emph{Normalisation.} $\vdagger\in\mathcal V$, that is, $\int_{\mathcal X}\vdagger\,d\pi=0$.
\item[(ii)] \emph{Measure-separability.} Every function of $X_1$ that is almost surely equal to a
function of $X_2$ is almost surely equal to a constant.
\item[(iii)] \emph{Global identification of $\vdagger$.} $\mathcal V_{CI}=\{\vdagger\}$, where
$\mathcal V_{CI}:=\{\varphi\in\mathcal V:\ \Delta Y-\Delta\varphi\upmodels W\mid\Delta Z\}$.
\item[(iv)] \emph{Identification of $\theta_\dagger$.} $E[\Delta Z\,\Delta Z']$ is nonsingular, and
$E\|\Delta Z\|^{2}<\infty$, $E[\Delta U^{2}]<\infty$, $E\|\Delta Z\,\Delta U\|^{2}<\infty$.
\end{enumerate}
\end{assumption}

% (iii) FILLS THE PLACEHOLDER at draft `:111', which reads `edddee'.  The
% intended content is unrecoverable; (iii) is the theorist's proposal, flagged
% for user confirmation in theory_memo.md Sec.7, item N-3.  Rationale: the
% sentence introducing the block promises "restrictions for identification";
% (i) and (ii) are normalisations removing the additive constant and the kernel
% of Delta; what is missing is injectivity, and injectivity of the map
% varphi |-> law of (Delta Y - Delta varphi, W, Delta Z) is exactly (iii).
%
% AN OPEN GAP, RECORDED RATHER THAN CLOSED.  (iii) IS the conclusion
% "varphi_dagger is identified", assumed.  NO PRIMITIVE ROUTE TO IT IS OFFERED
% HERE, and A1c below is not one: A1c is the LINEARISED condition, it neither
% implies nor is implied by (iii), and pointing at it as the "NPIV analogue"
% would gesture at a bridge this block has just said does not exist.  A methods
% referee will ask what observable feature of P delivers (iii); supplying a
% primitive sufficient condition is open work, recorded in theory_memo.md
% Sec.5.  It is NOT to be closed by weakening the statement.
%
% DISPATCH 2: THE FEVE--FLORENS ROUTE WAS TRIED AND IT DOES NOT REACH (iii).
% Feve & Florens (2014, Prop. 2.2) identify the structural function in a
% first-differenced panel with endogenous regressors from two primitives: strong
% identification of (Z_1,Z_2) by the instrument, and measurable separability of
% Z_1 and Z_2.  Transplanted here those become A1d and A1b(ii), and the
% transplant WORKS -- but it lands on A1c (see cor:a1cprim), not on (iii).  The
% obstruction is exact and is worth stating in the paper rather than hiding.
% In their model the restriction is a CONDITIONAL MEAN and the operator is
% LINEAR in varphi, so ker K = {0} IS identification.  Here the restriction is
% CONDITIONAL INDEPENDENCE: membership of V_CI is the single family
%     E[Wtilde varsigma(Delta Y - Delta varphi) | Delta Z] = 0,   varsigma bounded,
% indexed by test functions of the OBSERVED residual at the trial value, whereas
% running the linearised argument globally would need the larger family
%     E[Wtilde varsigma(Delta Y - Delta varphi) varkappa(Delta U) | Delta Z] = 0,
% indexed also by test functions of the TRUE residual -- which is not observed
% at a trial varphi and about which the restriction says nothing.  The two
% families coincide when varphi = varphi_dagger and not otherwise, which is
% precisely the gap between first-order and global identification.
%
% THIS MAKES RECORDING THE GAP DEFENSIBLE RATHER THAN LAZY.  Feve & Florens
% write of their own Proposition 2.2 that "the assumption ... is too strong".
% So the honest primitive is already hard in the LINEAR case, where the operator
% is fixed and the restriction is a conditional mean.  Citing that is the right
% way to present (iii): not "we could not find a primitive" but "the primitive
% is known to be demanding in the linear special case, and the nonlinearity
% introduced by conditional independence separates it from the linearised
% condition that the primitive does deliver".
%
% (iv) IS NEW IN ROUND 2 and is what actually identifies theta_dagger:
%   theta_dagger = (E[Delta Z Delta Z'])^{-1} E[Delta Z (Delta Y - Delta varphi_dagger)].
% Round 1 routed this through A3, which is CONDITIONAL ON A2.  That made
% identification of theta_dagger contingent on a source condition, which is
% false and contradicts the separation recorded in theory_memo.md Sec.3.
% Nothing in this dispatch, and nothing in the identification of theta_dagger,
% depends on A2.
%
% DISPATCH-13 CORRECTION TO THIS COMMENT, AND IT IS A CORRECTION AND NOT A
% RETRACTION.  Through dispatch 12 the sentence above read "...through A3, which
% is nonsingularity of a DIFFERENT matrix (J, involving ell_star) and which is
% itself conditional on A2."  THE PARENTHETICAL IS FALSE.  By lem:momzero(iv),
% under A1 the ell_star summand of J is the ZERO MATRIX, so J = E[Delta Z Delta Z']
% wherever J is defined: A3's matrix and A1b(iv)'s matrix ARE THE SAME MATRIX.
% What made round 2's repair correct is the CONDITIONALITY and nothing else --
% ell_star exists only under A2, so eq:JSig is a well-formed display only under
% A2, and identification must not be routed through a display that may not exist.
% THE REPAIR STANDS; ITS STATED GROUND IS HALVED.  A3 is now retired (ass:rank);
% A1b(iv) is what identifies theta_dagger AND, through lem:momzero(iv), what
% supplies the rank condition the one-step needs.  So (iv) does MORE work than
% round 2 assigned it, not less.

\begin{assumption}[A1c: First-order local identification]\label{ass:completeness}
The linearisation is injective: $\ker\mathcal A_{\mathrm{cond}}=\{0\}$, where
$\mathcal A_{\mathrm{cond}}:\mathcal V\to L^{2}(\Pi)$ is the Fr\'echet derivative of
\cref{prp:frechet} evaluated at $(\vdagger,\theta_\dagger)$.
\end{assumption}

% THE STATEMENT IS THE OPERATOR ONE.  THE PASSAGE TO THE DRAFT'S
% CONDITIONAL-MEAN FORM IS NOW A THEOREM (prp:kerchar), AND THE PARAGRAPH BELOW
% IS THE RECORD OF WHAT THAT THEOREM COST, NOT A LIVE DEBT.
%
% ROUND-7 DESTALING.  Through round 6 this comment opened by asserting IN
% CAPITALS that the two forms "are not equivalent" and that "the passage between
% them is not free", listed five ingredients as unavailable, and only forty lines
% later recorded that prp:kerchar -- in this same block's own results file --
% PROVES the passage.  A reader stopping at the head carried away a debt the file
% closes.  The head is therefore rewritten: what follows is the accounting of the
% passage, and the passage HOLDS under A1, A4(i), B1(i)-(iv), C1(ii) and C2'(i).
% What is NOT free, and remains the reason A1c is stated in operator form, is
% that the passage consumes those five tags -- so the operator form is the one
% every consumer can cite unconditionally.  Write, for a direction h,
%     m_1(u,z) := E[Delta h | Delta U = u, Delta Z = z, W = 1],
%     m(u,z)   := E[Delta h | Delta U = u, Delta Z = z],
% and let f_1(u|z), f(u|z) be the corresponding conditional densities of
% Delta U.  The conditional-mean form says m_1 = m; the operator form says
% (after the binary identity and the Bierens passage) f_1 m_1 = f m weighted
% and integrated.  m_1 = m DOES NOT GIVE f_1 m_1 = f m unless f_1 = f, which
% is A1.  Passing from one to the other consumes, at least: A1; A4(i) so that
% the W = 1 cell is non-null; a lem:bierens Step-5 passage from an
% L^2(Pi)-a.e. statement to a pointwise one; continuity in u; and a DECLARED
% reference measure on (u,z), since "for a.e. (u,z)" is meaningless until one
% is fixed.
% DISPATCH-25 SCOPE NOTE ON THIS PARAGRAPH, AND IT IS A NOTE AND NOT A REWRITE.
% The paragraph is written at sans-L = 2: m_1, f_1 and "the W = 1 cell" are the
% binary instance.  At general sans-L the same accounting runs cell by cell --
% m_l, f_l for l = 0, ..., sans-L - 1, and A4(i) makes EVERY cell non-null, not
% only the l = 1 one -- and the conclusion is unchanged, because the obstruction
% it records (m_l = m does not give f_l m_l = f m without A1) is per cell and
% does not interact across cells.  It is left in the binary spelling because
% prp:kerchar, the theorem it accounts for, lives in a Phase-2 file and the two
% must move together or not at all.  NONE OF THOSE IS AVAILABLE AT THE POINT
% WHERE A1c IS STATED -- which
% is why A1c is stated in operator form.  All of them ARE available by the time
% prp:kerchar is reached, and the corrections to this list are recorded below.
%
% The operator form is also the form every consumer reads: def:spaces, A2, H1
% and lem:adjoint all name ker A_cond or the singular system of A_cond, none
% of them the conditional-mean display.  So the operator form is primary and
% the draft's display is recorded below as an interpretation, not an
% equivalence.
%
% RENAMED, because "completeness" promises more than the condition delivers.
% A1c IS FIRST-ORDER ONLY, AND NO STRENGTHENING OF IT WILL BUY LOCAL
% UNIQUENESS.  By prp:compact, A_cond is COMPACT; an injective compact
% operator on an infinite-dimensional space is not bounded below, so
% inf over ||h|| = 1 of ||A_cond h|| = 0.  Take h = e_j, the j-th right
% singular function.  By prp:frechet(i), a competing solution at
% varphi_dagger + eps e_j would require
%     eps sigma_j = ||A_cond (eps e_j)|| <= C_A eps^2, i.e. eps >= sigma_j / C_A,
% and sigma_j decreases to 0.  The first-order argument therefore excludes
% competitors only at distances below sigma_j / C_A, a threshold that shrinks
% to zero along the singular directions: A1c does not deliver a neighbourhood
% of uniqueness at all.  GLOBAL IDENTIFICATION IS A1b(iii) AND MUST BE ASSUMED
% SEPARATELY; A1c neither implies nor is implied by it.
%
% The draft's `eq-condcomppanel' (`:127--130') is the conditional-mean display
% above, and it needs correcting in its own right: it conditions on (U,W) only,
% which is the unconditional reading abandoned in plan Sec.17, and it should
% condition on Delta Z as well.
% ROUND 2: retagged from "C6" to "A1c"; CLEANUP PASS: restated and renamed.
%
% DISPATCH 2 ADDS THE MISSING BRIDGE, AND IT IS NOW A THEOREM.  The paragraph
% above says that the draft's conditional-mean display and the operator form
% ker A_cond = {0} "are not equivalent, and the passage between them is not
% free", then lists what the passage would consume: A1; A4(i); a lem:bierens
% Step-5 passage from an L^2(Pi)-a.e. statement to a pointwise one; continuity
% in u; and a declared reference measure on (u,z).  prp:kerchar PROVES the
% equivalence, and the accounting is BETTER than that list, not merely equal to
% it.  Two corrections to the list itself:
%
%   (1) THE REFERENCE MEASURE WAS NEVER MISSING.  A first draft of this note
%       claimed the fifth ingredient was supplied by the strengthened B1(iv.a).
%       THAT IS FALSE.  prp:kerchar never uses varpi_u: clauses (b) and (c) are
%       stated P-almost surely and the reference measure on (u,z) is
%       P_{Delta U, Delta Z}, which is available unconditionally.  B1(iv.a)
%       enters only through lem:bierens Step 5, and only its ROUND-2
%       full-support form.  CONSEQUENCE, and it is the stronger statement:
%       prp:kerchar AND cor:a1cprim SURVIVE A COMPLETE ROLLBACK OF THE B1(iv.a)
%       STRENGTHENING.  The A1c bridge and the N-10 closure are independent
%       results that happen to have been proved in the same dispatch.
%
%   (2) A1 AND A4(i) ARE CONSUMED BY ONLY ONE HALF.  prp:kerchar is split:
%       (a) <=> (b) needs neither, and (b) <=> (c) needs both.  cor:a1cprim
%       reads clause (b), so IT NEEDS NEITHER A1 NOR A4(i) -- only B1(i)-(iv),
%       C1, C2'(i), A1b(i)-(ii) and A1d.  Listing A1 and A4(i) on the corollary,
%       as a first draft did, would have charged it for hypotheses it does not
%       use.
%
% What remains open is A1b(iii), which is a different statement.  See
% prp:kerchar, cor:a1cprim and rem:ffcompare in results.tex.

% ============================================================================
%  A1d (ass:instrcomp) MOVED ENTIRE TO assumptions_supp.tex, 2026-08-28, ON THE
%  AUTHOR'S SHORTENING INSTRUCTION AND ON THE A3 / A6 / C4' PRECEDENT.  It printed
%  357 words in the ARTICLE: a struck-through retirement carrying a full proof of
%  its own unsatisfiability.  That is a working-corpus record, not article text.
%  The tag, the \label, the struck banner and the whole ground survive in the
%  supplement under "A1d (RETIRED: instrument completeness given the residual),
%  moved entire", so a stale citation still lands on the reason.
%  NO COUNTER MOVES: \newtheorem*{assumption} is unnumbered.
%  NOTHING IN THE ARTICLE CITED ass:instrcomp -- the only article-side reference
%  was inside the block itself, and it went with it.
% ============================================================================
% THE SUBJECT WAS WRONG AND THE CENSUS WAS SHORT BY ONE, BOTH REPAIRED 2026-08-26.
%  (1) SUBJECT.  This paragraph read: "lem:nusigma(v) is the unconditional shadow
%      of the display above and was in this corpus throughout.  IT is not
%      maintained, is not available for use, and is a hypothesis of no result in
%      this programme; ... ITS only consumer was withdrawn with it."  The nearest
%      antecedent of "It" is lem:nusigma(v) -- a PROVED LEMMA CLAUSE, which is
%      maintained, is available, and is used in the A1d refutation printed four
%      lines above.  Read literally the sentence retired a lemma.  A1d is now the
%      grammatical subject and the sentence is a paragraph of its own.
%  (2) CENSUS.  A1d has carried a printed "RETIRED, 2026-08-05" banner since that
%      date, and the Discussion's retirement census in assumptions_supp.tex read
%      "four retirements resting on three grounds" and named A3, A2'(ii), A6 and
%      B1(vi).  FIVE, NOT FOUR.  The omission is this project's standing defect
%      species -- a closure claim short by one -- and it is what left A1d with no
%      standing in any of the block's own status registers even though its own
%      environment says RETIRED.  Both halves of the count now agree.

% ===========================================================================
%  A1d IS RETIRED, 2026-08-05, BECAUSE IT IS UNSATISFIABLE IN THIS MODEL.
%
%  It read: for every Lambda in L^2(P_{X_1,X_2}),
%      E[Wtilde Lambda(X_1,X_2) | Delta U, Delta Z] = 0 a.s.  =>  Lambda = 0 a.s.
%
%  Take Lambda == 1, which is in L^2 because P is a probability measure.  Wtilde
%  = W - rho(Delta Z) with rho(Delta Z) = P(W = 1 | Delta Z), and A1's first
%  clause is Delta U independent of W given Delta Z, so
%      E[Wtilde | Delta U, Delta Z] = E[W | Delta Z] - rho(Delta Z) = 0 a.s.
%  The antecedent holds and the consequent fails.  Hence A1 => not-A1d, and
%  since A1b(iii) implies A1's conditional independence at varphi-dagger, also
%  A1b(iii) => not-A1d.  The condition is empty on the model class every other
%  result in this block is stated on.
%
%  THE GENERAL FACT, which is the part worth keeping: Wtilde is CONDITIONALLY
%  CENTRED given Delta Z, so the operator Lambda |-> E[Wtilde Lambda | Delta U,
%  Delta Z] annihilates every Lambda that agrees a.s. with a sigma(Delta U,
%  Delta Z)-measurable function -- the constants being the trivial case.  ANY
%  completeness condition posed on that operator over a class containing the
%  constants is therefore unsatisfiable here.  This is what separates the
%  conditional-independence design from Feve & Florens's conditional-mean one,
%  whose operator g |-> E[g | W] does not annihilate constants: their hypothesis
%  transplants as a CONTRADICTION, not merely as an excess.
%
%  The corpus already contained the refutation and never connected it:
%  lem:nusigma(v) proves E[Wtilde Xi(Delta U, Delta Z)] = 0 under A1 for every
%  integrable Xi, which is the unconditional shadow of the display above.  Same
%  species as the retired A3, whose refutation was also already in the file.
%
%  CONSEQUENCE.  cor:a1cprim, A1d's only consumer, is withdrawn: see the marker
%  at its former site in results_s2_identification.tex.  A1c is a hypothesis of
%  no result in the corpus and now has no primitive route offered.  A1b(iii)
%  stands as an openly assumed global identification condition, which is what
%  the block's own discussion at :182-189 already anticipated a referee would
%  press on.  Nothing downstream changes: no theorem cited A1d or A1c.
% ===========================================================================

% A1d IS NOT MAINTAINED.  It is offered as a PRIMITIVE SUFFICIENT CONDITION for
% A1c and is cited by exactly one result, cor:a1cprim.  No theorem in the
% programme assumes it; A1c remains the maintained condition.  It is stated in
% the block rather than inline so that cor:a1cprim can cite it by tag, per house
% style, and so that the assumption-auditor sees a declared single consumer.
%
% WHAT IT IS.  A1d is the exact analogue, for the conditional-independence
% design, of the "strong identification of (Z_1,Z_2) by W" hypothesis of
% Feve & Florens (2014, Prop. 2.2): injectivity of the map
% Lambda |-> E[Wtilde Lambda | Delta U, Delta Z] on functions of the two periods'
% endogenous regressors.  In their linear conditional-mean panel model that
% hypothesis, together with measurable separability of Z_1 and Z_2 -- our
% A1b(ii) -- DELIVERS IDENTIFICATION OUTRIGHT, because their operator is linear
% in varphi.  Here it delivers A1c and no more: our operator equation is
% NONLINEAR in varphi, so injectivity of the linearisation is a first-order
% statement.  The separation is not a defect of the adaptation; it is what
% distinguishes a conditional-independence restriction from a conditional-mean
% one, and it is discussed at cor:a1cprim and rem:ffcompare.
%
% AND IT IS STRONG.  Feve & Florens themselves write that the assumption of
% their Proposition 2.2 "is too strong".  That judgement transfers: A1d is a
% completeness condition on a two-dimensional argument (X_1,X_2) with a
% (sans-L - 1)-dimensional contrast Wtilde doing the separating, so it will be
% harder here than there, not easier.  It is recorded as a route, priced
% honestly, and NOT adopted.
% DISPATCH 25 DOES NOT REVIVE A1d, AND THE REFUTATION IS CARDINALITY-FREE.  The
% clause was retired on 2026-08-05 because Wtilde is CONDITIONALLY CENTRED given
% Delta Z, so the operator Lambda |-> E[Wtilde Lambda | Delta U, Delta Z]
% annihilates the constants.  That argument reads E[Wtilde_l | Delta Z] = 0 for
% every l, which holds at every sans-L by construction of eq:contrast, so more
% cells do not repair it: at sans-L = 4 the operator annihilates the constants in
% all three coordinates at once.  The dimension of the contrast was NEVER what
% made the condition fail.

% ---------------------------------------------------------------------------
% PART II.  THE INDEX SET AND THE WEIGHT FAMILY
% ---------------------------------------------------------------------------

\begin{assumption}[B1: Bierens family and range space]\label{ass:bierens}~
\begin{enumerate}
\item[(i)] \emph{Transform.} $\Phi:\IR^{p}\to\IR^{p}$ is Borel, one-to-one, and bounded:
$c_\Phi:=\sup_{z}\|\Phi(z)\|_\infty<\infty$. The maintained choice is $\Phi$ componentwise $\arctan$,
for which $c_\Phi=\pi/2$. \emph{The value of $c_\Phi$ is not pinned by this block and no result reads
it}: only $c_\Phi<\infty$ is consumed, at \suppref{Lemma}{bierenscore} \textup{(compactness of the
image cube)} and through $\bar w$ at \textup{(ii)}. \emph{$c_\Phi$ is not separately identified from
the scale of $\mathcal T$} --- see \suppref{Lemma}{scaleinv}\textup{(vii)} --- \emph{and enters the
programme only through $c_\Pi$ of \textup{(iv.c)}.}
\item[(ii)] \emph{Weight.} $w=\exp$. Write $\bar w(t):=e^{c_\Phi\|t\|_1}$, so that
$|w(t'\Phi(z))|\le\bar w(t)$ for every $z$.
\item[(iii)] \emph{Index set.} $\mathcal T\subset\IR^{p}$ has strictly positive Lebesgue measure.
\item[(iv)] \emph{Range reference measure.} $\Pi=\Pi_u\otimes\Pi_t\otimes\Pi_{\mathsf w}$, where
\begin{enumerate}
\item[(iv.a)] $\Pi_u$ is a finite measure on $\IR$ that is absolutely continuous with respect to
Lebesgue measure, with a density $\varpi_u$ that is bounded, strictly positive at every point of $\IR$,
and continuously differentiable with $\|\varpi_u'\|_\infty<\infty$; in particular $\Pi_u$ has full
topological support, $\Pi_u(O)>0$ for every nonempty open $O\subseteq\IR$;
\item[(iv.b)] $\Pi_t$ is a finite measure on $\mathcal T$ for which there exists a Borel set
$\mathcal T_0\subseteq\mathcal T$ such that $\mathrm{Leb}|_{\mathcal T_0}$ and $\Pi_t|_{\mathcal T_0}$
are \emph{not mutually singular};
\item[(iv.c)] $c_\Pi^{2}:=\Pi_u(\IR)\int_{\mathcal T}\bar w(t)^{2}\,\Pi_t(dt)<\infty$;
\item[(iv.d)] $\Pi_{\mathsf w}$ is a \emph{probability} measure on $\{1,\dots,\mathsf L-1\}$ whose masses
$\varpi_{\mathsf w}(l):=\Pi_{\mathsf w}(\{l\})$ are strictly positive at every
$l\in\{1,\dots,\mathsf L-1\}$.
\end{enumerate}
\item[(v)] \emph{Correction measure.} $\gamma$ is a finite signed measure on $\mathcal T$ with
$\int_{\mathcal T}\bar w(t)\,|\gamma|(dt)<\infty$; and, where $\partial_u L$ is required,
$|\partial_u\lambda(u,t)|\le\Lambda_1(t)$ for all $u$, with
$\int_{\mathcal T}\Lambda_1(t)\,\bar w(t)\,|\gamma|(dt)<\infty$.
\end{enumerate}
\end{assumption}

% ---------------------------------------------------------------------------
% B1(ii) IS RESTRICTED TO w = exp, AND THAT IS A CORRECTION, NOT A CONVENIENCE.
%
% Round 1 stated B1(ii) as "real-analytic and non-polynomial", which is the
% textbook Bierens phrasing.  UNDER THAT CONDITION lem:bierenscore IS FALSE.
% Counterexample, inside the paper's own maintained transform: let V = +-1 with
% probability 1/2 each, let Phi be the rescaled arctan (an ODD function), so
% Phi(-1) = -Phi(1), and set epsilon := V.  Then E|epsilon| < infinity and
% E[epsilon | V] = V != 0, yet for w = cos and EVERY t,
%      E[epsilon cos(t Phi(V))] = (1/2)cos(t Phi(1)) - (1/2)cos(-t Phi(1)) = 0.
% (ROUND-8: the value Phi(1) was written "a" through round 7, which collided with
% the DEGREE OF ILL-POSEDNESS a of plan Sec.2.6 -- and glossary Sec.1.6 declares
% the sans-serif constant slot expressly to keep generic constants away from that
% very letter.  The value is written out; no symbol is introduced.  This also
% removes the juxtaposition cos(ta) of the RESERVED Bierens index t with a.)
% B1(i)-(iii) all hold and the converse fails.  Structurally: Phi ODD against
% w EVEN means the family sees only the EVEN part of the image measure, and
% Step 4 of the proof of lem:bierenscore recovers only the even moments.  The
% failure is not abstract -- it breaks lem:bierens itself: with Delta Z
% two-valued at +-c and w = cos, condition (c) holds whenever the UNCONDITIONAL
% restriction does, while the Delta-Z-conditional restriction can fail through
% opposite-signed violations that cancel.
%
% w = exp is the maintained choice in plan Sec.2.2, it is what Bierens (1990)
% Thm 1 actually characterises, and it is what any implementation uses.  The
% general-w statement is recorded as rem:swintercept in results.tex, with the
% intercept-indexed Stinchcombe--White family that repairs the failure; it is
% not claimed here because it would change the range space and is consumed
% nowhere in the programme.
%
% CONSEQUENCE FOR PLAN Sec.1: the notation table listed "cos s / sin s" as
% admissible Bierens weights.  Both are inadmissible against an odd Phi.
% ---------------------------------------------------------------------------
%
% Reading of the rest of B1.
%  (i) BOUNDEDNESS does THREE jobs.  It makes w = exp integrable against
%      epsilon with NO moment condition whatsoever on Delta Z, since
%      |w(t'Phi(z))| <= bar-w(t) uniformly in z.  It makes the image measure
%      B |-> E[eps 1{Phi(V) in B}] COMPACTLY SUPPORTED, which is what makes its
%      moment problem determinate in the proof of lem:bierenscore.  And it is
%      why the arctan device may not be skipped.
%      ONE-TO-ONE is what makes sigma(Phi(Delta Z)) = sigma(Delta Z), so that a
%      statement about Phi(Delta Z) is a statement about Delta Z.  Borel image
%      is NOT assumed: an injective Borel map between Polish spaces is
%      automatically bi-measurable (Lusin--Souslin), and for the maintained
%      choice Phi is a homeomorphism onto an open cube, so this is elementary.
%
%  (iv) COMPACTNESS OF T IS SUFFICIENT BUT NOT NECESSARY, and (iv) is now
%      stated at the strength the proofs consume and no more.
%      (iv.a) Step 5 of lem:bierens needs the zero set of M(.,t) to be DENSE
%          in R, which follows from full topological support alone; equivalence
%          to Lebesgue measure, assumed in round 1, is far stronger and idle.
%
%          DISPATCH-2 STRENGTHENING, AND IT IS WHAT CLOSES N-10.  (iv.a) now
%          also requires Pi_u to be ABSOLUTELY CONTINUOUS with a bounded,
%          everywhere strictly positive density varpi_u.  This IMPLIES the old
%          full-support clause (a positive density gives every nonempty open set
%          positive measure), so every proof written against the round-2 form of
%          (iv.a) -- lem:bierens Step 5 is the only one -- goes through verbatim.
%
%          WHY IT IS NEEDED.  The adjoint A*_cond produces, in raw form, an
%          object measurable with respect to F_XZW = sigma(X_1,X_2,Delta Z,W):
%          the u-integration against Pi_u sits against the conditional density
%          f_{Delta U | F_XZW}(u), and until Pi_u(du) can be written varpi_u(u)du
%          that pairing is NOT a conditional expectation and has no
%          (Delta U, Delta Z)-representative.  With the density in hand,
%              int f_{Delta U | F}(u) ell(u,t) Pi_u(du)
%                    = E[ varpi_u(Delta U) ell(Delta U, t) | F_XZW ],
%          and the Bierens induction map of lem:adjoint(ii) is a CONSTRUCTION
%          rather than the assumption A2'(0) used to be.  See lem:adjoint and
%          theory_memo.md Sec.7 item N-10.
%
%          WHY THE DERIVATIVE CLAUSE IS THERE, AND IT IS THE N-10 CLOSURE'S OWN
%          DEBT.  Once lem:adjoint(ii) makes ell_star = W ell^Pi_star, the
%          factor varpi_u(u) sits inside ell_star, so
%              partial_u ell_star(u,z)
%                 = varpi_u'(u) int ell^Pi(u,t) w(t'Phi(z)) Pi_t(dt)
%                 + varpi_u(u) int partial_u ell^Pi(u,t) w(t'Phi(z)) Pi_t(dt).
%          A2'(i) bounds ||partial_u ell_star||_inf.  WITHOUT varpi_u in C^1 THAT
%          DERIVATIVE NEED NOT EXIST, and A2'(i) would be a bound on an object
%          the block does not deliver.  Before N-10 was closed the defect was
%          invisible: A2'(0) made ell_star primitive, so its derivative was
%          whatever the assumption said it was.  Constructing ell_star exposed
%          the debt.  The clause also turns the Pi_u SELECTION LEVER from an
%          aspiration into a theorem: with varpi_u in C^1 and ||varpi_u'||_inf
%          finite, the display above HOLDS whenever partial_u ell^Pi exists, so
%          the design's influence on q and on sigma_q -- hence on
%          prp:sharpquad's sharp constant -- is computable rather than merely
%          conjectured.  BINDING ON STEP 4.
%
%          WHAT IT COSTS: NOTHING ABOUT P.  Pi_u is a reference measure on the
%          RANGE side, chosen by the analyst, exactly like Phi and w.  Any
%          logistic, Gaussian or Cauchy Pi_u satisfies all four clauses --
%          bounded, positive, C^1, bounded derivative -- and every
%          implementation already uses such a Pi_u.  It is a restriction on the
%          normalisation, not on the data-generating process, and it is the only
%          place in the block where that distinction does real work.
%
%          WHAT DOES *NOT* DEPEND ON ANY OF THIS.  prp:kerchar and cor:a1cprim
%          use NO clause of (iv.a) beyond the round-2 full-support one, and they
%          use that only through lem:bierens Step 5.  Their statements are P-a.s.
%          and their reference measure on (u,z) is P_{Delta U, Delta Z}, which is
%          available unconditionally.  THEY SURVIVE A ROLLBACK OF THIS ENTIRE
%          STRENGTHENING.  An earlier note claimed the opposite; see the
%          correction under A1c.
%      (iv.b) Step 5 needs: no Pi_t-null set swallows T_0.  That is EXACTLY
%          non-mutual-singularity of Leb|_{T_0} and Pi_t|_{T_0}.  (If N is
%          Pi_t-null and Leb(T_0 \ N) = 0, then Leb|_{T_0} sits on N and
%          Pi_t|_{T_0} on its complement, i.e. they are mutually singular;
%          conversely non-singularity forbids such an N.)  Equivalently, by the
%          Lebesgue decomposition, (iv.b) SAYS THAT Pi_t HAS A NONZERO
%          ABSOLUTELY CONTINUOUS PART.  Non-mutual-singularity already forces
%          Leb(T_0) > 0 and Pi_t(T_0) > 0, since the zero measure is mutually
%          singular with everything.  Round 1 assumed EQUIVALENCE on T_0, which
%          is two-sided; only one direction was ever invoked, and even that is
%          stronger than needed.
%
%          WHAT (iv.b) EXCLUDES -- corrected in the cleanup pass.  It excludes
%          EVERY LEBESGUE-SINGULAR Pi_t, not merely the atomic ones: a Cantor
%          measure is non-atomic, has uncountable support, and is excluded.
%          Round 2's gloss ("it excludes a purely atomic Pi_t, and that
%          exclusion is the point") named the smallest of the excluded classes.
%
%          AND (iv.b) IS SUFFICIENT, NOT NECESSARY.  What the proof really
%          needs is an ANALYTIC-UNIQUENESS condition on supp Pi_t: by
%          lem:bierenscore Step 2 the map t |-> M(u,t) is real-analytic on all
%          of R^p, so it is enough that its zero set cannot contain supp Pi_t
%          up to a Pi_t-null set unless it vanishes identically.  For p = 1 any
%          Pi_t whose support has an accumulation point suffices -- a Cantor
%          measure included, since a Cantor set is perfect -- by the
%          one-variable identity theorem, and lem:bierens survives such a Pi_t.
%          For p >= 2 the condition is that supp Pi_t not lie inside a proper
%          real-analytic variety; positive Lebesgue measure guarantees that,
%          and so do many Lebesgue-null sets, but perfectness alone does not
%          (a hyperplane is perfect and is the zero set of a coordinate).
%          (iv.b) is retained as the clean sufficient version.
%
%          What (iv.b) does exclude in every dimension is a FINITE grid.  THAT
%          WAS THE GROUND on which (vi) was classed an identification assumption
%          and not an approximation device.  (vi) IS RETIRED (dispatch 14); the
%          exclusion (iv.b) performs is unaffected, and the sentence is kept in
%          the corrected register so that a later dispatch reinstating the
%          growing-grid design finds it.
%      (iv.c) is the integrability that makes L^2(Pi) the right range space;
%          for w = exp it reads Pi_u(R) int e^{2 c_Phi ||t||_1} Pi_t(dt)
%          < infinity, and any compact T with a finite Pi_t satisfies it.
%          IT IS UNCHANGED BY DISPATCH 25's THIRD FACTOR.  The full-measure
%          constant is Pi_u(R) int bar-w^2 dPi_t times Pi_w({1,...,sans-L - 1}),
%          and B1(iv.d) fixes that last factor at 1, so c_Pi is the same number
%          at every sans-L and every constant carrying it -- prp:frechet's
%          operator-norm bound, prp:compact's Hilbert-Schmidt bound, C_A -- is
%          unmoved.  THAT IS WHY (iv.d) SAYS "PROBABILITY" AND NOT "FINITE".
%      (iv.d) IS NEW IN DISPATCH 25.  APPENDED, SO NO CLAUSE NUMBER MOVES; this
%          is the A2'(iii) and A5(vi) precedent, both of which appended a clause
%          under a live tag.  WHAT IT COSTS: NOTHING ABOUT P, exactly as (iv.a)
%          costs nothing about P.  Pi_w is a reference measure on the range side
%          and the analyst chooses it; the uniform choice
%          varpi_w(l) = 1/(sans-L - 1) satisfies it.
%          WHY STRICT POSITIVITY IS A CONDITION AND NOT A CONVENTION.  The
%          L^2(Pi) norm of T(varphi,theta) is sum_l varpi_w(l) times the
%          L^2(Pi_u (x) Pi_t) norm of its l-th contrast component.  If some
%          varpi_w(l) were zero that contrast would be INVISIBLE to the norm,
%          ||T|| = 0 would not imply the l-th restriction, and lem:bierens would
%          lose the direction it needs -- the same failure mode (iv.b) excludes
%          on the t-side, one factor over.  A dispatch that lets a mass go to
%          zero has silently coarsened the instrument.
%          EFFECT ON HYPOTHESIS LISTS.  Any list block-citing B1(iv) carries
%          (iv.d) as well, on the register rule stated in the header.  [SCOPE
%          NARROWED 2026-08-24, SECOND PASS: six of the lines that block-cited
%          B1(iv) now cite BY CLAUSE and so no longer reach this sentence --
%          prp:wellposed, prp:frechet, prp:compact, lem:opnoise, lem:unifop,
%          lem:aerr.  They carry (iv.d) EXPLICITLY, which is stronger than
%          carrying it by block, and they no longer carry (iv.b) at all.  The
%          sentence is still live at the ten block-form lines listed in the
%          header, lem:adjfubini among them.]  This is
%          harmless and is recorded rather than left to be discovered: at
%          sans-L = 2 the clause is FORCED (delta_1 is the only probability
%          measure on a one-point set) and so is satisfied by every earlier
%          round's Pi; at sans-L >= 3 it is a condition on a chosen object.
%          NO EXISTING LIST IS STRENGTHENED IN ANY WAY THAT BEARS ON P.
%          AND THE CONVERSE HALF IS THE AUTHOR'S 2026-08-24 RULING, STATED IN
%          THE HEADER: a DEFINITION's block citation -- def:spaces's bare B1(iv)
%          -- is a LOCATING POINTER and imposes nothing, so a result whose proof
%          reads (iv.d) must carry it on its OWN line and cannot borrow it from
%          def:spaces.  Eight did not; all eight are repaired that day and the
%          list of them is in the header.
%          [SCOPE, 2026-08-26: the pointer half holds HERE, at def:spaces's bare
%          B1(iv), which locates.  IT IS NOT A UNIVERSAL ABOUT DEFINITIONS, and
%          the printed rule no longer states it as one: a CONSTITUTIVE citation
%          -- A2b, which IS "admits a version satisfying A2'(i.a)-(i.b)", and
%          def:hatT, which builds eq:trunset out of A4(i)'s constant -- carries
%          what it cites, because deleting the citation leaves no object.]
%
%  (v)  By lem:density the estimator SOLVES FOR L rather than integrating a
%      weighting measure, so the first half of (v) is nearly moot on the
%      theta-side: no t-integral survives there.  It binds only if one insists
%      on the integral representation of L.  The SECOND half is the explicit
%      dominating function for differentiating L under the t-integral in u --
%      round 1 asserted that the exponential-moment condition "licenses that
%      directly" without naming either the integral or the dominating function.
%
%  (vi) IS RETIRED (dispatch 14) AND IS NOT AVAILABLE AS A HYPOTHESIS.  The
%      reading below is the reading it HAD while it was maintained, kept because
%      the clause statement is kept: it was the Dominguez--Lobato requirement,
%      an IDENTIFICATION assumption and not a regularity condition, since with a
%      FIXED finite grid the operator equation need not identify
%      varphi_dagger.  Precedent: Poirier (2017); lead case Donald, Imbens &
%      Newey (2009).  WHAT REPLACED IT: nothing -- identification of
%      varphi_dagger in this block is A1b(iii), assumed outright, and the
%      counterexample-hunter's finite-grid witness is recorded at the struck
%      clause so that the statement and the witness are found together.

% ---------------------------------------------------------------------------
% PART III.  SOURCE CONDITIONS AND THE SMOOTHNESS LINK
% ---------------------------------------------------------------------------

\begin{assumption}[A2: Source condition for $\theta$]\label{ass:sourcecond}
Each coordinate of $\nu/\pi$ lies in $\mathcal R(\mathcal A^\ast_{\mathrm{cond}})$, the range of the
adjoint of $\mathcal A_{\mathrm{cond}}:\mathcal V\to L^{2}(\Pi)$, where $\nu:=\nu_{\Delta Z}$ and
$\nu_g(x):=E[g\mid X_2=x]f_{X_2}(x)-E[g\mid X_1=x]f_{X_1}(x)$ as in \eqref{eq:nudef}. Equivalently
(Picard):
$\nu/\pi\perp\ker\mathcal A_{\mathrm{cond}}$ in $L^{2}(\pi)$ and
$\sum_{j\ge1}\sigma_j(\mathcal A_{\mathrm{cond}})^{-2}\,\langle\nu/\pi,\,e_j\rangle_{L^{2}(\pi)}^{2}<\infty$.
\end{assumption}

% ---------------------------------------------------------------------------
% THE OBJECT IS nu / pi, NOT nu.  ROUND-2 CORRECTION, AND IT PROPAGATES.
%
% prp:compact computes the kernel of A_cond as k(u,t;x) = nu_{g_{u,t}}(x)/pi(x).
% The adjoint therefore maps into functions of the form nu_{tilde g}/pi, so
% R(A*_cond) consists of nu-objects DIVIDED BY pi.  A membership test settles it
% independently of any range computation.  By prp:compact Step 2 every kernel
% slice lies in V, so R(A*_cond) is contained in V BY CONSTRUCTION OF THE
% ADJOINT, and V is exactly {zeta : int zeta pi dx = 0}.  Now:
%     zeta = nu     requires int nu pi dx = E[Delta Z (pi(X_2) - pi(X_1))] = 0,
%                   which is FALSE for a generic reference density;
%     zeta = nu/pi  requires int nu dx = E[Delta Z (1 - 1)] = 0, TRUE always.
% So nu/pi is in the right space and nu is not.  Equivalently: nu_g/pi lies in V
% automatically, for every g, since int (nu_g/pi) d pi = int nu_g dx = 0.
%
% NOTE the reason.  This does NOT come from R(A*) subset (ker A)^perp together
% with Delta 1 = 0: once def:spaces fixes the domain at V, the constants are not
% in the domain and that orthogonality argument is unavailable.  It comes from
% the construction of the adjoint, which is simpler and is what prp:compact
% Step 2 actually proves.
%
% This is a bookkeeping consequence of taking the adjoint in L^2(pi), which is
% forced by plan Sec.19.1's dual-pairing principle: the rate for h is in
% L^2(pi), so the residual and the range condition belong there too.  The
% memo's Sec.0.1 identity "A*_cond ell = nu_{Wtilde ell}" is an L^2(dx)
% statement; under the L^2(pi) convention it carries a pi^{-1}.  ONE OF THE TWO
% HAS TO GIVE, AND THE L^2(pi) CONVENTION WINS.  lem:adjoint (Step 3) inherits
% this and must state A*_cond ell = nu_{Wtilde ell}/pi.  Recorded because tags
% are declared stable and Step 3 must not re-open it.
%
% DISPATCH 2: DISCHARGED, AND WITH ONE FURTHER CORRECTION.  lem:adjoint(iii)
% states A*_cond ell = nu_{Wtilde (W ell)} / pi, with the induction map W of
% lem:adjoint(ii) between ell's range-space form and its observable-space form.
% The pi^{-1} is as recorded above; what the round-2 note could not yet say is
% that the "Wtilde ell" inside nu_g is NOT the solution of the adjoint equation
% but its image under W.  Writing nu_{Wtilde ell} without the W is a type error
% of the same kind the round-2 note corrected for pi, one level in.
%
% A2 ALSO ACQUIRES A PRIMITIVE READING.  By lem:adjoint(iii), A2 holds if and
% only if nu_{Delta Z + Wtilde (W ell)} == 0 for some ell in L^2(Pi) -- i.e. if
% and only if the ORTHOGONALITY EQUATION the estimator targets is solvable
% WITHIN THE REACHABLE CLASS R(W).  Since R(W) is dense but not known to be all
% of L^2(P_{Delta U, Delta Z}), A2 IMPLIES solvability of that equation and the
% converse is open.  A2 may therefore be strictly stronger than what the sieve
% estimator needs.  RECORDED, NOT CLAIMED (theory_memo.md finding F-20);
% claiming it would require characterising R(W).
% ---------------------------------------------------------------------------

\begin{assumption}[A2$'$: Source condition with usable representer]\label{ass:sourcecondp}
A2 holds, and the selected solution
$\ell^\Pi_\star\in\ell^\Pi_\dagger+\ker\mathcal A^\ast_{\mathrm{cond}}$ of
$\mathcal A^\ast_{\mathrm{cond}}\ell^\Pi=-\nu/\pi$, where $\ell^\Pi_\dagger$ denotes any one
particular solution, so that the displayed set is the whole solution coset and $\ell^\Pi_\star$ is a
selection from it, \emph{admits a version} satisfying
\begin{enumerate}
\item[(i)] \emph{Range-space regularity.}
\begin{enumerate}
\item[(i.a)] $\Lambda^{(0)}_\star:=\displaystyle\max_{1\le l\le\mathsf L-1}\ \int_{\mathcal T}\Big(\sup_{u\in\IR}\big\|\ell^\Pi_\star(u,t,l)\big\|\Big)\,\bar w(t)\,\Pi_t(dt)<\infty$;
\item[(i.b)] for every $l\in\{1,\dots,\mathsf L-1\}$ and $\Pi_t$-almost every $t$ the map
$u\mapsto\ell^\Pi_\star(u,t,l)$ is continuously
differentiable, and
$\Lambda^{(1)}_\star:=\displaystyle\max_{1\le l\le\mathsf L-1}\ \int_{\mathcal T}\Big(\sup_{u\in\IR}\big\|\partial_u\ell^\Pi_\star(u,t,l)\big\|\Big)\,\bar w(t)\,\Pi_t(dt)<\infty$;
\end{enumerate}
and $\ell_\star:=\mathcal W\ell^\Pi_\star$ is defined \emph{pointwise} from that version, through
\suppeq{Wdef}, columnwise in $l$;
\item[(iii)] \emph{Anchor.} The debiasing weight attached to $\ell_\star$ is
$L_\star:=L^{\,b_\star}_{\ell_\star}$ of \cref{def:weight} for a measurable
$b_\star:\IR^{p}\to\IR^{p\times(\mathsf L-1)}$ with $E\big\|b_\star(\Delta Z)\big\|^{2}<\infty$, the norm
being Frobenius.
\end{enumerate}
\emph{The two budgets aggregate the $p(\mathsf L-1)$ slicewise quantities in \textbf{two different
ways}}: the $\|\cdot\|$ inside each supremum is the \textbf{Euclidean} norm of the vector
$\ell^\Pi_\star(u,t,l)\in\IR^{p}$, and the aggregation over the contrast index $l$ is by the
\textbf{maximum}. \emph{The asymmetry is deliberate}, and why each half is what it is is in the
Discussion of the assumptions, under \emph{A$2'$}. At $p=1$ the Euclidean norm on $\IR$ is the modulus,
and at $\mathsf L=2$ the maximum is over a one-element set.
\end{assumption}

% DISPATCH-25 TYPING OF A2', AND THE ONE CHOICE IT MAKES.
%
% WHAT MOVED: the contrast index l joins (u,t) as a third argument of ell^Pi,
% and b_star's CODOMAIN becomes R^{p x (sans-L - 1)}.  Its DOMAIN is untouched:
% that R^p is Delta Z's space and has nothing to do with sans-L.  The p in the
% codomain is theta's dimension and is also untouched; what is new is the second
% factor, one column per contrast.  AT sans-L = 2 BOTH DISPLAYS ARE THE PRINTED
% ONES: max over a one-element set is that element, ell^Pi_star(u,t,1) is
% ell^Pi_star(u,t), R^{p x 1} is R^p, and Frobenius is Euclidean.
%
% THE CHOICE, AND IT IS FLAGGED RATHER THAN BURIED.  The budgets Lambda^{(0)}_star
% and Lambda^{(1)}_star are single CONSTANTS everywhere they are consumed --
% cor:elldiff bounds ||ell_star||_inf by ||varpi_u||_inf Lambda^{(0)}_star, and
% lem:reprsmooth produces one number -- so the p(sans-L - 1) slicewise budgets
% have to be aggregated into one.
%
% *** AMENDED 2026-08-09, AUTHOR DECISION, AND THIS REVERSES HALF OF DISPATCH 25.
% *** DISPATCH 25 TOOK THE MAXIMUM OVER BOTH INDICES.  Over l that is right and
% *** stands.  OVER j IT WAS A REGRESSION, and the corrected aggregation is
% *** EUCLIDEAN OVER j, MAX OVER l.
% ***   THE GROUND.  Pre-dispatch-25, ell^Pi_star was R^p-valued and the budget
% *** read int (sup_u |ell^Pi_star(u,t)|) bar-w dPi_t with |.| the EUCLIDEAN norm
% *** on R^p (see the retired mirror assumptions.tex:717-720).  Dispatch 25 wrote
% *** the coordinate index j out and replaced that Euclidean norm by a MAXIMUM
% *** over j.  Every consumer of the two budgets, without exception, reads a
% *** Euclidean or Frobenius norm of the ASSEMBLED object -- cor:elldiff's two
% *** sup-norms, rem:orthoghyp(3), rem:anchor, lem:solset(iii), eq:rprodbound --
% *** so under the max-over-j form each of those constants is SHORT BY sqrt(p),
% *** and the shortfall is visible already at sans-L = 2.  Phase 4 found it,
% *** printed the sqrt(p) at cor:elldiff rather than absorb it, and escalated;
% *** this is the upstream repair the author chose instead.
% ***   WHY NOT EUCLIDEAN OVER l TOO.  eq:Wdef acts DIAGONALLY in l -- slice l
% *** generates column l and no other -- so no downstream bound ever sums across
% *** l; the max is what those bounds read, and it is free of Pi_w, which matters
% *** because A2' prices the representer and B1(iv.d) prices the reference
% *** measure and the block keeps those two apart everywhere else.
% ***   AT p = 1 AND AT sans-L = 2 NOTHING MOVES: a Euclidean norm on R^1 is the
% *** modulus and a maximum over a one-element set is that element, so the
% *** corrected displays are character-identical to the dispatch-25 ones there,
% *** and identical to the PRE-sans-L ones at p = 1, sans-L = 2.
% ***   THE ONE CLAIM DISPATCH 25 MADE THAT DOES NOT SURVIVE, and it was about
% *** sans-L rather than p: "every bound downstream survives with at most a
% *** factor sqrt(sans-L - 1) from passing to the Frobenius norm".  Repair R1
% *** (2026-08-09) put the masses varpi_w(l) inside eq:Wdef, and the l-aggregation
% *** then costs (sum_l varpi_w(l)^2)^{1/2} <= 1 rather than sqrt(sans-L - 1).
% *** R1 PAYS FOR IT; no downstream constant carries sqrt(sans-L - 1).
%
% CONSUMERS RE-READ AND THEIR VERDICTS ARE IN THE PHASE-4 REPORT.  The four sites
% amended here are A2'(i.a), A2'(i.b), A2b and eq:lambdamu, and nothing else in
% this file is touched.  DECLARED AND NOT REPAIRED HERE, both out of scope:
%   * results_s2_identification.tex:2146 (def:weight(a)) says the budgets are
%     aggregated "over j and l by the maximum" and "therefore serve every entry".
%     The first half is now wrong and the second is now WEAKER than the truth --
%     a Euclidean-in-j budget bounds the Euclidean norm of a COLUMN, not merely
%     each entry.  Phase 2a's file.
%   * ~~proofs.tex:2850-2853 and 4050-4087 run the two derivations at the scalar
%     arity; :4086 already writes the MAX form Lambda^{(0)} v Lambda^{(1)} <=
%     c_mu'' Lambda^{(mu'')}, which is what survives.  Phase 5's file.~~
%     STRUCK 2026-08-10 (Phase C).  FALSE IN BOTH HALVES, AND THE POINTERS ARE
%     STALE ON TOP OF THAT.
%       (a) The derivations no longer run at the scalar arity.  Phase 4's
%           close-out deleted lem:reprsmooth's coordinate suppression and ran
%           Steps 0, 1, 1', 2, 3 and 5 at the R^p arity; the budgets there now
%           carry the Euclidean norm on R^p and the contrast index l.
%       (b) The MAX form is NOT what survives.  eq:l21budget
%           (results_supp.tex:2122) prints the SUM form,
%           Lambda^(0)_star + Lambda^(1)_star <= (c^(0) + c^(1)) Lambda^(mu'')
%           <= c_mu''/sqrt2 Lambda^(mu''), which is STRICTLY SHARPER; the MAX
%           form survives only where a proof explicitly weakens to it, and
%           lem:reprsmooth's Step 5 labels it as such ("a weakening of
%           eq:l21budget, each budget being nonnegative and so at most their
%           sum").  Restoring the + form was Phase 4 round 2's finding M1.
%     Cited by LABEL from here on: eq:l21budget and lem:reprsmooth's Steps 1
%     and 5.  Line numbers into proofs.tex from this file have been stale by
%     +308, +315 and more over this track, and a comment nobody compiles is
%     exactly where such a pointer rots unnoticed.

\begin{assumption}[A2b: Admissible representer (a \emph{property}, not maintained)]\label{ass:admrep}
A function $\ell^\Pi\in L^{2}(\Pi)$ is an \emph{admissible representer} if it admits a version satisfying
A2$'$(i.a)--(i.b) with $\ell^\Pi$ in place of $\ell^\Pi_\star$, that is, with
\[
\begin{aligned}
\Lambda^{(0)}&:=\max_{1\le l\le\mathsf L-1}\int_{\mathcal T}\Big(\sup_{u\in\IR}\big\|\ell^\Pi(u,t,l)\big\|\Big)\bar w(t)\,\Pi_t(dt)<\infty,\\
\Lambda^{(1)}&:=\max_{1\le l\le\mathsf L-1}\int_{\mathcal T}\Big(\sup_{u\in\IR}\big\|\partial_u\ell^\Pi(u,t,l)\big\|\Big)\bar w(t)\,\Pi_t(dt)<\infty,
\end{aligned}
\]
the norms and the aggregation over $l$ being A$2'$(i.a)--(i.b)'s,
and $u\mapsto\ell^\Pi(u,t,l)$ continuously differentiable for every $l\in\{1,\dots,\mathsf L-1\}$ and
$\Pi_t$-almost every $t$, and if
$\ell:=\mathcal W\ell^\Pi$ is defined \emph{pointwise} from that version through \suppeq{Wdef},
columnwise in $l$.
\emph{A2b is NOT MAINTAINED}: it is a hypothesis a result may impose on a representer it
quantifies over. Under A2$'$(i) the selected solution $\ell^\Pi_\star$ satisfies
A2b, with $\Lambda^{(0)}=\Lambda^{(0)}_\star$ and $\Lambda^{(1)}=\Lambda^{(1)}_\star$; that
implication is the \emph{only} relation between the two tags.
\end{assumption}

\begin{assumption}[A2c: Source condition of index $\mu''$ on the representer (a \emph{property}, not maintained)]\label{ass:reprsource}
Fix $\mu''>\tfrac32$. For $g\in L^{2}(\IR,du)$ write $(\mathcal Fg)(\xi):=\int_{\IR}g(u)e^{-\mathrm iu\xi}du$
for the Fourier transform in $u$, and let
\begin{equation}\label{eq:usobolev}
H^{\mu''}(\IR)\;:=\;\Big\{g\in L^{2}(\IR,du):\
\big\|g\big\|^{2}_{H^{\mu''}(\IR)}:=\tfrac{1}{2\pi}\!\int_{\IR}\!\big(1+\xi^{2}\big)^{\mu''}\big|(\mathcal Fg)(\xi)\big|^{2}d\xi<\infty\Big\}
\end{equation}
be the $u$-Sobolev scale of order $\mu''$ on $\IR$, \emph{taken with respect to Lebesgue measure $du$ and
not with respect to $\Pi_u$}. For an $\IR^{p}$-valued $g$ on $\IR$ write
$H^{\mu''}(\IR;\IR^{p})$ for the same scale with $|\cdot|$ replaced inside \eqref{eq:usobolev} by the
Euclidean norm $\|z\|:=\big(\sum_{j=1}^{p}|z_j|^{2}\big)^{1/2}$ on $\IC^{p}$, which is where
$(\mathcal Fg)(\xi)$ takes its values, so that
$\|g\|^{2}_{H^{\mu''}(\IR;\IR^{p})}=\sum_{j=1}^{p}\|g_j\|^{2}_{H^{\mu''}(\IR)}$ and
$g\in H^{\mu''}(\IR;\IR^{p})$ if and only if every coordinate lies in $H^{\mu''}(\IR)$;
\emph{at $p=1$ they are the same object}.
A function $\ell^\Pi\in L^{2}(\Pi)$ \emph{satisfies the source condition of
index $\mu''$} if it admits a jointly measurable representative, \emph{fix one; everything below is
read at it}, whose $\mathsf L-1$ slices satisfy: there is a Borel set
$\mathcal T_\ast\subseteq\mathcal T$, \emph{one set, serving every $l$}, with
$\Pi_t(\mathcal T\setminus\mathcal T_\ast)=0$ such that the
slice $\ell^\Pi(\cdot,t,l)$ lies in $H^{\mu''}(\IR;\IR^{p})$ for \emph{every} $t\in\mathcal T_\ast$ and
every $l\in\{1,\dots,\mathsf L-1\}$, and
\begin{equation}\label{eq:lambdamu}
\Lambda^{(\mu'')}\;:=\;\max_{1\le l\le\mathsf L-1}\ \int_{\mathcal T_\ast}\big\|\ell^\Pi(\cdot,t,l)\big\|_{H^{\mu''}(\IR;\IR^{p})}\;\bar w(t)\,\Pi_t(dt)\;<\;\infty ,
\end{equation}
the budget being aggregated as in A2$'$(i).
% EXTRACTED FROM THE STATEMENT OF A2c ON 2026-08-03.  Seven lines stood here
% saying what A2c IS -- not maintained, a hypothesis lem:reprsmooth imposes on a
% representer it quantifies over, its only role, and the geometry it is stated in.
% Status and motivation are not conditions, so they are now rem:reprsourceread.
% WHAT DELIBERATELY STAYED, because it disambiguates rather than comments: "taken
% with respect to Lebesgue measure du and not with respect to Pi_u", and "fix one;
% everything below is read at it" for the jointly measurable representative.  Both
% fix WHICH object the clause is about, so a reader who drops them reads a
% different condition.
% NO \label MOVED: eq:usobolev and eq:lambdamu both stayed in the statement.
\end{assumption}

\noindent\emph{Notes on A2c.} That it is a property and not a maintained condition, whose sole role is to
imply A2$'$(i) at the starred representer, and the $\bar w$-weighted geometry it is stated in, are \suppref{Remark}{reprsourceread}.

% ---------------------------------------------------------------------------
% WHY A2c EXISTS, WHY IT IS NOT MAINTAINED, AND WHAT IT DOES NOT DO.
%
% THE OBLIGATION IT DISCHARGES.  (O1) of rem:obligations named L-21 as the
% discharge site for A2'(i.a)-(i.b) IN SUP-INSIDE FORM ON ell^Pi_star.  L-21 is
% now lem:reprsmooth, and A2c is the condition it runs on.  A2'(i) REMAINS A
% MAINTAINED ASSUMPTION; what changes is that it is no longer an unverifiable
% primitive -- it is implied by one smoothness condition on one object, stated in
% the language the block already uses for varphi_dagger (H1, a Hilbert scale).
%
% WHY NOT THE ROUTE THE MEMO REGISTERS.  proof_strategy_memo.md's L-21 row
% proposes "a source condition of index mu'' ... giving ell_star in
% R((A A*)^{mu''}) -> C^1".  THAT ROUTE IS THREE STEPS SHORT, and each gap is
% recorded here rather than papered over:
%   (a) WRONG REFERENT.  It targets ell_star, the OBSERVABLE-space object.  The
%       round-5 repair moved A2'(i) onto the RANGE-space object ell^Pi_star
%       (see clause (c) of the long note below), so a condition on ell_star
%       cannot discharge it -- and cor:elldiff, which bridges the two, CONSUMES
%       A2'(i) and so cannot be used to get there.
%   (b) NO SCALE LINK ON THE RANGE SIDE.  R((A_cond A*_cond)^{mu''}) is a
%       subspace of L^2(Pi) defined by the singular system of A_cond.  Membership
%       in it says NOTHING about sup_u or about u-differentiability unless the
%       A_cond A*_cond scale is MATCHED to a u-smoothness scale.  H1 supplies
%       exactly such a matching on the DOMAIN side (Sobolev order s generated by
%       D, matched to sigma_j asymp j^{-a}); there is no analogue on the range
%       side, and inventing one would be a SECOND Hilbert-scale link, assumed and
%       not derived, about how the left singular functions tilde-e_j oscillate in
%       u.  That is a strictly larger assumption than A2c, on an object the
%       analyst does not choose.  ROUTE REJECTED, and the reason is recorded
%       because it is the obvious thing a reader will propose.
%   (c) NO t-INTEGRATION.  (i.a) and (i.b) are integrals over T of a sup over u,
%       against bar-w d Pi_t.  The memo's row states only the embedding and omits
%       the t-integration entirely, which is half of what the clauses assert.
%       eq:lambdamu supplies it, in the SAME geometry.
%
% WHAT A2c COSTS BEYOND A2, STATED PLAINLY.  ell^Pi_star in L^2(Pi) gives, for
% a.e. t, ell^Pi_star(., t) in L^2(Pi_u) = L^2(varpi_u du).  B1(iv.a) makes
% varpi_u bounded and strictly positive POINTWISE, NOT bounded below, so
% L^2(du) is contained in L^2(Pi_u) and NOT conversely: a slice can lie in
% L^2(Pi_u) with fat tails and fail to lie in L^2(du).  A2c therefore adds a
% TAIL-DECAY requirement in |u| on top of smoothness.  That is the same species
% of condition as C4''s residual truncation, and it is stated rather than hidden.
% The asymmetry is the mirror image of the one recorded at A2 for pi: a bound on
% a density in one direction does not give the reverse containment.
%
% WHY A2c IS CHECKABLE, WHICH IS THE WHOLE POINT.  By eq:usobolev, A2c is a
% weighted summability condition on the u-Fourier transform of the slices -- the
% same KIND of object as A2's Picard sum, and reportable in the application in
% the same way (plan Sec.9.2 item 3).  Picard summability in the A_cond system
% cannot supply it (see item 2 of the FOUR THINGS note below); a summability
% condition in the u-frequency variable can, and does.
%
% mu'' IS INDEPENDENT OF mu.  mu = s/(2a) is a smoothness index for
% varphi_dagger on X; mu'' is a smoothness index for ell^Pi_star in u.  Nothing
% links them, and the smoothness-budget table lists them separately.
%
% THE INTEGER CASE NEEDS NO TRANSFORM, and rem:l21route records it: at mu'' = 2,
% H^2(R) is {g, g', g'' in L^2(du)} and eq:embed holds ~~with c_2 = 2~~ WITH THE
% CONSTANT 2 by an elementary argument.  TWO CORRECTIONS TO THE ROUND-9 FORM OF
% THIS NOTE, both from the dispatch-6 review.
%   (1) THE SYMBOL WAS MISUSED.  c_2 is eq:embedconst at mu'' = 2, and its value
%       is SQRT(2), since int (1+xi^2)^{-1} dxi = pi.  The elementary route gives
%       the WEAKER constant 2.  Writing "c_2 = 2" put two different numbers on
%       one symbol.  Both constants are correct in their own route; the sharper
%       one is c_2 = sqrt(2) and it is the one eq:l21budget is stated with.
%   (2) THE H^2 NORM CARRIES A FACTOR 2 ON ||g'||.  Under eq:usobolev the weight
%       at mu'' = 2 is (1 + xi^2)^2 = 1 + 2 xi^2 + xi^4, so
%       ||g||^2_{H^2} = ||g||^2 + 2||g'||^2 + ||g''||^2 -- eq:h2norm.  The SET is
%       {g, g', g'' in L^2(du)} either way, the two quadratic forms being
%       equivalent with ratio in [1, 2]; the NORM is not.  Step 1' of
%       lem:reprsmooth's proof is correct against the corrected norm, so the
%       comparison sqrt(2) < 2 is like for like.
% An applied reader who does not want the Fourier
% definition can read A2c at mu'' = 2 and lose only mu'' in (3/2, 2).
% ---------------------------------------------------------------------------

% ---------------------------------------------------------------------------
% WHY A2b IS A SEPARATE TAG (round-8 fix, auditor Finding 4 -- IT GATED STEP 4).
%
% THE TAG LETTER WAS SWEPT BEFORE IT WAS CLAIMED, which is the round's own
% organising discipline.  MY FIRST CHOICE WAS "R1", following the auditor's
% "N1' treatment" wording.  R1 IS ALREADY OCCUPIED: results.tex:270 carries
% \begin{remark}[R1]\label{rem:R1}, R1 being a RESULT name from the plan (as are
% L-16, R2, O-11).  Claiming it would have reproduced, inside the fix for
% Finding 4, exactly the defect the warden's structural observation describes --
% a resolution that reasons about the symbol it renames and not about the slot it
% claims.  A2b is used instead: letter-suffixed A-family tags are the project's
% established device for a related-but-distinct condition (A1b, A1c, A1d), the
% string "A2b" occurs nowhere else in any file, and it is visually separable from
% A2' in a way that "N2" would not have been from the memo's escalation items
% N-2, N-7, N-10, N-18.
%
% THE DEFECT.  def:weight(a) defined "admissible representer" by re-using the
% TEXT of A2'(i.a)-(i.b) with ell^Pi in place of ell^Pi_star.  The tag A2'(i)
% therefore named TWO objects at once:
%   - a MAINTAINED ASSERTION about the selected solution (consumed by
%     cor:elldiff at ell^Pi_star, lem:solset, lem:secondord, prp:orthog(ii)-(iii),
%     lem:balance);
%   - a PROPERTY PREDICATE over generic representers (consumed by def:weight(a),
%     hence lem:solset(ii), prp:orthog(i), eq:PsiL).
% N1 got two things this did not: an explicit status sentence, and a SEPARATE TAG
% (N1') when applied to a different object.  A2'(i) got neither.
%
% WHY IT GATED STEP 4.  L-21's target is A2'(i.a)-(i.b) IN SUP-INSIDE FORM ON
% ell^Pi_star.  With the same text doing duty as a predicate over generic
% representers, and cor:elldiff's hypothesis line naming ell^Pi_star while FOUR
% sites applied it generically (results.tex def:weight(a) and eq:PsiL;
% proofs.tex lem:solset(ii) and prp:orthog(i)), a Step-4 author could not read
% off WHICH OBJECT L-21 discharges.  Splitting the tag makes that mechanical:
% L-21 discharges A2'(i) at ell^Pi_star and nothing else; A2b is what a result
% ASSUMES about a representer it quantifies over, and is never discharged by
% L-21 because it is never asserted.
%
% CONSEQUENCE FOR cor:elldiff.  The corollary is now STATED at a generic
% admissible representer under A2b, with ell^Pi_star as the SPECIALISATION.  That
% is a strictly stronger result on strictly weaker hypotheses -- the same
% move prp:orthog(i) received in round 7 -- and it removes the four generic
% applications' reliance on a hypothesis line that named ell^Pi_star.
%
% ORDERING NOTE.  A2b is stated HERE, in the assumption block, and NOT by
% forward reference to def:weight(a), because cor:elldiff (results.tex) precedes
% def:weight in the file and may not cite it.  def:weight(a) now DEFINES the
% phrase "admissible representer" TO MEAN A2b, which is a naming convention and
% not a second statement of the condition.
%
% ROUND-9 CORRECTION TO THIS COMMENT BLOCK, FOUND BY SWEEPING AND NOT BY A
% FINDING.  Three lines above read "R1" where they meant A2b -- the tag this very
% block records as REJECTED, on the ground that R1 already names rem:R1.  The
% round-8 rename was applied to every statement and every citation and NOT to the
% prose that explains the rename, so the block declared a slot unclaimable in one
% paragraph and claimed it in the next.  That is glossary Sec.1.5b's defect
% species one level in -- the sweep for the letter being CLAIMED was run over the
% .tex statements and not over the comment that governs them -- and it is
% corrected in place rather than left as a comment "nobody compiles", because a
% comment block in the file of record is what the next author reads.
% ---------------------------------------------------------------------------

% ---------------------------------------------------------------------------
% CLAUSE (iii) IS NEW IN ROUND 7, AND IT CLOSES A CRITICAL: L_star WAS USED
% EVERYWHERE AND DEFINED NOWHERE.
%
% THE DEFECT.  Through round 6 the symbol L_star appeared in eq:Psitau,
% lem:secondord, prp:orthog(ii)-(iii), A5(iv) and the glossary's eps_n row, and
% was DEFINED IN NO FILE.  The only thing resembling a definition was the aside
% "W ell^Pi_star = partial_u L_star" inside a proof step, which makes L_star an
% ANTIDERIVATIVE IN u OF ell_star WITH A FREE CONSTANT OF INTEGRATION b(Delta Z).
% Nothing constrained b.  Under N1 the gap closes in two lines -- N1 pins b --
% BUT prp:orthog DELIBERATELY WITHHOLDS N1, which is the whole content of its
% part (iii).  With N1 withheld and b arbitrary: eq:Psitau's expectation at tau = 0
% need not exist; E[L_star | Delta Z] on the right of prp:orthog(iii)'s
% equivalence need not exist; and the "=>" direction CONSTRUCTS ITS PERTURBATION
% OUT OF THAT NON-EXISTENT OBJECT, r = eps_varrho sign(E[L_{star,j} | Delta Z]).
%
% THE REPAIR, IN TWO PARTS.  def:weight in results.tex DEFINES the weight
% attached to a representer as L^b_ell(u,z) = int_0^u ell(v,z) dv + b(z), an
% anchor b and a representer ell being exactly the data that determine it.
% Clause (iii) then prices the anchor, and NOTHING ELSE.
%
% WHY THIS IS THE WEAKEST FORM THAT CARRIES THE PROOF, AND WHAT MORE WOULD BUY.
% By cor:elldiff, ||ell_star||_inf <= ||varpi_u||_inf Lambda^(0)_star < infinity
% under (i).  Hence, for EVERY u and z,
%       ||L_star(u,z)||  <=  ||b_star(z)||  +  ||ell_star||_inf |u| ,
% so that, with A1b(iv)'s E[Delta U^2] < infinity and lem:nusigma(i)'s
% E[(Delta h)^2] <= 4 bar-c_pi ||h||^2,
%   (a) E ||L_star(Delta U - tau_1 Delta h, Delta Z)|| < infinity for EVERY
%       tau_1 -- no shift restriction is needed on the L_star term, which is why
%       lem:secondord's shift proviso is imposed on the DIRECTION D and not here;
%   (b) E[ ||L_star|| | Delta Z ] < infinity a.s., which is what makes
%       prp:orthog(iii)'s conditional expectation and its sign perturbation
%       well defined;
%   (c) E[ Wtilde^2 ||L_star||^2 ] <= 2 E||b_star||^2
%                                   + 2 ||ell_star||_inf^2 E[Delta U^2] < infinity,
%       which is the VARIANCE of the moment -- see the correction under A4(d).
% Anything stronger -- L_star in L^2(P) as a standing convention, or a sup-norm
% bound on L_star -- would buy NOTHING beyond (a)-(c) and would be false for the
% affine weights of cor:udep, whose L(u,z) = c(z) u is unbounded in u.
%
% AND IT COSTS NOTHING AT EITHER ANCHOR OF INTEREST, WHICH IS THE POINT.
%   canonical anchor  b_star == 0:  E||b_star||^2 = 0.
%   N1 anchor  b_star(z) = -E[ int_0^{Delta U} ell_star(v,z) dv | Delta Z = z ]:
%       ||b_star(z)|| <= ||ell_star||_inf E[ |Delta U| | Delta Z = z ], so by
%       conditional Jensen E||b_star||^2 <= ||ell_star||_inf^2 E[Delta U^2]
%       < infinity by A1b(iv).
% So (iii) BINDS ONLY on an anchor the analyst has chosen badly.  It is stated
% because the anchor is a free design object, not because either canonical choice
% is in doubt.
%
% VERSIONS (m-3).  ell^Pi_star is an equivalence class in L^2(Pi), and neither
% "sup_u |ell^Pi_star(u,t)|" nor "continuously differentiable in u" is invariant
% under Pi-null modification -- while cor:elldiff's conclusions are evaluated at
% the RANDOM point (Delta U, Delta Z).  A2' therefore says "admits a version"
% and fixes ell_star pointwise from that version.  Without this the sup-norm
% clauses are not well posed as stated.
%
% VECTOR NORM IN (ii).  (ii) was written E[Wtilde^2 ell_star^2] for an
% R^p-valued ell_star.  It is now E[Wtilde^2 ||ell_star||^2], which is the
% coordinatewise statement summed and is well typed; nothing downstream changes.
% ---------------------------------------------------------------------------
%
% ---------------------------------------------------------------------------
% ROUND-2's CLAUSE (0) IS DELETED, BECAUSE IT IS NOW A THEOREM.  Round 2 added
% an "evaluability" clause asserting ell_star in closure(G_gamma), so that
% ell_star(Delta U, Delta Z) would be a well-defined random variable.  The
% clause repaired a real type collision -- a solution of A*_cond ell = -nu/pi
% lives on (u,t) with t the BIERENS INDEX, while (i), (ii) and eq:JSig (through
% dispatch 12: "A3") evaluate
% ell_star at (u,z) with z a VALUE of Delta Z -- but it repaired it by
% ASSERTING the referent instead of constructing it.  Escalated as N-10.
%
% lem:adjoint(ii) NOW CONSTRUCTS THE MAP.  Under the strengthened B1(iv.a) the
% Bierens induction operator
%     (W ell^Pi)(u,z) := varpi_u(u) int_T ell^Pi(u,t) w(t'Phi(z)) Pi_t(dt)
% is a BOUNDED linear map L^2(Pi) -> L^2(P_{Delta U, Delta Z}) with dense range,
% and lem:adjoint(iii) shows that it COMMUTES WITH THE ADJOINT EQUATION:
%     A*_cond ell^Pi = nu_{Wtilde (W ell^Pi)} / pi.
% So for EVERY solution ell^Pi of the adjoint equation the observable-space
% object ell_star := W ell^Pi is automatically square integrable and
% ell_star(Delta U, Delta Z) is automatically a well-defined random variable.
% An assumption that has become a theorem must leave the block; leaving it in
% would show up in the usage matrix as an unused-or-redundant clause and would
% overstate what is assumed.  A2'(i) and A2'(ii) KEEP THEIR NUMBERS, so every
% downstream citation of A2'(i) or A2'(ii) is unaffected.
%
% THREE THINGS THE CONSTRUCTION DOES NOT DO, recorded so that Step 4 does not
% inherit them silently.
%   (a) W IS NOT SURJECTIVE.  Its range is dense (lem:adjoint(ii)) but need not
%       be all of L^2(P_{Delta U, Delta Z}), so a variance-minimising selection
%       posed over the full space may be unreachable.  That is prp:effbound's
%       problem (Step 8), not lem:adjoint's.
%   (b) W IS NOT INJECTIVE in general; what IS proved is
%       ker W subset ker A*_cond, which is all the selection argument of L-16
%       needs, since a direction killed by W is killed by the adjoint too.
%   (c) CLAUSE (i) IS NOW STATED ON THE RANGE-SPACE OBJECT, AND THAT IS A
%       ROUND-3 REPAIR, NOT A RESTATEMENT.  Rounds 1-2 wrote (i) as
%       ||ell_star||_inf < infinity AND ||partial_u ell_star||_inf < infinity.
%       AFTER N-10 CLOSED, THAT WAS ILL-POSED: ell_star = W ell^Pi_star is
%       CONSTRUCTED, and nothing in the block made partial_u ell_star EXIST.
%       B1(iv.a)'s C^1 clause bounds only the varpi_u' term of eq:Wderiv; the
%       second term needs partial_u ell^Pi to exist with a dominating function,
%       WHICH B1(v) DOES NOT SUPPLY -- B1(v) dominates partial_u lambda in the
%       L-representation, and since W ell^Pi_star = partial_u L_star, that is one
%       derivative short of what eq:Wderiv consumes.
%
%       THE REPAIR PUTS THE CONDITION WHERE IT BELONGS.  Differentiability is a
%       property of the SELECTED SOLUTION, not of an object the analyst chooses,
%       so it cannot go in B1 at all.  (i.a) and (i.b) are conditions on
%       ell^Pi_star; cor:elldiff then DERIVES
%           ||ell_star||_inf        <= ||varpi_u||_inf Lambda^(0)_star,
%           ||partial_u ell_star||_inf
%              <= ||varpi_u'||_inf Lambda^(0)_star + ||varpi_u||_inf Lambda^(1)_star,
%       together with eq:Wderiv itself.  What used to be assumed is now a
%       conclusion, and (i.b) is exactly the dominating function that licenses
%       differentiating under the Pi_t-integral.
%
%       B1(v) IS NOT AMENDED.  Its second half remains the dominating function
%       for the GENERIC L-representation quoted in rem:Tcond; at the SELECTED
%       ell_star the domination now comes from (i.b).  Adding a second-derivative
%       clause to B1(v) was the alternative and was rejected: it would impose a
%       smoothness condition on lambda, an object the analyst does not choose
%       once lem:adjoint(ii) fixes gamma = Pi_t and partial_u lambda =
%       varpi_u ell^Pi.  THE TWO SYMBOLS ARE RELATED BY ONE u-INTEGRATION AND
%       ARE NOT INTERCHANGEABLE; C13's reservation of lambda stands.
%
%       THE SUPREMUM IS INSIDE THE INTEGRAL, AND THAT IS A ROUND-5 CORRECTION.
%       Round 4 wrote (i.a) and (i.b) as  sup_u int (.) dPi_t < infinity.  THAT
%       DOES NOT DOMINATE.  Leibniz's rule needs int sup_u (.) dPi_t < infinity;
%       the sup-outside form gives L^1(Pi_t)-boundedness in u WITHOUT uniform
%       integrability, and cor:elldiff's proof silently used the stronger one.
%       Two repairs were available:
%         (alpha) move the sup inside -- STRICTLY STRONGER, since
%                 sup_u int <= int sup_u, so the round-4 clauses are implied and
%                 every printed conclusion survives with the same constants;
%         (beta)  keep the sup outside and weaken the conclusion to: ell_star is
%                 LIPSCHITZ in u with constant Lambda^(1)_star, and eq:Wderiv
%                 holds for a.e. u rather than every u.
%       (alpha) IS ADOPTED, AND THE DECISIVE REASON IS NOT THE ONE FIRST GIVEN.
%       Round 5 justified (alpha) by pointing at lem:secondord(i) and cor:udep,
%       which evaluate partial_u ell_star at Delta U, a RANDOM point.  cor:udep
%       is a weak instance: its use of partial_u ell_star is at the FUNCTION
%       level (whether it is identically zero), not at a random argument.
%       THE DECISIVE CONSUMER IS THE C2'(i) REMOVAL.  Under (beta), eq:Wderiv
%       holds only for a.e. u, so using it at Delta U requires Delta U to charge
%       no Lebesgue-null set -- i.e. requires an absolute-continuity condition,
%       which is exactly what C2'(i) supplies and what this round REMOVED from
%       lem:secondord and prp:orthog as idle.  (beta) would put it straight
%       back.  SO (alpha) AND THE C2'(i) REMOVAL STAND OR FALL TOGETHER, and
%       choosing (beta) would have undone a separate correction in the same
%       round.  A2'(i) is a sup-norm condition in any case, so tightening it in
%       the sup direction is in keeping.
%
%       WHAT REMAINS OPEN, AND THE TARGET HAS BEEN RAISED.  (i.a) and (i.b) are
%       still primitives.  L-21 -- a source condition of index mu'' on
%       ell^Pi_star plus a Sobolev embedding -- is the route to them and is
%       ~~unproved~~ **PROVED IN ROUND 9 as lem:reprsmooth, on the condition A2c.**
%       **L-21 MUST NOW DELIVER THE SUP-INSIDE FORM.**  theory_memo.md
%       Sec.D5 item 2 previously described L-21's target as the sup-outside
%       form; that is the form that does NOT dominate, so aiming L-21 at it would
%       leave cor:elldiff unproved after Step 4 had apparently discharged it.
%       Recorded before Step 4 runs, which is when it is cheap.
%       **AND IT IS THE SUP-INSIDE FORM THAT lem:reprsmooth DELIVERS**: A2c is
%       stated in the bar-w L^1(Pi_t) geometry of (i.a)-(i.b) themselves, and
%       lem:reprsmooth(ii) is a SLICEWISE inequality, so the sup never has to be
%       moved across the t-integral at all.  The raised target is met, not
%       approximated.
%       The gain over rounds 1-2 is unchanged: the clauses are now WELL FORMED
%       and are stated on an object whose existence the block guarantees.
%
% ONE CONSEQUENCE FOR STEP 4.  By eq:Wderiv the DESIGN CHOICE Pi_u enters
% partial_u ell_star, hence q, hence BOTH sigma_q AND the operator K with kernel
% q f_{X_1,X_2}, hence the second-order form of lem:secondord(i).  Pi_u is
% therefore a THIRD selection lever alongside the choice of ell^Pi within its
% affine set, and SG-15 is where it is swept.
%
% DO NOT WRITE "(1/2)||sigma_q/pi||_inf IS prp:sharpquad's SHARP CONSTANT."  An
% earlier version of this comment did.  FOUR things are wrong with it and all
% four are recorded at lem:solset(iii):
%   (a) sigma_q is R^p-valued (F-15), so a NORM presupposes a direction
%       vartheta in R^p with ||vartheta|| = 1: the object is
%       ||vartheta' sigma_q / pi||_inf, never ||sigma_q/pi||_inf.  (The direction
%       letter is vartheta and NOT varsigma; varsigma is a reserved LOCAL bound
%       variable and round 6 wrongly made it binding.  Round-7 rename.)
%   (b) the form on L^2(pi) is B = (vartheta' sigma_q/pi). - 2 Ktilde, and B is
%       NOT self-adjoint, since q f_{X_1,X_2} is not a symmetric kernel;
%   (c) a quadratic form sees only the symmetric part, so the sharp constant is
%       (1/2)||B_sym|| with B_sym = (B + B*)/2, and (1/2)||B|| is only an UPPER
%       bound;
%   (d) ||vartheta' sigma_q/pi||_inf is a LOWER bound on ||B_sym||, via Weyl's
%       theorem and the essential spectrum -- useful, and it is what makes the
%       form non-compact, but it is neither the constant nor an upper bound.
%
% AND THE SPACE IS L^2(pi), NOT V -- ROUND-7 CORRECTION, AND IT GATES STEP 4.
% def:spaces says in italics that every operator in this paper has domain V.  B
% IS THE ONE DECLARED EXCEPTION, and it has to be, because V = {h : int h pi = 0}
% is NOT INVARIANT under multiplication by vartheta' sigma_q / pi.  Three
% consequences that Step 4 must not inherit silently:
%   (i)  ||vartheta' sigma_q / pi||_inf < infinity is not automatic and is now
%        proved, not assumed: |sigma_q(x)| <= ||partial_u ell_star||_inf
%        (f_{X_1}(x) + f_{X_2}(x)) by (P1), so
%        ||vartheta' sigma_q/pi||_inf <= 2 bar-c_pi ||partial_u ell_star||_inf,
%        finite by C1(ii) and cor:elldiff.  Without it the essential-spectrum
%        statement would be about an UNBOUNDED operator.
%  (ii)  Weyl and the norm ||B_sym|| are statements on ALL of L^2(pi), whereas
%        lem:secondord(i)'s quadratic form is evaluated only at h in V, so the
%        form's sharp constant is (1/2)|| P_V B_sym P_V |_V ||.  THE TWO ARE NOT
%        THE SAME OPERATOR.
% (iii)  The lower bound transfers anyway, and lem:solset(iii)'s proof now says
%        why: V has codimension one, so B_sym - P_V B_sym P_V has rank at most
%        two, a finite-rank perturbation changes no essential spectrum, and
%        sigma_ess(P_V B_sym P_V |_V) = sigma_ess(B_sym).
% BINDING ON STEP 4.
%
% NOTE the status of eq:Wderiv: it is a CONDITIONAL identity, and (i.b) is its
% hypothesis, not a theorem.  Calling the lever "delivered" overstates it.

% FOUR THINGS ABOUT A2'.
%
% 1. (ii) IS IMPLIED BY (i).  Since ||Wtilde||_inf <= 1 -- which is free from
%    W taking finitely many values and varrho_l in [0,1], and needs NO overlap
%    condition -- (i) gives E[Wtilde^2 ||ell_star||^2] <= || ||ell_star|| ||_inf^2
%    < infinity.
%    The memo lists both without noting the implication.  Both are retained
%    deliberately: tier T3 (the trimmed estimator, memo Sec.6.6) RELAXES (i)
%    and keeps (ii), so (ii) must exist as a separate, citable clause.
%    Recorded in theory_memo.md Sec.6, finding F-1.
%    DISPATCH-25 FORM OF THE BOUND, AND THE CONSTANT IS NOT 1 AT sans-L >= 3.
%    Componentwise, Wtilde_l = 1{W = l} - varrho_l lies in [-1, 1] whatever
%    sans-L is, so ||Wtilde||_inf <= 1 SURELY at every sans-L, and THAT is the
%    (P1) bound that generalises verbatim.  The EUCLIDEAN norm does not: at
%    W = k >= 1, ||Wtilde||_2^2 = (1 - varrho_k)^2 + sum_{l != k, l >= 1}
%    varrho_l^2 <= 1 + 1 = 2, and at W = 0 it is sum_{l>=1} varrho_l^2 <= 1, so
%    ||Wtilde||_2 <= min{sqrt 2, sqrt(sans-L - 1)}.  AT sans-L = 2 THAT MINIMUM
%    IS 1 and the printed bound |Wtilde| <= 1 is recovered exactly; at
%    sans-L >= 3 it is sqrt 2 and the implication above carries the factor 2,
%    which is a constant and not a rate.  THE BOUND IS UNIFORM IN sans-L -- it
%    does NOT grow with the number of cells -- which is worth recording because
%    the naive "sans-L - 1 terms each bounded by 1" reading would suggest it
%    does.
%
% 2. PICARD CANNOT SUPPLY (i).  The Picard criterion in A2 delivers only
%    ell^Pi_star in L^2(Pi) -- an L^2 membership statement.  It says nothing
%    about sup-norms, and (i) is precisely a sup-norm statement.  The route to
%    (i) (memo L-21, ~~to be proved at Step 4~~ **PROVED at Step 4, round 9:
%    lem:reprsmooth on the condition A2c**) is a SOURCE CONDITION OF INDEX mu''
%    ON ell^Pi_star, followed by a Sobolev embedding valid once mu'' exceeds the
%    embedding threshold for the relevant scale -- **which is mu'' > 3/2, the
%    threshold for H^{mu''}(R) -> C^1_b(R) in one dimension**.  mu'' is genuinely
%    independent of mu: mu is a smoothness index for varphi_dagger, mu'' one for
%    ell^Pi_star.  ROUND-7 REFERENT CORRECTION: this note, the smoothness-budget
%    table and the glossary all said "mu'' for ell_star" after the round-5 repair
%    had moved (i) onto the RANGE-SPACE object ell^Pi_star.  L-21 must deliver
%    (i.a)-(i.b), which are conditions on ell^Pi_star, so mu'' is an index for
%    ell^Pi_star.  **ROUND-9: the glossary row for mu'' was the last uncorrected
%    instance of that defect and is now fixed; the round-7 correction reached this
%    note and the budget table and not the row.  Found by sweeping the referent,
%    not by a finding.**  Without L-21, (i) is an unverifiable primitive and the
%    entire second-order analysis rests on it -- which is why it is stated as an
%    assumption here and discharged, not asserted.
%    WHAT THE DISCHARGE DOES AND DOES NOT CHANGE.  A2'(i) is STILL MAINTAINED and
%    still carried as a hypothesis by every consumer.  What lem:reprsmooth removes
%    is its status as an UNVERIFIABLE primitive: it is now implied by one
%    smoothness-plus-tail condition on one object, in a language the block already
%    uses (H1's Hilbert scale) and with a reportable diagnostic.  A2c does NOT
%    become maintained, on the A1d model.
%
% 3. (iii) IS PRICED SEPARATELY FROM (ii) AND THE TWO ARE NOT THE SAME OBJECT.
%    (ii) is about ell_star = partial_u L_star; (iii) is about L_star itself,
%    through its anchor.  Neither implies the other: an L^2 bound on a function
%    does not bound its derivative (lem:levels(iv)), and a sup bound on a
%    derivative does not bound the function (the anchor is free).  The moment's
%    VARIANCE is governed by E[Wtilde^2 ||L_star||^2] and therefore by (i)+(iii),
%    not by (ii).  See the round-7 correction under A4(d).
%
% 4. A2'' IS NOT WRITTEN HERE, deliberately.  The round-2 range-inclusion
%    condition R(B) subset R(A*A) is empty unless sigma_q == 0 (memo Sec.3.5,
%    verified by the strategist-critic).  Do not reintroduce it.
%
\begin{assumption}[H1: Hilbert-scale link]\label{ass:hilbscale}
$\mathcal A_{\mathrm{cond}}:\mathcal V\to L^{2}(\Pi)$, read into the \emph{whole} three-factor range
space of \cref{def:spaces}, not contrast by contrast, has infinite rank and its singular values satisfy
$\sigma_j(\mathcal A_{\mathrm{cond}})\asymp j^{-a}$ for some $a>0$; the source set for $\vdagger$ is a
ball of the \emph{classical} Sobolev space $H^{s}(\mathcal X)$, for non-integer $s$, the interpolation
space between the neighbouring integer orders, \textbf{with no boundary condition imposed at either
endpoint of $\mathcal X$}; and the operator is \emph{matched} to that scale in the sense of
\citet[\S8.5]{engl2000}, which is here the following stated comparison: there is a finite $\mathsf c$,
not depending on $\zeta$, such that every $\zeta\in\mathcal V\cap H^{s}(\mathcal X)$ admits
\begin{equation}\label{eq:h1match}
\zeta\;=\;\big(\mathcal A^{\ast}_{\mathrm{cond}}\mathcal A_{\mathrm{cond}}\big)^{s/(2a)}h
\qquad\text{with }\ \|h\|_{L^{2}(\pi)}\;\le\;\mathsf c\,\|\zeta\|_{H^{s}} .
\end{equation}
\end{assumption}

\noindent\emph{Notes on H1.} That $a$ is an exponent of the \textbf{stacked} operator and not of any one
contrast, that it does not depend on the range weights $\Pi_{\mathsf w}$, what a per-contrast reading of
it does and does not license at $\mathsf L\ge3$, and why $\mathsf L=2$ reduces with no exception, are the
supplement's Discussion paragraph \emph{H1}. The disposition of the two numbers the glyph $\mu$ names is
\suppref{Remark}{pasrc}.

% WHY THE STATEMENT NOW NAMES THE RANGE SPACE, AND IT IS ONE CLAUSE.  Through
% dispatch 25 H1 read "A_cond : V -> L^2(Pi) has infinite rank ...", which is
% CORRECT at every sans-L -- def:spaces makes L^2(Pi) the three-factor space --
% but which a reader arriving from the binary rounds will read as the binary
% operator, since that is the object the same six tokens named through round 11.
% At sans-L >= 3 they name a DIFFERENT OPERATOR, and the difference is a RANGE
% ENLARGEMENT and not the normalisation (which was reversed; the contrast is raw
% on both sides at every sans-L).  The clause is inserted so that the reader is
% told which operator before `a' is quantified over, not after.
% NOTHING ELSE IN THE STATEMENT MOVES.  eq:h1match was ALREADY at the stacked
% Gram operator A*_cond A_cond, so it needed no amendment; the domain V, the
% scale H^s and the constant sans-c are untouched.
%
% WHY THE sans-L >= 3 CONTENT IS IN THE SUPPLEMENT AND NOT HERE.  It is four
% paragraphs -- the Pi_w invariance, the per-contrast pinch a = min_l a_l, the two
% over-readings it invites, and the sans-L = 2 reduction -- and none of it is a
% CONDITION.  Under the house style a maintained assumption carries conditions
% only; and the article is page-budgeted against OBES while the Discussion
% paragraph H1 already carries the rest of this assumption's reading.

% WHY H1 MATTERS, and it does -- BUT ROUND 12 NARROWS THE CLAIM, AND THE
% NARROWING IS THE CRITICAL FINDING OF THAT ROUND.
%
% WHAT H1 STILL DOES, and it is most of what the block uses it for:
%   (a) it FIXES THE SCALE H^s in which def:srcball, A5(vi'), (P-a.1), (P-c) and
%       (P-g) are written -- without it those five conditions name no norm;
%   (b) it supplies A2's PICARD FORM, via the singular system sigma_j asymp
%       j^{-a} of A_cond on V;
%   (c) it delivers the mild-ill-posedness and mu > 1/2 lines of the budget
%       table AS AN INTERPRETATION: s > a + 1/2 gives
%           mu = s/(2a) > (a + 1/2)/(2a) = 1/2 + 1/(4a) > 1/2,
%       with room, and mu > 1/2 is necessary for A5(iii)'s frontier to be
%       satisfiable at any r whatsoever.
%
% WHAT IT DOES NOT DO, AND ROUNDS 1-11 ASSUMED IT DID.  The clause
% "consequently the source index in prp:phirate is mu = s/(2a)" IS WITHDRAWN FROM
% THE RATE ROUTE.  Three facts force this and each is proved inside the block.
%   (1) (F-f) -- the classical Landweber/discrepancy theorem the rate rests on --
%       consumes a SPECTRAL source condition, v = (L* L)^mu varsigma with
%       ||varsigma|| <= R_*, at the operator being regularised.  H1 supplies a
%       HILBERT-SCALE condition plus a link condition.  Those are different
%       hypotheses.
%   (2) THE BRIDGE BETWEEN THEM IS OPERATOR INTERPOLATION AT EXPONENT 2 mu, AND
%       IT IS AVAILABLE ONLY AT 2 mu <= 1 -- Heinz's range, which is also the
%       interpolation range of engl2000 Sec. 8.5, THE VERY PIN H1 NAMES.
%       eq:muhalf forces 2 mu > 5/3; and lem:srctransfer(ii.a) CONSTRUCTS a 2x2
%       witness refuting operator monotonicity of z -> z^p at EVERY p > 1.  So no
%       configuration this block admits escapes, and the file already contains
%       the proof that the device is unavailable.
%   (3) FINITE-DIMENSIONALITY OF S_n DOES NOT RESCUE IT.  On S_n range membership
%       is vacuous -- every positive-definite operator there is surjective -- so
%       all the step's content is the uniform-in-n PRE-IMAGE NORM bound, and that
%       bound IS the interpolation inequality.  Finite rank makes the failure
%       quantitative rather than qualitative; it does not remove it.
%
% THE REPAIR, AND WHY IT IS A WEAKENING RATHER THAN A LOSS.  (P-a.3) now states
% the SPECTRAL condition directly, at the compression eq:compress, with mu and
% R_* PRIMITIVE.  The alternative -- keep the Hilbert-scale pair and add a
% COMMUTATIVITY proviso between A-hat* A-hat and the compressed scale operator --
% would preserve mu = s/(2a) exactly, and lem:srctransfer(ii.b) already proves the
% transfer under it.  BUT OPTION A = OPTION B + A PROVISO, so the hypothesis
% actually assumed here is STRICTLY WEAKER and needs strictly less proved.  The
% proviso is not diagnosable -- it asserts spectral alignment of a random,
% rank <= n_m operator built from 2n design points with a spline-Sobolev scale --
% and it would add a THIRD constant def:lfstep must print before the data are
% seen.  rem:pasrc(4) records it as the alternative, with its price stated.
%
% AND THE COLLAPSE WAS ALREADY PARTLY ILLUSORY, WHICH IS THE HONEST FORM OF THE
% VERDICT.  H1's matching is a statement about A_cond ON ALL OF V.  (P-a) is about
% a FINITE-RANK, SIEVE-COMPRESSED operator at a RANDOM base point, and
% lem:srctransfer(iii.a) records that a link condition CANNOT hold on V at all,
% because A-hat has finite rank (prp:hatAprop(ii)).  So H1 never delivered
% mu = s/(2a) for the operator (P-a) is about: (P-a) HAS BEEN ASSUMING ITS OWN
% MATCHING ALL ALONG.  Round 12 makes an existing assumption visible rather than
% adding one.
%
% CONSEQUENCE FOR ESCALATION ITEM 15, stated rather than absorbed.  It does NOT
% close favourably on the rate route: the parameter count there is (mu, s, a, r)
% with mu primitive at (P-a.3) and s, a entering through the scale, the budget
% table and A2 -- FOUR parameters, one of them soft.  That is exactly the cost
% the regularity-checker's third remedy names, and it is paid in the open at
% rem:pasrc(4) and rem:phirscope.  It DOES still close favourably for A2's Picard
% form and for the balancing route, neither of which invokes (F-f).
% Reference: Engl, Hanke & Neubauer (2000), Sec. 8.5.
%
% NOTE THE DOMAIN.  The singular system is that of A_cond restricted to V.  On
% all of L^2(pi) the constants lie in the kernel and the system would carry a
% zero singular value, which is not the system H1 or A2 are about.  The domain
% is fixed once, in def:spaces of results.tex.
%
% NOTE ALSO THE RANK, added in the cleanup pass.  "sigma_j asymp j^{-a} for all
% j" SILENTLY FORCES INFINITE RANK, and finite rank is not a pathology here:
% both counterexample-hunter DGPs produce a RANK-ONE A_cond satisfying every
% other hypothesis of the block.  Nothing else breaks -- A2's two-clause Picard
% form is well posed at finite rank, since the sum is then finite -- but the
% singular system spans (ker A_cond)^perp, NOT V, so with finite rank H1's
% hypothesis is simply false and H1 must not be invoked.  The infinite-rank
% requirement is therefore stated rather than smuggled in through the decay.

% ---------------------------------------------------------------------------
% PART IV.  DESIGN, OVERLAP, NORMALISATION, RATES
% ---------------------------------------------------------------------------


% DISPATCH 13.  A3 IS RETIRED, AND THE RETIREMENT IS A STRENGTHENING.  The
% history matters more than the disposition and is recorded in the statement.
%
% WHAT THE THREE EARLIER ROUNDS SAID, IN ORDER.
%   Round 1  routed identification of theta_dagger through A3.  WRONG: that made
%            identification contingent on A2.
%   Round 2  (M-3) repaired it by routing identification through A1b(iv), giving
%            as the ground that A3 is "nonsingularity of a DIFFERENT matrix (J,
%            involving ell_star) and is itself conditional on A2".  THE
%            CONDITIONALITY HALF IS RIGHT.  THE "DIFFERENT MATRIX" HALF WAS
%            ASSERTED AND NEVER COMPUTED.
%   Dispatch 10  displayed J at eq:JSig for the first time and hardened the
%            uncomputed half into an ITALICISED claim inside def:moment: "J is
%            the matrix of A3 and it is NOT E[Delta Z Delta Z']".
% So a warning about CONDITIONALITY hardened, across three rounds, into a false
% claim about IDENTITY.  THE DIAGNOSIS IS A SIGMA-FIELD SCOPE SLIP of exactly the
% species (P4b)'s warning and (P4c) exist to police: E[Wtilde . | .] kills a
% sigma(Delta U, Delta Z)-measurable factor and does NOT kill a
% sigma(X_1, X_2)-measurable one.  The nearby object that genuinely does NOT
% vanish is q = E[Wtilde partial_u ell_star | X_1, X_2] of the glossary, which
% lem:secondord and prp:sharpquad consume; a reader who has q in mind will read
% E[Wtilde ell_star Delta Z'] as non-vanishing too, and that is the error.
%
% THE FILE ALREADY CONTAINED THE REFUTATION, AT TWO SITES, WHICH IS WHY THIS IS A
% BOOKKEEPING FAILURE AND NOT A MATHEMATICAL ONE:
%   lem:nusigma(v)     E[Wtilde Xi(Delta U, Delta Z)] = 0 for every integrable
%                      Xi, under A1.  PROVED since dispatch 1.
%   rem:rootnhyp(1)    names that identity as exactly what A1's first clause
%                      buys, in Part I itself.
% Neither was ever evaluated at Xi(u,z) = ell_star(u,z) z'.
%
% FOUR FAILURE MODES CHECKED BEFORE ACTING, ALL NEGATIVE.
%  (a) SELECTION.  Does ell_star depend on W through the SELECTION of ell^Pi_star
%      from ell^Pi_dagger + ker A*_cond?  It may depend on the JOINT LAW,
%      including the law of W.  IRRELEVANT: by eq:Wdef, ell_star = W ell^Pi_star
%      is a fixed Borel map (u,z) -> R^p, and a law-dependent choice of a
%      deterministic function is still a deterministic function.  Measurability
%      with respect to sigma(Delta U, Delta Z) is a property of the map's
%      ARGUMENTS, of which W is not one.
%  (b) N1.  Does the normalisation reintroduce W?  NO, twice over: N1 constrains
%      the ANCHOR b of def:weight(b) and leaves ell untouched, and eq:JSig
%      carries ell_star and not L_star.
%  (c) TRUNCATION.  A4(i) truncates varrho-hat to [eps, 1 - eps] at def:hatT, so
%      Wtilde-hat is NOT conditionally centred.  CORRECT, AND THE FINDING DOES
%      NOT EXTEND TO IT: J is defined at the TRUE varrho (def:moment says so in
%      its own text), A4(i) makes the truncation free THERE, and the second
%      summand of eq:Jhat does not vanish at any n.  lem:sandwich Step 3 pays its
%      cost and is unchanged.  Recorded at lem:momzero(iv.2).
%  (d) INTEGRABILITY for the tower step.  |Wtilde| <= 1 by (P1),
%      ||ell_star||_inf < oo by cor:elldiff, E||Delta Z|| < oo by A1b(iv) --
%      which is the SAME bound lem:momzero Step 3 already used to conclude
%      ||J|| < oo.  Nothing new is needed.
%
% WHAT DOES NOT FOLLOW, AND THE BLOCK WOULD COLLAPSE IF IT DID.  The vanishing
% does NOT extend to the varphi-direction: Delta h is sigma(X_1, X_2)-measurable
% and NOT sigma(Delta U, Delta Z)-measurable, so E[Wtilde ell_star Delta h] is
% not zero -- indeed A2 is precisely the assertion that it equals
% -E[Delta Z Delta h] for every h in V (prp:orthog(i), eq:orthphi).  Were it zero
% for every h one would get nu_{Delta Z} = 0, i.e. the instrument does not move
% the regressor.  NEYMAN ORTHOGONALITY IS UNTOUCHED: prp:orthog is about the
% three NUISANCE directions, not about partial_theta.  Sigma is untouched and in
% fact SPLITS, eq:sigsplit, with neither summand vanishing.

\begin{assumption}[A4: Overlap and propensity regularity]\label{ass:overlap}~
\begin{enumerate}
\item[(i)] \emph{Overlap.} \textbf{Fix} $\varepsilon_\varrho\in\big(0,\tfrac1{\mathsf L}\big]$, once and
for all in this block --- \emph{that much is A4(i)'s \textbf{declaration}, and it is what
\cref{def:hatT} builds \eqref{eq:trunset} from}. \textbf{The condition} is that
$\varrho_l(\Delta Z)\ge\varepsilon_\varrho$ almost surely, \emph{for every}
$l\in\{0,1,\dots,\mathsf L-1\}$: \emph{every cell of the instrument is reached, at every value of the
controls, with probability bounded away from zero}, which rules out a design in which some cell empties
out over part of the support of $\Delta Z$. Consequently
\begin{equation}\label{eq:omegafloor}
\lambda_{\min}\big(\Omega(\Delta Z)\big)\;\ge\;\varepsilon_\Omega\;:=\;\frac{\varepsilon_\varrho}{\mathsf L}\;>\;0
\qquad\text{almost surely.}
\end{equation}
\item[(ii)] \emph{Smoothness.} Each of $\varrho_1,\dots,\varrho_{\mathsf L-1}$ belongs to a H\"older ball
of order $r>0$, the \emph{same} $r$ for every $l$, on the support of the
non-degenerate coordinates of $\Delta Z$, whose number is $d_Z$ with $1\le d_Z\le p-1$.
\item[(iii)] Estimation, and the measurability that makes the display a probability statement.
$\hat\varrho=(\hat\varrho_1,\dots,\hat\varrho_{\mathsf L-1})'$ is computed on an auxiliary fold; its
\emph{own pointwise function}
$(\omega,z)\mapsto\hat\varrho(\omega)(z)$, the one \cref{def:hatT} evaluates at each $\Delta Z_i$, is
\textbf{jointly measurable}, with respect to the product of the $\sigma$-algebra generated by the
auxiliary fold with $\mathfrak B(\IR^{p})$; and
\[
\big\|\hat\varrho-\varrho\big\|_{L^2}
\;:=\;\Big(\int\big\|\hat\varrho(z)-\varrho(z)\big\|^{2}\,P_{\Delta Z}(dz)\Big)^{1/2}
=O_P\big(n^{-r/(2r+d_Z)}\big)=:O_P(r_n),
\]
the inner $\|\cdot\|$ being Euclidean on $\IR^{\mathsf L-1}$ and the integral running over $z$ alone, so
that the left side is a function of the auxiliary fold.
\end{enumerate}
\end{assumption}

% ===========================================================================
%  DISPATCH-25 RESTATEMENT OF A4.  ALL THREE CLAUSES MOVE; THE TAG, THE CLAUSE
%  NUMBERS AND THE CONTENT AT sans-L = 2 DO NOT.  Phase 1 of
%  quality_reports/plans/2026-08-08_multinomial-instrument.md.
%
%  WHAT (i) SAID:   eps_varrho in (0, 1/2) with
%                   varrho(Delta Z) in [eps_varrho, 1 - eps_varrho] a.s.
%  WHAT (i) SAYS:   eps_varrho in (0, 1/sans-L] with
%                   varrho_l(Delta Z) >= eps_varrho a.s. for EVERY
%                   l in {0, 1, ..., sans-L - 1}, and eq:omegafloor.
%
%  THE UPPER END WENT FROM STRICT TO WEAK 2026-08-27, AND IT IS A WEAKENING WITH
%  NO CONSUMER.  eq:trunset = {q: q_l >= eps for l = 1..L-1, sum q_l <= 1 - eps}
%  is non-empty IFF L eps <= 1: at q_l = eps the constraint reads
%  (L-1) eps <= 1 - eps.  So the strict inequality the block printed was strictly
%  inside the referent boundary, at the single spelling eps = 1/L, where every
%  symbol in the corpus keeps a referent and the strict inequality is FALSE --
%  which is exactly the printed split-first test's definition of a CLAUSE.  A
%  clause printed inside a DECLARATION is the register defect the block condemns
%  for B1(vi), read on a constant instead of on a tag.
%   WHY STRUCK AND NOT MOVED TO A HYPOTHESIS LINE.  It excludes no law: if P
%  satisfies overlap at some eps it satisfies it at eps/2, so the admitted class
%  of P is IDENTICAL under either upper end.  And no step consumes it -- the two
%  consumers of the range are eq:trunset's non-emptiness (weak suffices) and
%  eq:omegafloor (needs eps > 0 only), while rem:hatTread(2)'s two-sided per-cell
%  bound eps <= |Wtilde-hat_{i,l}| <= 1 holds at eps = 1/L as well.  THAT LAST IS
%  A READING AND NOT A CENSUS.
%   WHAT THE WEAK END ADMITS THAT THE STRICT ONE DID NOT: the single choice
%  eps = 1/L, at which the CONDITION forces varrho == (1/L, ..., 1/L) a.s. and the
%  instrument is degenerate.  Nothing is wrongly admitted -- that law fails A1b(iii)
%  and A2 -- and the degeneracy is caught where degeneracies belong, at
%  identification, not inside a tuning constant's range.
%
%  2026-08-26 RE-SPELLING, AND IT CHANGES NO LAW AND NO RESULT.  (i) read "THERE
%  IS eps_varrho in (0, 1/sans-L) with ..." and now reads "FIX eps_varrho in
%  (0, 1/sans-L)" + "the condition is varrho_l >= eps_varrho a.s. for every l".
%  THE ADMITTED CLASS OF P IS THE SAME -- a block that fixes the constant and
%  imposes the inequality at it admits exactly the laws for which some such
%  constant works -- and every result carrying A4(i) reads the inequality at the
%  declared constant, as before.  WHY IT WAS DONE: the existential form has no
%  declarative half to cite, so def:hatT's CONSTITUTIVE use of eps_varrho had no
%  register to land in, and the block's own (superseded) negation test then ruled
%  OVERLAP a declaration carrying no hypothesis -- see the header's
%  "SPLIT FIRST, THEN TEST THE REMAINDER" entry.  eq:trunset needs eps_varrho and
%  eps_varrho <= 1/sans-L, both in the declaration [WEAK FROM 2026-08-27; the
%  2026-08-26 spelling of this line said "< 1/sans-L", which was the defect the
%  entry above repairs]; the PROJECTION's property that
%  it cannot increase the L^2 error needs varrho itself to lie in eq:trunset,
%  which is the CONDITION, and that is charged where it is used and not here.
%
%  THE COLLAPSE AT sans-L = 2 IS EXACT AND IT IS THE FIRST THING TO CHECK.
%   - The parameter range.  (0, 1/sans-L] at sans-L = 2 IS (0, 1/2], which is the
%     interval earlier rounds printed WITH ITS ENDPOINT RESTORED -- the same SET
%     up to that one point; the glyphs differ, as they do everywhere the
%     cardinality is carried symbolically.  At the endpoint eq:trunset is
%     [1/2, 1/2] and the condition forces varrho_1 == 1/2, the degenerate case
%     the entry above disposes of.
%   - The quantifier.  "for every l in {0, 1}" is the CONJUNCTION of
%     varrho_1 >= eps and varrho_0 >= eps.  With varrho_0 = 1 - varrho_1 the
%     second is varrho_1 <= 1 - eps.  So the two-sided interval
%     [eps, 1 - eps] IS the sans-L-fold conjunction read at sans-L = 2, and
%     nothing was added or dropped.  THE BASE CELL l = 0 IS IN THE QUANTIFIER
%     FOR EXACTLY THIS REASON: dropping it would silently weaken the binary
%     condition to one-sided.
%   - eq:omegafloor.  At sans-L = 2, Omega = varrho(1 - varrho), and the clause
%     asserts varrho(1-varrho) >= eps/2, which is TRUE and WEAKER than the sharp
%     eps(1-eps) the block used before.  The floor is used only as a positive
%     lower bound, never at its value, so nothing that consumed eps(1-eps)
%     degrades; see the sharpness note below.
%
%  WHY eq:omegafloor IS PRINTED IN THE CLAUSE THOUGH IT IS NOT A RESTRICTION.
%  "Consequently" is a derivation, and the derivation is in assumptions_supp.tex
%  under the A4 heading -- the A8 precedent, where the same word carries the
%  same obligation and the same treatment.  IT IS NOT A SECOND HYPOTHESIS: it
%  follows from the first sentence with no further condition, so no hypothesis
%  list is strengthened by its presence and every list citing A4(i) may cite
%  eq:omegafloor without acquiring anything.  IT IS PRINTED BECAUSE THE
%  NORMALISATION DIVIDES BY Omega and a result that divides must be able to
%  point at the bound it divides under, by tag, in one place.
%    THE PLAN ASKED FOR IT AS A CONJUNCT ("... AND lambda_min(Omega) >= eps_Omega").
%    IT IS STATED AS A CONSEQUENCE INSTEAD, AND THE DEPARTURE IS DELIBERATE:
%    an assumption that states its own corollary as a second hypothesis is
%    exactly what the assumption-auditor deducts for, and this block has already
%    retired two clauses (A3, A2'(ii)) on the ground that they were theorems.
%    Stating a third one would re-create the defect the retirements removed.
%
%  THE CONSTANT eps_Omega := eps_varrho / sans-L IS VALID AND NOT SHARP.  The
%  derivation runs through Omega^{-1} = diag(varrho)^{-1} + varrho_0^{-1} 1 1',
%  whose norm is at most max_l varrho_l^{-1} + (sans-L - 1) varrho_0^{-1}
%  <= sans-L / eps_varrho.  Sharper constants exist -- at sans-L = 2 the exact
%  minimum of varrho(1-varrho) on [eps, 1-eps] is eps(1-eps) -- and none is
%  needed: no bound in this corpus is stated AT the value of the floor, only
%  under it.  A dispatch that starts to track the constant should re-derive it
%  rather than sharpen this one in place.
%
%  (ii) IS PER COMPONENT AND CARRIES THE SAME r AT EVERY l, AND THAT IS A
%  CHOICE.  A Holder ball is a condition on a SCALAR function, so the vector
%  statement is sans-L - 1 scalar statements and there is no third option.  What
%  is chosen is the COMMON ORDER: allowing r_l to differ by cell would make
%  alpha = r/(2r + d_Z), hence h_u, delta-hat_n and the whole rate architecture,
%  cell-dependent, and every downstream display that carries one alpha would
%  have to carry sans-L - 1 of them.  The common r is the minimum of the
%  cellwise orders and costs nothing: a function in a Holder ball of order r_l
%  is in one of order min_l r_l on a bounded support.  At sans-L = 2 there is
%  one component and the clause is the printed one.
%
%  (iii) IS STATED IN THE VECTOR NORM, AND THE ALTERNATIVE WAS COMPONENTWISE.
%  The two are the SAME O_P class at fixed finite sans-L:
%      max_l ||hat-varrho_l - varrho_l||_{L^2}
%        <= ||hat-varrho - varrho||_{L^2}
%        <= sqrt(sans-L - 1) max_l ||hat-varrho_l - varrho_l||_{L^2},
%  so nothing is gained or lost in RATE either way and the choice is about which
%  symbol propagates.  THE VECTOR FORM IS TAKEN, for three reasons.
%   (1) EVERY CONSUMER READS ONE SCALAR.  A5(i)'s branch delta-hat_n^{(varrho)}
%       := r_n, lem:opnoise's Step 4, lem:secondord(iv)'s sans-r direction and
%       C5'(iii) each want a single number; a componentwise statement would
%       force max_l back in at each of them, in four places instead of one.
%   (2) IT IS THE BLOCK'S OWN IDIOM.  A2'(ii) "is written with the Euclidean
%       norm, which is the coordinatewise statement summed" -- dispatch-2 change
%       log, item 4.  Same device, same reason.
%   (3) AT sans-L = 2 IT IS THE PRINTED DISPLAY, since the Euclidean norm of a
%       one-vector is its absolute value.  The ":=" spelling out the norm is new
%       and is there because "||.||_{L^2}" of a VECTOR-valued function has two
%       readings and the block may not leave that to the reader.
%  THE FACTOR sqrt(sans-L - 1) IS AN sans-L-DEPENDENT CONSTANT.  It is not a
%  rate and it does not touch alpha, but per the plan's "Rates" item no constant
%  in this corpus may be claimed sans-L-free without checking, and this is one
%  that is not.
%
%  WHAT IS *NOT* IN A4, AND MUST NOT BE PUT THERE.  A4 says nothing about how
%  hat-varrho is computed.  A multinomial logit, a series multinomial fit and
%  sans-L - 1 separate binary regressions all satisfy (iii) if they meet the
%  rate; the block prices the rate and the measurability and nothing else, which
%  is the same standing (iii) had when the fit was a single binary regression.
%  THE ONE THING THE FIT MUST DELIVER IS A POINT IN THE SIMPLEX -- non-negative
%  and summing to at most one -- and that is NOT assumed here either: def:hatT
%  imposes it by TRUNCATION, for the reason recorded at (e) below and at
%  rem:hatTread(2).
% ===========================================================================

% WHAT A4(i) DOES, AND -- ROUND-2 CORRECTION -- WHAT IT DOES NOT.
%
%  (a) IT MAKES THE sans-L CONDITIONAL LAWS WELL DEFINED.  Every W-cell is
%      non-null given Delta Z = z for a.e. z, so F_0(.|z), ..., F_{sans-L-1}(.|z)
%      are pinned down a.e. rather than arbitrary on a null set.  Used in
%      lem:quantu.  (THE COUNT WAS "BOTH" THROUGH DISPATCH 24, WHICH WAS THE
%      BINARY READING OF THE SAME SENTENCE.)
%
%  (b) IT MAKES THE PREFACTOR INVERTIBLE -- AND AT sans-L >= 3 THE PREFACTOR IS A
%      MATRIX.  (THE HEADING SAID "CANCELS" AND NAMED THIS AS "WHERE THE
%      NORMALISATION DECISION LIVES"; the estimator does not cancel it, and the
%      decision was taken the other way.)  At sigma(Delta Z), and
%      only there, writing F_m(u|z) := P(Delta U <= u | W = m, Delta Z = z) and
%      d(u|z) := (F_1 - F_0, ..., F_{sans-L-1} - F_0)'(u|z),
%          E[Wtilde 1{Delta U <= u} | Delta Z = z]  =  Omega(z) d(u|z),
%      which is eq:omegaprefac of the Discussion, where it is derived.  AT
%      sans-L = 2 THIS IS THE OLD DISPLAY VERBATIM:
%          varrho(z)(1 - varrho(z)) { F_1(u|z) - F_0(u|z) }.
%      Under A4(i), eq:omegafloor bounds Omega below and the conditional
%      restriction can be divided through -- now by a MATRIX INVERSE and not by
%      a scalar.  THAT DIVISION IS WHAT rem:htread's reading uses.  IT IS NOT
%      APPLIED IN THE ESTIMATOR: the operator and psi carry the RAW contrast, and
%      the division appears only in the interpretation of what the restriction
%      says, never as a weight.
%
%      *** BUT A4(i) IS NOT WHAT MAKES lem:bierens AN EQUIVALENCE, AND THE FIRST
%      DISPATCH-25 FORM OF THIS NOTE SAID IT WAS.  THE STRUCK TEXT READ: "This is
%      what makes lem:bierens an EQUIVALENCE rather than an implication: without
%      A4(i) the restriction is vacuous wherever some cell is null, the degenerate
%      region being {z : det Omega(z) = 0} = {z : min_l varrho_l(z) = 0}."  BOTH
%      HALVES ARE FALSE AT sans-L >= 3 AND THE SECOND IS THE DANGEROUS ONE. ***
%
%      THE COUNTEREXAMPLE, AT sans-L = 3.  Take varrho = (q, 0)' with q in (0,1),
%      so cell 2 is null and det Omega = varrho_0 varrho_1 varrho_2 = 0.  W is NOT
%      conditionally degenerate -- it takes cells 0 and 1 with positive
%      probability -- Wtilde is not zero, and the restriction still forces
%      F_1 = F_0.  So {det Omega = 0} is the set where Omega is SINGULAR, and the
%      restriction is NOT vacuous there; it merely says less, restricting only the
%      non-null cells.  At sans-L = 2 the two sets coincide, which is why the
%      binary sentence was true and its generalisation through the determinant is
%      not.
%
%      THE CORRECT STATEMENT, AND THE EQUIVALENCE NEEDS NO OVERLAP AT ALL.  Fix z
%      and put A := {l >= 1 : varrho_l(z) > 0}.  Omega vanishes off A x A, and its
%      A-block is Omega_A = diag(varrho_A) - varrho_A varrho_A' with
%      det Omega_A = varrho_0 prod_{l in A} varrho_l.  Two cases.
%        varrho_0 > 0:  Omega_A is nonsingular, so Omega d = 0 gives d_A = 0, i.e.
%           F_l(.|z) = F_0(.|z) at every non-null l.  Null cells impose nothing and
%           their F_l is an arbitrary version.
%        varrho_0 = 0:  ker Omega_A = span{1_A} -- from varrho_l(c_l - varrho_A'c)
%           = 0 with varrho_l > 0 on A, c is constant on A -- so d is constant on
%           A, which again says all non-null cells share one conditional law.
%      IN BOTH CASES the restriction holds iff the conditional law of Delta U given
%      (W, Delta Z = z) is the same across all NON-NULL cells, which IS conditional
%      independence at z.  NO CONDITION ON varrho IS USED.  Read at test functions:
%          E[Wtilde_l ς | z] = varrho_l ( E[ς | W = l, z] - sum_m varrho_m E[ς | W = m, z] ),
%      whose l-th component is identically zero when varrho_l = 0 and otherwise
%      says E[ς | W = l, z] equals the common mixture; averaging that over l with
%      weights varrho_l returns varrho_0 (mixture - E[ς | W = 0, z]) = 0, which
%      closes the varrho_0 > 0 case, and the varrho_0 = 0 case is the display
%      itself.
%
%      SO THE ONLY SET ON WHICH THE RESTRICTION IS VACUOUS is where W is
%      CONDITIONALLY DEGENERATE, i.e. where Omega(z) = 0, i.e.
%          Z_0 = {z : max_l varrho_l(z) = 1} = {z : Omega(z) = 0},
%      which at sans-L = 2 is {z : varrho(z)(1 - varrho(z)) = 0} -- the set Part C
%      already carries under the name Z_0, recovered verbatim.  {z : det Omega = 0}
%      is a STRICTLY LARGER set at sans-L >= 3 and is NOT the degenerate region.
%
%      WHY THIS MATTERS BEYOND THIS COMMENT.  lem:bierens states the equivalence
%      under B1(i)-(iv) ALONE, and results_s2_identification.tex:282-290 records
%      A4(i)'s removal from it as a deliberate cleanup-pass correction -- but the
%      GROUND it gives there is the binary one ("on D the prefactor
%      varrho(1-varrho) is zero, so both sides vanish whatever version is used").
%      THAT GROUND DOES NOT GENERALISE; the CONCLUSION does, by the argument above.
%      PHASE 2 MUST REPLACE THE GROUND AND KEEP THE CONCLUSION, and must define D
%      as {Omega = 0} and not as {det Omega = 0}.  A dispatch that generalises the
%      printed reason instead of the printed result will import an overlap
%      hypothesis lem:bierens does not have and does not need.
%
%      THE HAND-OFF IS UNQUALIFIED, AND ROUND 4 REMOVES A QUALIFICATION ROUNDS
%      2-3 ATTACHED TO IT.  Under the normalisation they had to say that
%      "lem:bierens needs no overlap" was true of (a) <=> (b) and false of the
%      clauses naming M and T, because those objects did not exist without
%      Omega^{-1}.  WITH THE RAW CONTRAST RESTORED, M AND T ARE DEFINED
%      EVERYWHERE AND BOUNDED, so NO CLAUSE OF lem:bierens NEEDS A4(i) AND THE
%      HAND-OFF IS UNQUALIFIED AT EVERY CLAUSE.  A reader who finds the round-3
%      clause-by-clause carve-out is reading a struck instruction.
%      (THE ROUND-3 PHRASING WAS ALSO WRONG ON ITS OWN TERMS AND THE CORRECTION
%      IS WORTH KEEPING, because it recurs whenever an IPW object is discussed:
%      Wtilde^Omega is NOT "non-integrable" where a cell vanishes -- it is
%      UNDEFINED there, Omega having no inverse.  Where it IS defined it is
%      perfectly integrable, and with a varrho-FREE bound: rem:htread computes
%      E[|Wtilde^Omega_l| | Delta Z] = 2 exactly.  Non-integrability and
%      undefinedness are different failures and the file must not conflate them.)
%
%  (b') THE COMPONENTWISE ALTERNATIVE IS WRONG, AND THAT SURVIVES THE REVERSAL.
%      Whenever the prefactor of (b) is inverted -- which the ESTIMATOR does not
%      do, but which rem:htread's reading does -- the inverse is Omega^{-1} and
%      NOT diag(varrho_l(1 - varrho_l))^{-1}.  The map in (b) is Omega, whose
%      off-diagonal entries are -varrho_l varrho_m: the NEGATIVE cross-cell
%      correlation that a multinomial has and a collection of binaries does not.
%      Dividing component l by varrho_l(1 - varrho_l) therefore does NOT return
%      d_l, and what it returns is not the contrast between cell l and the base
%      cell at all.  Two identities make Omega^{-1} the right inverse to take,
%      and both are derived in the Discussion at eq:omegainv: Omega^{-1} Wtilde
%      is the vector of INVERSE-PROBABILITY contrasts
%      1{W = l}/varrho_l - 1{W = 0}/varrho_0; and the quadratic form
%      Wtilde' Omega^{-1} Wtilde = (1 - varrho_W)/varrho_W is INVARIANT TO WHICH
%      CELL IS THE BASE, which the componentwise weighting is not.
%      AT sans-L = 2, Omega^{-1} Wtilde = (W - varrho)/(varrho(1 - varrho))
%          = 1{W = 1}/varrho - 1{W = 0}/(1 - varrho),
%      which is the identity a binary design writes down without comment.  IT
%      DID NOT APPEAR ANYWHERE IN THIS FILE before dispatch 25, in that or any
%      other spelling -- swept then; bar-w in this file is B1(ii)'s Bierens
%      envelope and never a propensity.
%
%      *** AND IT IS NOT PLACED IN THE ESTIMATOR.  Omega^{-1} APPEARS ONLY IN
%      rem:htread. ***  The three candidate homes and their final dispositions:
%        THE theta MOMENT psi -- NOT NORMALISED.  eq:psidef carries the RAW
%          contrast.  (Round 2 said the normalisation there was REDUNDANT via a
%          bijection of def:weight's class; that argument is false -- clause (a)
%          is a RANGE condition -- and round 3 accordingly normalised psi too.
%          Round 4 normalises NEITHER side, so the question the bijection was
%          invented to settle no longer arises.)
%        THE varphi OPERATOR T AND ITS SAMPLE VERSION T-hat -- NOT NORMALISED.
%          eq:hatT carries the raw contrast, and def:operator keeps its own.
%          THE AUTHOR'S GROUND IS THAT NORMALISING MOVES A_cond, and A2, A2',
%          A1c and H1 are all stated at A_cond, with mu = s/(2a), eq:frontier,
%          the feasibility table and the CFF comparison downstream of a.
%        THE RANGE MEASURE Pi_w -- A DIFFERENT OBJECT, and it is the one place a
%          contrast weight does legitimately live.  Pi_w weights the contrasts in
%          the L^2(Pi) NORM and is chosen by the analyst; B1(iv.d) prices it.
%          It is NOT Omega^{-1} and must not be set to "the efficient weights".
%
%  (c) IT DOES *NOT* MAKE A_cond A REWEIGHTING OF THE RAW W-CONTRAST OPERATOR.
%      Round 1 asserted -- following plan Sec.18.3 -- that A_cond and the raw
%      W-contrast operator "differ by the bounded, boundedly-invertible weight
%      varrho(1-varrho), hence have the same range and the same kernel".
%      THAT IS FALSE AND IS WITHDRAWN.  The identity in (b) holds only at
%      sigma(Delta Z).  The operator's kernel requires conditioning on
%      (X_1, X_2, Delta Z), where, writing
%      p_{X,m} := P(W = m | X_1, X_2, Delta Z) and
%      f_m(u) := f_{Delta Y | X_1,X_2,Delta Z,W = m}(u), the l-th component is
%          E[Wtilde_l f_{Delta Y | X_1,X_2,Delta Z,W}(u) | X_1,X_2,Delta Z]
%             = p_{X,l} f_l(u) - varrho_l sum_m p_{X,m} f_m(u),
%      which equals the l-th component of Omega (f_1 - f_0, ..., f_{sans-L-1} -
%      f_0)' ONLY IF p_{X,m} = varrho_m for every m, i.e. only if W is
%      independent of (X_1,X_2) given Delta Z.  (AT sans-L = 2 THE DISPLAY IS THE
%      OLD ONE: with p_X := p_{X,1} it reads p_X(1 - varrho) f_1(u)
%      - (1 - p_X) varrho f_0(u), and the condition is p_X = varrho.)  That says
%      THE INSTRUMENT DOES NOT SHIFT THE ENDOGENOUS REGRESSOR GIVEN Delta Z, and
%      it MUST NOT BE ASSUMED.  MORE CELLS DO NOT WEAKEN THIS: the requirement is
%      an equality of sans-L - 1 functions instead of one, so the false route
%      round 1 took is harder to take, not easier.  See rem:draftrec(ii) in results.tex
%      for the correct comparison, which is between the two operators' KERNELS
%      and claims neither equality of range nor equality of kernel.
%      CONSEQUENCE: plan Sec.18.3's bullet "A4's mechanism" is refuted as
%      stated, and the reconciliation it asked for is delivered in the weaker,
%      correct form.  Escalated in theory_memo.md Sec.7, item N-7.
%
%  (d) IRREGULARITY -- AND THE REFERENT IS L_star, NOT ell_star.  ROUND-7
%      CORRECTION.  As overlap degrades, A2 can continue to hold AS A RANGE
%      STATEMENT while the moment's variance diverges: that is irregular
%      identification in the sense of Khan & Tamer (2010).  THE VARIANCE OBJECT
%      IS E||L_star Wtilde||^2, because psi = Delta Z Delta U
%      + L_star(Delta U, Delta Z) Wtilde (eq:psidef) carries L_star and not
%      its derivative.
%      TWO TYPE SLIPS IN THE ROUND-2 FORM OF THIS PARAGRAPH ARE CORRECTED HERE
%      AND THEY RECUR BELOW, so the rule is stated once: (1) THE CONTRAST STANDS
%      TO THE RIGHT of the p x (sans-L - 1) weight -- "Wtilde L_star" is
%      dimensionally ILL FORMED AT sans-L >= 3 and was inherited from the scalar
%      era.  (AT sans-L = 2 IT IS WELL FORMED AND IS WHAT eq:psidef PRINTS: both
%      orders agree when the contrast is a scalar.  Round 3 wrote "sans-L >= 2"
%      here, which condemned the printed display.)  (2) "Wtilde^2" NAMES NOTHING
%      for a vector.  Where the old text wrote
%      E[Wtilde^2 ||A||^2] read E||A Wtilde||^2, and where it wrote a scalar
%      square read the appropriate quadratic form.  Every occurrence below is to
%      be read under this rule; they are not rewritten one by one because the
%      paragraph is an archaeological record of four earlier corrections and
%      re-typing it would destroy what it records.
%      Rounds 1-6 named
%      E[Wtilde^2 ell_star^2] here -- THE DERIVATIVE -- which is the object
%      lem:balance(i) and lem:secondord(iv) consume through eps_n, and is a
%      different quantity.  This is the same referent slip that A5(iv)'s round-5
%      correction fixed in the opposite direction (there the block priced the
%      LEVEL where the consumers needed the DERIVATIVE; here it priced the
%      DERIVATIVE where the variance needs the LEVEL).
%
%      WHAT SURVIVES -- REWRITTEN IN DISPATCH 14, BECAUSE THE ROUND-7 VERSION OF
%      THIS PARAGRAPH CONTAINED TWO FALSE STATEMENTS AND ONE LEGITIMATE GROUND.
%      The two false ones are struck; the legitimate one is kept and is now the
%      only ground stated.
%
%        FALSE 1 (struck).  "A2'(ii) stays as written, because it is the clause
%        the second-order and bias terms consume."  A2'(ii) IS CONSUMED BY NO
%        RESULT IN THIS BLOCK.  results.tex says so three times -- at
%        lem:balance's hypothesis line, at lem:solset and at lem:momzero's closing
%        paragraph -- and the ROUND-8 CHANGE LOG OF THIS VERY FILE says so:
%        "A2'(ii) is consumed by none of them and is a Step-8 clause."  A comment
%        block asserting the opposite of the change log twelve hundred lines below
%        it is the defect, not the citation.
%
%        FALSE 2 (struck).  The bare "(i) + (iii)" in the sentence beginning "What
%        the irregularity discussion now points at" was ambiguous between A2'(i),
%        A2'(iii) and -- inside a comment attached to A4 -- A4(i), A4(iii); read
%        the second way it attributes the finiteness of E[Wtilde^2 ||L_star||^2]
%        to OVERLAP, and THAT IS WRONG.  A4(i) PLAYS NO ROLE IN THE VARIANCE
%        BOUND.  The bound is
%            E[Wtilde^2 ||L_star||^2] <= 2 E||b_star||^2
%                                        + 2 ||ell_star||_inf^2 E[Delta U^2],
%        whose FOUR ingredients are: |Wtilde| <= 1, which is (P1) and is FREE
%        from W finitely supported and varrho_l in [0,1] -- the NOTE below this
%        comment has said so since round 2, and dispatch 25 replaces its
%        "W in {0,1}" ground by the weaker and correct one; at sans-L >= 3 the
%        Euclidean bound is sqrt 2 rather than 1, so the display above carries a
%        further factor 2, which is a CONSTANT and not a rate;
%        ||ell_star||_inf < infinity, which is cor:elldiff
%        under A2'(i) IN FULL -- (i.a) AND (i.b), since cor:elldiff is stated at
%        A2b and A2b requires both, and since A2'(i)'s closing clause is what
%        makes ell_star pointwise well defined at all -- with B1(i), (ii),
%        (iv.a); E||b_star||^2 < infinity, which is A2'(iii); and
%        E[Delta U^2] < infinity, which is A1b(iv).  The tags are therefore
%        A2'(i), A2'(iii), A1b(iv), B1(i), (ii), (iv.a) -- and NOT A4(i).  The
%        counterexample-hunter reached this from the other side and put it in one
%        line: "A4(i) plays no role in the variance bound."
%
%        WHAT SURVIVES, AND IT IS ENOUGH TO RETAIN THE POINT.  Overlap does bear
%        on the SIZE of the variance and not on its finiteness, and the mechanism
%        is exact rather than qualitative: with the two conditional laws given W
%        held fixed and varrho == eps, A_cond^{(eps)} = eps(1-eps) B with B free of
%        eps, so RANGES ARE INVARIANT and A2 holds at every eps > 0 or at none --
%        the range statement overlap does not touch -- while
%        ell_star^{(eps)} = {eps(1-eps)}^{-1} ell^{(eps)} and, under A1,
%        E[Wtilde^2 ell_star^2] = E[(ell^{(eps)})^2]/{eps(1-eps)} asymp
%        eps_varrho^{-1} -> infinity.  That is Khan-Tamer irregularity, it lives
%        where A4(i) FAILS, and it is why A4(i) is retained.  NOTE THE OBJECT: the
%        divergence above is displayed for E[Wtilde^2 ||ell_star||^2], the
%        DERIVATIVE, which is A2'(ii)'s object; the same computation runs for
%        E[Wtilde^2 ||L_star||^2], the LEVEL, through ||ell_star||_inf.
%
%        DISPATCH-25 FORM, AND THE MECHANISM IS THE SAME ONE.  The scalar
%        prefactor eps(1-eps) becomes the MATRIX Omega, so the display above is
%        A_cond^{(Omega)} = Omega B with B free of Omega.
%        *** THE RANGE-INVARIANCE SENTENCE IS SCOPED, ROUND 3.  It read "ranges
%        are invariant under multiplication by a boundedly invertible matrix
%        exactly as under multiplication by a nonzero scalar, so A2's range
%        statement is untouched at every sans-L".  THAT IS TRUE AS WRITTEN ONLY
%        WHEN Omega IS CONSTANT IN z -- which is the case this paragraph
%        constructs, varrho == eps.  For a z-VARYING Omega the multiplication is
%        by a non-constant matrix function, and the same fact that refutes the
%        weight-class bijection at the placement block applies here: right- or
%        left-multiplication by a merely Hoelder Omega(z) need not preserve the
%        range of an operator whose sections are real-analytic.  THE PARAGRAPH'S
%        OWN CONCLUSION SURVIVES, because it is computed at varrho == eps; what
%        is withdrawn is the unqualified "at every sans-L" generality. ***
%        Meanwhile
%        ell_star = ell Omega^{-1} -- THE MATRIX ACTS ON THE RIGHT, ell being
%        p x (sans-L - 1); round 2 wrote Omega^{-1} ell, which is ill-formed at
%        sans-L >= 3 -- and the variance carries ||Omega^{-1}||.  NOTE THAT THIS
%        Omega^{-1} IS NOT THE STRUCK NORMALISATION: it arises from SOLVING the
%        adjoint equation at a prefactor Omega, not from a weight the estimator
%        applies, and it is present in the unnormalised formulation exactly as it
%        was at sans-L = 2.  THE
%        DIVERGENCE IS EXACT AND NEEDS NO CHOSEN PATH: from
%        Omega^{-1} = diag(varrho)^{-1} + varrho_0^{-1} 1 1' one reads
%            ||Omega^{-1}|| >= e_l' Omega^{-1} e_l = varrho_l^{-1} + varrho_0^{-1}
%                           >= max_{m in {0, l}} varrho_m^{-1},
%        so ||Omega^{-1}|| -> infinity as ANY ONE cell probability tends to zero,
%        base cell included.  At sans-L = 2 this is 1/(varrho(1 - varrho)) asymp
%        eps^{-1}, the printed rate.  KHAN-TAMER IRREGULARITY THEREFORE LIVES
%        EXACTLY WHERE A4(i) FAILS AT ANY ONE l, which is why (i) quantifies over
%        every l and not only over the non-base cells.  THE APPLIED READING
%        MATTERS AND IS RECORDED AT THE PLAN'S Sec.5: with sans-L = 4 and one
%        thin cell, it is that cell alone that drives the variance, and the
%        realised min_l hat-varrho_l is the diagnostic to report.
%
%      A2'(ii) IS THEREFORE RETIRED IN PLACE (see the clause itself), and the only
%      ground the block ever had for keeping it -- a FORWARD Step-8 need at
%      prp:effbound, where the variance is computed and where the T3 tier relaxes
%      A2'(i) -- is a forward obligation and not a present consumer.  It is
%      recorded at the retirement note and flagged for Step 8.
%
%  NOTE, AND IT STANDS UNQUALIFIED.  (THE HEADING READ "NOW HALF STALE -- READ
%  THE SECOND PARAGRAPH BEFORE ACTING ON THE FIRST" THROUGH ROUND 4, WHILE THE
%  PARAGRAPH THREE LINES BELOW SAID ROUND 4 HAD REMOVED THAT VERY QUALIFICATION.
%  The heading was the round-3 normalisation's and outlived it.)
%  A4(i) is NOT needed for ||Wtilde||_inf <= 1, which is free from W
%  taking finitely many values and varrho_l in [0,1].  Round 1 listed A4(i) as a
%  hypothesis of prp:wellposed, prp:frechet and prp:compact, where it is idle; it
%  has been removed from all three.  (THE GROUND WAS WRITTEN "W in {0,1} and
%  varrho in [0,1]" THROUGH DISPATCH 24.  It was TRUE and it was NARROWER THAN
%  THE FACT: nothing in the argument uses two cells, only that each indicator and
%  each cell probability lies in [0,1].  Dispatch 25 states the fact at the
%  strength it holds, which is also the strength the sans-L-category design
%  needs.)
%
%  THE PARAGRAPH ABOVE STANDS UNQUALIFIED, AND ROUND 4 REMOVES A QUALIFICATION
%  ROUNDS 2-3 PUT ON IT.  They flagged it "half stale", on the ground that the
%  operator had been re-written at the normalised contrast, whose sup-norm bound
%  is eps_varrho^{-1} rather than 1.  THE NORMALISATION IS REVERSED, so:
%  ||Wtilde||_inf <= 1 IS FREE AGAIN, A4(i) IS IDLE IN prp:wellposed,
%  prp:frechet AND prp:compact, AND M-10's REMOVAL STANDS.  NOTHING IS TO BE
%  REVERSED THERE.  A reader who finds the round-3 instruction to reinstate A4(i)
%  in those three results is reading a struck instruction.
%
%  WHAT IS KEPT FROM ROUNDS 2-3, BECAUSE IT IS ABOUT THE CONTRAST AND NOT ABOUT
%  THE ESTIMATOR: the arithmetic of the two bounds, which is now recorded at
%  rem:htread and which a later reader may need if the normalisation is ever
%  reconsidered.  In brief, and both verified:
%     RAW CONTRAST.  ||Wtilde||_inf <= 1 and ||Wtilde||_2 <= min{sqrt 2,
%       sqrt(sans-L - 1)} -- at W = k >= 1, (1 - varrho_k)^2 + sum_{l != k}
%       varrho_l^2 <= 2; at W = 0, sum_l varrho_l^2 <= 1.  UNIFORM IN sans-L and
%       FREE of varrho.  THIS IS THE BOUND THE ESTIMATOR USES.
%     NORMALISED CONTRAST (rem:htread only).  ||Wtilde^Omega||_inf <=
%       eps_varrho^{-1}, and ||Wtilde^Omega||_2 <= sqrt(sans-L - 1)
%       eps_varrho^{-1}, SHARP at W = 0.  ROUND 2 WROTE "at most two nonzero
%       entries" AND DREW min{sqrt 2, sqrt(sans-L - 1)} eps_varrho^{-1} FROM IT;
%       BOTH ARE FALSE.  The pattern is "one entry, or all of them": at W = k >= 1
%       exactly the k-th entry is nonzero, at W = 0 ALL sans-L - 1 entries equal
%       -1/varrho_0.  The round-2 constant is right at sans-L = 2 and 3 and
%       understates at sans-L >= 4, i.e. at Section 6's design, by sqrt(3)/sqrt(2).
%  THE ERROR WAS TRANSPORTING THE RAW BOUND: the raw contrast is DENSE AND SMALL
%  and its Euclidean bound does not grow with sans-L; the normalised one is
%  SPARSE-OR-CONSTANT AND LARGE and grows like sqrt(sans-L - 1).  Recorded here
%  because the correction was made under the normalisation and must not be lost
%  with it.
%
%  (e) DISPATCH 7: A4(i) NOW ALSO LICENSES A CONSTRUCTION, AND THE CONSTRUCTION
%      IS WHAT KEEPS THE SAMPLE OPERATOR NON-DEGENERATE.  def:hatT truncates
%      varrho-hat to the eps_varrho-interior simplex eq:trunset -- which at
%      sans-L = 2 IS the interval [eps_varrho, 1 - eps_varrho], since the two
%      inequalities q_1 >= eps and q_1 <= 1 - eps are its two faces.  A4(i) is
%      what makes the truncation FREE -- varrho already lies in that set, so the
%      metric projection onto it is a pointwise contraction and A4(iii) is
%      preserved -- and what makes it USEFUL: |Wtilde-hat_{i,l}| >= eps_varrho
%      SURELY, for every i AND EVERY l.  THAT COMPONENTWISE BOUND IS THE
%      DEGENERACY GUARD AND IT IS THE ONE THE DEGENERACY BULLETS BELOW ARE READ
%      AT, because eq:hatT carries the RAW contrast.  (Rounds 2-3 replaced it
%      here by a guard on the normalised contrast, ||hatWtilde^Omega_i||_inf >= 1,
%      and by a surely-bound lambda_min(Omega-hat_i) >= eps_Omega supporting an
%      inverse eq:hatT no longer takes.  Both go with the normalisation.  The
%      eigenvalue bound remains TRUE on eq:trunset and is recorded at item 3 of
%      the def:hatT comment as an unconsumed consequence.)  THE PROJECTION IS
%      ONTO A CLOSED CONVEX SET
%      AND THAT IS WHY THE CONTRACTION ARGUMENT SURVIVES: the round-7 wording
%      "metric projection onto an interval" was the sans-L = 2 instance of the
%      Hilbert-space fact, and eq:trunset is closed, convex and non-empty
%      (non-empty exactly because A4(i) requires eps_varrho < 1/sans-L).
%      WITHOUT IT the block has a real gap, and it is not a formality.  A4(i)
%      constrains varrho and says NOTHING about varrho-hat, so with an untruncated
%      (or simplex-truncated) first stage the event Wtilde-hat_{i,l} = 0 has
%      positive probability at each i and each l.  Two consequences, of different
%      sizes.
%        - TOTAL degeneracy, Wtilde-hat_i = 0 for EVERY i in I_m, makes T-hat
%          identically zero, on which event every statement about it -- including
%          prp:tccfails -- is vacuous.  This event is negligible under overlap and
%          cross-fitting, but "negligible" is a theorem, not a convention.
%        - PARTIAL degeneracy is NOT negligible.  For any first stage whose raw
%          fitted values leave [0,1] on a set of probability of order r_n^2 -- a
%          series linear-probability fit, say -- the expected number of main-fold
%          units with Wtilde-hat_i = 0 is of order n r_n^2 = n^{1 - 2 alpha},
%          which DIVERGES for every alpha < 1/2.  Each such unit annihilates
%          g-hat_{.,.,i}, so a row and a column of the Gram matrix eq:gram vanish
%          and lambda_min(G_n) = 0 EXACTLY.  prp:tccrepair(ii)'s nonsingularity
%          hypothesis then fails and eq:tccrad returns a radius of zero, so the
%          file's only demonstration that A6b is non-vacuous would have no content
%          on an event of non-vanishing probability.
%      The truncation removes the exact-zero mechanism DETERMINISTICALLY.  It does
%      NOT make lambda_min(G_n) bounded below -- near-collinear g-hat_i remain
%      possible, which is Q-3 and is a simulation question (probe CP-8).
%      DISPATCH-25 READING OF THE TWO BULLETS.  "Wtilde-hat_i = 0" is now a
%      statement about the VECTOR in R^{sans-L - 1}: a row and a column of
%      eq:gram's Gram matrix vanish iff ALL sans-L - 1 components vanish at that
%      i, since g-hat_{.,.,i} carries them together.  The mechanism is therefore
%      HARDER to trigger at sans-L >= 3 and the truncation still removes it
%      deterministically, so both bullets stand with their conclusions unchanged
%      and their probabilities only smaller.  Nothing here is claimed to improve
%      with sans-L; the point is that nothing degrades.
%
%  (f) AND A4(iii)'s AUXILIARY FOLD IS LOAD-BEARING FOR THE SAME PROPERTY, not
%      only for the rate.  Drop cross-fitting -- estimate varrho on the full
%      sample with any rule that interpolates -- and
%      varrho-hat_l(Delta Z_i) = 1{W_i = l} for every i and every l almost
%      surely, so T-hat is identically zero and prp:tccfails
%      is FALSE with probability one rather than on a null set.  The fold split
%      is therefore listed in prp:tccfails's hypotheses and not only in
%      def:hatT's convention.  (AT sans-L = 2 THE DISPLAY IS THE OLD ONE,
%      varrho-hat(Delta Z_i) = W_i, since 1{W_i = 1} = W_i there.)
%
%  d_Z is the EFFECTIVE dimension: Delta Z contains a constant, which
%  contributes nothing to a nonparametric regression, so d_Z <= p - 1.

\begin{assumption}[N1: Conditional centring of a weight]\label{ass:centering}
A debiasing weight $L$ \emph{satisfies N1} if
$E\big[L(\Delta U,\Delta Z)\mid\Delta Z\big]=0$ almost surely, \emph{the zero on the right being the
zero matrix of $\IR^{p\times(\mathsf L-1)}$}: N1 is $p(\mathsf L-1)$ scalar conditions, one per pair
$(j,l)$, and it is \textbf{not} a condition on any contraction of the $\mathsf L-1$ columns. For a
fixed $\mathsf v\in\IR^{\mathsf L-1}$ with $\mathsf v\ne0$ the contracted condition
$E\big[L\mid\Delta Z\big]\,\mathsf v=0$ is $p$ conditions; it is implied by N1 and implies N1
\emph{only} at $\mathsf L=2$, where the columns number one and the two statements are the same $p$
conditions. \cref{cor:centarity} makes the gap explicit and prices it. N1 is a property that a weight
may or may not have; it is not maintained, and a result that needs it says so.
\end{assumption}

\begin{assumption}[N1$'$: Conditional centring of a direction]\label{ass:centeringdir}
A perturbation direction $D$ in the weight \emph{satisfies N1$'$} if
$E\big[D(\Delta U,\Delta Z)\mid\Delta Z\big]=0$ almost surely.
\end{assumption}

% N1'S STATUS IS FIXED HERE, ONCE, AND EVERY OTHER SITE CITES THIS ONE.
% ROUND-7 DECISION.  Through round 6 the tag N1 carried THREE INCOMPATIBLE
% STATUSES across five sites -- a CONSTRUCTION here and in the glossary ("L IS
% NORMALISED so that ..."), an ASSERTION in the Discussion ("N1 costs nothing and
% BUYS orthogonality"), and an unused HYPOTHESIS on lem:secondord, lem:solset and
% lem:balance -- plus, at lem:secondord(iii), a fourth use in which the tag was
% applied to the DIRECTION D rather than to the weight L.
%
% THE COST OF THE CONSTRUCTION READING IS THAT IT VOIDS THE RESULT N1 EXISTS FOR.
% If L_star satisfies N1 BY DEFINITION then prp:orthog(iii)'s "only if" has no
% content, and the sentence "do not assume N1" in its statement is incoherent --
% one cannot suspend a definitional normalisation.  That is defect M-8, which
% round 6 repaired at prp:orthog and at no other site.  ONLY THE PROPERTY
% READING KEEPS (iii) NON-VACUOUS, and (iii) is the load-bearing claim of the
% whole rate architecture: it is what demotes varrho from a root-n nuisance to a
% product-rate one.  So the property reading is adopted, here, and:
%   - N1 is STRUCK from lem:secondord, lem:solset and lem:balance, where it was
%     listed and never consumed;
%   - N1 is NOT assumed anywhere in prp:orthog, which characterises it;
%   - the direction condition gets its own tag, N1', so that the tag N1 names
%     exactly one object;
%   - lem:density(iii) reads "the subspace of weights satisfying N1", which is
%     the property reading and needed no change;
%   - cor:udep continues to ASSUME N1, which is now a genuine hypothesis.
%
% WHAT REMAINS TRUE, AND IS WHY N1 IS CHEAP RATHER THAN FREE.  By def:weight the
% weight attached to a representer ell is determined by its anchor b, and
% re-anchoring at b = -E[ int_0^{Delta U} ell(v, Delta Z) dv | Delta Z ] produces
% a weight satisfying N1 with the SAME ell -- so it changes neither partial_u L
% nor, by (P3) and A1, the value of E[psi].  A2'(iii) shows that anchor is
% admissible.  So N1 CAN ALWAYS BE ARRANGED; it does not HOLD automatically, and
% the difference is exactly what prp:orthog(iii) measures.  Perturbing
% varrho -> varrho + tau_3 r gives d/d tau_3 E[psi] = -E[r(Delta Z) L], which
% equals -E[r(Delta Z) E[L | Delta Z]] and vanishes for every bounded r under N1
% and only under N1.  Without it varrho-hat needs a root-n rate.

\begin{assumption}[A5: Nuisance rates]\label{ass:rates}~
Let $h:=\hat\varphi-\vdagger$, let $\hat\delta_n$, $\delta_n$ and $\epsilon_n$ be the rates the clauses
below define, and let $r_n$ be the propensity rate of A4(iii).
\begin{enumerate}
\item[(i)] Operator noise, at a declared level with two branches. The level is the
\emph{deterministic} sequence
\[
\hat\delta_n\;:=\;\max\big\{\hat\delta_n^{(T)},\hat\delta_n^{(\varrho)}\big\},
\qquad
\hat\delta_n^{(T)}\;:=\;n^{-1/2}\sqrt{\ln n},
\qquad
\hat\delta_n^{(\varrho)}\;:=\;r_n ,
\]
so that $\hat\delta_n\ge n^{-1/2}\sqrt{\ln n}$ at every $n$; and the operator error is
$O_P(\hat\delta_n)$.
\item[(ii)] Structural rate, with the logarithm of \eqref{eq:phirate} carried \textbf{in the
symbol}. $\|h\|_{L^2(\pi)}=O_P(\delta_n)$ with
\[
\delta_n\;:=\;\big(\hat\delta_n\ln n\big)^{2\mu/(2\mu+1)} ,
\]
and $\delta_n=o_P(n^{-1/4})$. \emph{The first half is \textbf{assumed} and is not delivered by
\cref{prp:phirate}}: $h=\hat\varphi-\vdagger$ is formed at the estimator \cref{def:lfstep}\textup{(b)}
instructs, while that proposition concludes at the infeasible reading $\hat\varphi^{\pi}$, and the
passage between them is \suppref{Remark}{phirscope}\textup{(7)}, two of whose five parts are open. The
second half is a statement about the deterministic $\delta_n$ alone and \emph{is} delivered, under
\eqref{eq:frontier}.
\item[(iii)] \emph{Propensity frontier.} Writing $\hat\delta_n=n^{-r/(2r+d_Z)}$, which is
\textup{(i)}'s maximum for all $n$ large enough, A4(ii)'s $d_Z\ge1$ giving $r/(2r+d_Z)<\tfrac12$, 
condition (ii) holds if and only if
\begin{equation}\label{eq:frontier}
\frac{2\mu}{2\mu+1}\cdot\frac{r}{2r+d_Z}\;>\;\frac14
\qquad\Longleftrightarrow\qquad
r\;>\;\frac{d_Z\,(2\mu+1)}{2\,(2\mu-1)},
\end{equation}
which in particular requires $\mu>\tfrac12$, whatever the value of $r$.
\item[(iv)] \emph{Debiasing weight, on the derivative.} The estimator $\hat\ell$ of the representer, a
random element of $L^{2}(P_{\Delta U,\Delta Z})$ measurable with respect to the $\sigma$-algebra generated
by the auxiliary fold and read at one fixed \emph{jointly} measurable representative, satisfies
\[
\big\|\hat\ell-\ell_\star\big\|_{L^{2}(P)}=O_P(\epsilon_n),
\qquad \delta_n\,\epsilon_n=o_P\big(n^{-1/2}\big).
\]
% WHICH CONVENTION FIXES THAT FOLD AND THAT MEASURABILITY, and why naming what the
% measurability is OF is not decoration: assumptions_supp.tex, "A5 (nuisance
% rates)".
\item[(v)] \emph{Plug-in.} $\sqrt{\hat m}\,n^{-1/2}=o_P(\delta_n)$.
\item[(vi)] Sobolev stability of the first-stage error, at the final iterate.
$\|h\|_{H^{s}}=O_P(1)$, with $s$ as in H1.
\item[(vi$'$)] Sobolev stability along the linearisation path, as a w.p.$\to1$ envelope at a
\textbf{fixed} radius. With $\hat\varphi^{(0)},\dots,\hat\varphi^{(\bar k)}$ the stages of
\cref{def:lfstep} and $\bar k$ finite, there is one finite $\mathsf R$, not depending on $n$, such that
% BOTH MEMBERSHIPS ARE NEEDED for the display below to be a probability statement
% about a finite quantity: the stages lie in S_n once hat-varphi^{(0)} does, by
% def:lfstep(b),(d), and hence in H^s by eq:sieveinhs, while varphi_dagger is in
% H^s by H1's source set.  The accounting is in assumptions_supp.tex under "A5
% (nuisance rates)".
% ARCHAEOLOGY.  This clause named only the first membership through dispatch 20.
\[
P\Big(\max_{0\le k\le\bar k}\big\|\hat\varphi^{(k)}-\vdagger\big\|_{H^{s}}\;\le\;\mathsf R\Big)
\;\longrightarrow\;1 .
\]
\end{enumerate}
\end{assumption}

% THE FOLD READING OF A5(ii), A5(iv), A5(vi) -- DECLARED IN dispatch 11, IN
% results.tex def:crossfit(d), AND RECORDED HERE BECAUSE THIS IS THE CLAUSES' HOME.
% All three clauses above are stated at a SINGLE estimated nuisance: h is
% varphi-hat - varphi_dagger, ell-hat is one function.  Part I's estimator
% (def:crossfit) computes F of each, one per fold complement, so every Part-I
% result that cites them consumes them at max_{k <= F}.  THAT IS eq:foldrates,
% and it is NOT a silent widening:
%    A5(ii) first half and A5(vi) at the folds are DERIVED, in thm:rootn Step 0,
%      from (T-5) with F fixed and n_{-k} asymp n -- so they cost nothing;
%    A5(iv) at the folds is a STRENGTHENING IN FORM of the clause as written
%      (same order for fixed F, but a different statement), and it is therefore
%      carried by hypothesis (R-e.1) of thm:rootn EXPLICITLY, not read into this
%      clause.
% A4(iii) received the same treatment one clause earlier, in def:crossfit(a), and
% (R-a) and (R-e.2) already wrote max_k.  THE POINT OF RECORDING IT HERE: a reader
% who meets A5(iv) in this file and then meets it cited inside Part I must be able
% to see that the two are not the same statement and where the difference is paid.

% A5(vi) IS NEW IN ROUND 7 AND IS A TAG FOR AN INPUT THAT WAS ALREADY BEING
% CONSUMED UNTAGGED.  lem:balance(iii) applies the spline direct theorem to
% zeta = h = varphi-hat - varphi_dagger and needs ||h||_{H^s} = O_P(1); through
% round 6 that condition sat in the lemma's statement as the last untagged
% substantive hypothesis in Part F.  Two things must be said about it.
%   (a) IT IS NOT IMPLIED BY THE SOURCE CONDITION.  A2 and H1 constrain
%       varphi_dagger; h is an LF ITERATE, and its Sobolev norm is a property of
%       the regularisation path, not of the truth.  Recorded as N-18(a).
%   (b) ROUND-9 TEXT, NOW STRUCK: "IT IS CONSUMED ON ONE ROUTE ONLY.
%       lem:balance(ii), the rate route, does not use it..."  THAT IS FALSE FROM
%       STEP 6 ONWARD.  lem:pathcont needs the H^s bound to place the iterate in
%       the SOURCE-CONDITION ball V_n of def:srcball, and it needs it AT EVERY
%       STAGE, not at the final one.  So the Route column changes to "both" AND
%       the content changes -- which is why (vi') is added rather than (vi)
%       merely re-routed.
% A5(vi') IS NEW IN DISPATCH 7 AND IT IS NOT A RESTATEMENT OF (vi).  (vi) prices
% the FINAL iterate; lem:pathcont(c) and prp:phirate Step 5 consume the maximum
% over k = 0, ..., bar-k.  (vi) is (vi') at k = bar-k, so (vi') implies (vi) and
% not conversely.  bar-k is FINITE and does not grow with n (prp:phirate), so a
% finite union of O_P(1) statements is O_P(1) and (vi') is not a uniformity
% requirement over a growing family -- which is the only reason it is affordable
% as a maintained clause.
% WHAT (vi') IS NOT.  It is not discharged by prp:phirate.  prp:phirate delivers
% the L^2(pi) half of containment; the H^s half is assumed, here.  The
% alternative -- Bernstein's inequality for splines on S_n, giving
% ||zeta||_{H^s} <= c L_x^s ||zeta||_{L^2(pi)} -- would derive it at the cost of
% L_x^s delta_n = O(1), a genuine extra restriction that is NOT imposed.
% It is stated as an assumption rather than left inline so that the
% assumption-auditor sees a declared condition with a named consumer.

% A5(i) IS REPAIRED ON 2026-08-04.  THE TWO BRANCHES ARE DECLARED (":="), NOT
% BOUNDED ("O_P"), AND THE CLOSING SENTENCE IS STRUCK.  Accounting, quantifier
% and the five sites that consumed the floor: assumptions_supp.tex, "A5
% (nuisance rates)", subsection "Clause (i)".
%   THE DEFECT.  Through 2026-08-03 the clause read
%   "delta-hat_n^{(T)} = O_P(n^{-1/2} sqrt(ln n))" -- an UPPER bound -- while five
%   sites read it two-sidedly to obtain a LOWER bound on delta-hat_n, and
%   lem:opnoise(i) is a Chebyshev tail bound and cannot rescue the step either.
%   FROM X = O_P(a_n) ONE CANNOT INFER a_n <~ X.  The gap is what let an
%   implementation drop the T branch and stay formally consistent with the clause;
%   rate_delta_n in app_tsp_common.R did exactly that and printed a wrong number
%   in a shipped table (corrected 2026-08-03).
%   WHY ":=" IS THE RIGHT STATUS AND NOT A STRENGTHENING OF WHAT IS ASSUMED.
%   delta-hat_n is the noise level FED TO the discrepancy principle and must be
%   computable before the data are seen -- eq:lfsched builds the whole stagewise
%   schedule from it.  The superscripts are COMPONENT TAGS naming which summand
%   of the max is meant (glossary, parenthesised-superscript census, reading 4),
%   not estimators.  So the branches are deterministic sequences and "O_P" was a
%   category slip on them.  The block's own preamble already says the clauses
%   DEFINE these rates.  What the clause asserts ABOUT THE DGP is unchanged and
%   still upper-only: the operator error is O_P(delta-hat_n), which is
%   eq:opnoise's last equality and is DELIVERED by lem:opnoise, so every
%   "A5(i) is discharged, not assumed" downstream survives verbatim.
%   WHAT IS GAINED: the floor delta-hat_n >= n^{-1/2} sqrt(ln n) is now
%   ARITHMETIC -- a property of a max whose branches are named -- rather than a
%   two-sided reading of an O_P.
%
% A5(i) IS CONFIRMED BY lem:opnoise, AND ONE HALF OF IT IS SHARPER THAN WRITTEN
% -- BUT THE SHARPENING IS TUNING-DEPENDENT AND MUST BE READ WITH C5'(i).
%  (a) WHAT IS SHARPER.  The T-component's L^2(Pi) NORM is O_P(n^{-1/2}) with NO
%      logarithm (lem:opnoise(i)): L^2(Pi) is a Hilbert space, the summands are
%      bounded in norm, and a second-moment computation plus Chebyshev suffices.
%      No uniformity in (u,t) is taken, so no logarithm is bought.  The DECLARED
%      branch delta-hat_n^{(T)} = n^{-1/2} sqrt(ln n) therefore DOMINATES the
%      realised T-component with a sqrt(ln n) to spare -- which is what a
%      discrepancy-principle noise level is required to do (A7).
%  (b) IT IS NOT IDLE.  prp:stopfeas(i) consumes the sqrt(ln n) to turn a
%      second-moment bound into a FIXED-CONSTANT exceedance statement: Chebyshev
%      at s = tau n^{-1/2} sqrt(ln n) leaves O(1 / ln n) -> 0.  Delete the
%      logarithm and the statement holds only at a tau depending on epsilon.
%  (c) THE GUARD, AND IT IS WHY (a) IS NOT WRITTEN INTO THE CLAUSE.  The
%      SMOOTHING BIAS is O(h_u^{p_K}) and under the ROUND-9 tuning
%      h_u ~ n^{-1/16} that is exactly n^{-1/2} = o(n^{-1/2} sqrt(ln n)).  Under
%      the AMENDED tuning of C5'(i) below, h_u ~ delta-hat_n^{1/p_K}, the bias is
%      asymp delta-hat_n and is NO LONGER o(n^{-1/2}) -- it is the largest term
%      in eq:opnoise, by design.  So the two clauses interact: the T-component
%      closes at a fixed constant AGAINST delta-hat_n under EITHER tuning, and
%      against the sharper delta-hat_n^{(T)} = n^{-1/2} sqrt(ln n) only under the
%      round-9 one.  prp:stopfeas(i) is stated at delta-hat_n for that reason.
%      DO NOT HARD-CODE A BANDWIDTH INTO A5(i).

% A5(iv) IS PRICED ON THE DERIVATIVE, AND THAT IS A ROUND-5 CORRECTION.  Rounds
% 1-4 wrote it as ||L-hat - L_star||_{L^2} = O_P(eps_n).  BOTH CONSUMERS NEED
% THE DERIVATIVE INSTEAD:
%    lem:balance(i)-(ii) bound ||nu_g / pi|| by 2 bar-c_pi^{1/2} ||ell-hat -
%        ell_star||, because the first-order bias is generated by ell = partial_u L
%        and not by L;
%    lem:secondord(iv) bounds the (varphi, L) cross term by
%        2 bar-c_pi^{1/2} ||partial_u D|| ||h||, again the derivative.
% AN L^2 RATE FOR A FUNCTION DOES NOT TRANSFER TO ITS DERIVATIVE -- the same
% one-directional transfer lem:levels(iv) records for D^{-1}_0, here in the
% nuisance rather than in the structural direction.  So the round-4 A5(iv) was
% priced on an object no consumer uses, and stating the derivative version is
% not a strengthening for its own sake but a correction of the referent.
% NOTE what is NOT claimed: the LEVEL rate ||L-hat - L_star|| is not implied by
% the derivative rate either, and nothing at Step 5 needs it.  If a later step
% needs the level rate it must be added as a separate clause, not inferred.
%
% ROUND-7 POSTFIX: NO A5 CLAUSE IS ADDED FOR def:dirclass MEMBERSHIP, AND THE
% DECISION IS RECORDED HERE BECAUSE THIS IS WHERE A READER WILL LOOK FOR IT.
% lem:secondord(iii)-(iv) are stated at a FIXED admissible direction pair, and
% Step 6 or 7 will read them at D = L-hat - L_star, r = varrho-hat - varrho,
% h = varphi-hat - varphi_dagger.  That instantiation needs MEMBERSHIP in
% def:dirclass -- ||D||_{0,h} < infinity, and for (iv) differentiability with
% ||D||_{1,h} < infinity -- which is a DIFFERENT statement from any rate, and
% which A5 does not contain.  Three reasons it does not become a clause:
%   (a) [CORRECTED IN DISPATCH 11 -- THE ROUND-7 GROUND WAS FALSE AND THIS FILE
%       REFUTES IT NINETY LINES ABOVE.]  The round-7 text read "every clause of
%       A5 prices a RATE IN n; these are qualitative finiteness statements and an
%       O_P statement about them prices the wrong kind of object."  A5(vi) IS
%       ||h||_{H^s} = O_P(1) -- qualitative, no rate in n -- and A5(vi') is a
%       w.p.->1 envelope at a fixed radius.  So "A5 prices rates only" is not a
%       ground this block can use, and it was used twice more downstream
%       (results.tex rem:dirmember(a), rem:obligations) before dispatch 11 caught
%       it.  THE CORRECT DISCRIMINANT: every clause of A5 constrains an object
%       THE MAINTAINED BLOCK ITSELF INTRODUCES, at a single estimator -- h,
%       ell-hat, delta-hat_n, varrho-hat.  ||D||_{0,h} and ||D||_{1,h} are
%       seminorms of def:dirclass evaluated at a DIRECTION PAIR no maintained
%       clause fixes, so a clause about them would be a clause about a
%       definition's argument.  See rem:rootnhyp(5) for the same correction
%       applied to (R-a)-(R-e) of thm:rootn;
%   (b) both seminorms are indexed by (h, tau-bar) with h RANDOM, so a clean
%       clause would have to quantify over a deterministic class containing h
%       w.p.a.1 -- the ball V_n of L-19, which Step 6 constructs and which does
%       not exist in this block yet;
%   (c) it is IMPLIED, not assumed: rem:dirmember's eq:dirmember shows both
%       seminorms finite for EVERY h in V and every tau-bar as soon as L-hat is
%       an admissible weight (def:weight(b)) generated by a bounded, u-continuous
%       ell-hat.  Assuming an implied condition is the addition this round
%       removed everywhere else.
% SO WHAT A5(iv) PRICES IS ||ell-hat - ell_star||_{L^2(P)} AND NOTHING ELSE.
% SCOPE OF THE CLAIM -- ROUND-8 NARROWING (auditor Finding 7).  The round-7 text
% here read: "... so that NO BOUND IN THE BLOCK is denominated in a norm this
% assumption does not price."  READ LITERALLY THAT IS FALSE, and the file
% refutes it two hundred lines later: eq:crossshift is denominated in
% ||r||_infty and ||D||_{0,h}; lem:secondord(iv)'s SECOND bound was denominated
% in ||r||_infty (round-8 restates it in ||r||_{L^2(P)}, which A4(iii) does
% price -- Finding 5); and proofs.tex says in italics "Nothing in the assumption
% block prices ||D||_{0,h} or ||D||_{1,h}".  THE CORRECT, SCOPED CLAIM IS THIS:
%
%   In the D-DIRECTION of clause (iv) -- and there only -- lem:secondord(iv) is
%   printed in the norm A5(iv) prices, via eq:prodbound's FIRST inequality, with
%   ||D||_{1,h} shown as a stated weakening.
%
% The block makes NO block-wide claim that every bound in it is priced.  Three
% seminorms in Part F are deliberately unpriced -- ||D||_{0,h}, ||D||_{1,h} and,
% before round 8, ||r-hat - r||_infty -- and each is either a MEMBERSHIP
% condition on the direction class (the first two) or was a MIS-DENOMINATION
% since repaired (the third).
%
% The forward obligation -- a VERIFICATION, not a hypothesis -- sits in
% rem:dirmember, alongside A2'(i.b)/L-21 and N1' as the third of its kind; all
% three are inventoried in ONE PLACE, rem:obligations in results.tex, which is
% where a reader should be sent and where A5(iv)'s own presupposition (Finding 2)
% is now visible from the same page that says the obligation is open.

% A5 (iii) IS A FRONTIER, NOT A THRESHOLD.  The condition r > d_Z/2, which the
% memo headlines in four places, is the mu -> infinity LIMIT of
% \eqref{eq:frontier} and is NOT the operative condition; the correct
% derivation is the memo's own Sec.3.3 Step 3.  DO NOT WRITE r > d_Z/2.
%
% Smoothness in varphi_dagger and smoothness in varrho SUBSTITUTE along
% \eqref{eq:frontier}; the substitution is real and is worth displaying,
% because hiding it makes the block look weaker than it is:
%
%      mu       d_Z = 3           d_Z = 4           d_Z = 5
%      -----    ---------------   ---------------   ---------------
%      0.75     r > 7.5           r > 10            r > 12.5
%      1        r > 4.5           r > 6             r > 7.5
%      1.5      r > 3             r > 4             r > 5
%      2        r > 5/2           r > 10/3          r > 25/6
%      3        r > 2.1           r > 2.8           r > 3.5
%      infinity r > 1.5           r > 2             r > 2.5   (unattainable)
%
% THE mu = 2 ROW IS PRINTED AS EXACT FRACTIONS, DELIBERATELY.  Round 2 wrote
% "r > 3.33" against a true threshold of 10/3 = 3.333..., ROUNDING A STRICT
% LOWER BOUND DOWN: every r in (3.33, 10/3] passes the printed test and fails
% the frontier.  ("4.17" rounded 25/6 = 4.1666... UP and was safe, but the two
% must not be printed to different conventions.)  All other entries in the
% table are exact.
% Read the table by ROWS, not by its last line.  At mu = 1, d_Z = 3 the
% requirement is r > 4.5; at d_Z = 5 it is r > 7.5.  A propensity that
% is 7.5 times differentiable in five dimensions is not a condition a referee
% will grant, and the honest response is that the frontier is where the
% paper's feasibility lives: BUY mu, NOT r.  Under H1, mu = s/(2a), so buying
% mu means assuming varphi_dagger smoother or the problem less ill-posed.
% DISPATCH 25: THE FRONTIER'S EXPONENT IS UNCHANGED AND ITS BURDEN IS NOT.
% alpha = r/(2r + d_Z) carries no sans-L, so every number in the table above is
% the same number at every sans-L; and the vector-norm form of A4(iii) means
% r_n is one rate and not sans-L - 1 of them.  WHAT DOES CHANGE IS WHAT MUST BE
% GRANTED: A4(ii) now asks for the SAME order r on EACH of sans-L - 1 cell
% probabilities, so the smoothness hypothesis is sans-L - 1 conditions where it
% was one.  That is a real cost and it is not visible in alpha; the word
% "binary" was struck from the sentence above because it made the cost invisible
% by naming the wrong object.  The cost that IS visible in the application is
% the OVERLAP one, not the smoothness one -- see the plan's Sec.5, where the
% thin cell binds A4(i) before anything binds A4(ii).
%
% mu > 1/2 is a NEW NECESSARY CONDITION, absent from the memo's list.  Under
% H1 it is implied by plan Sec.2.6's own scope s > a + 1/2, with room:
% mu > 1/2 + 1/(4a).
%
% A WARNING FOR THE SIMULATION GRID (memo Sec.7, SG-20).  "A propensity
% estimator of order 2 clears n^{-1/4}" is a statement about delta-HAT_n, not
% about A5(iii).  A5(ii) constrains
% delta_n = (delta-hat_n ln n)^{2 mu / (2 mu + 1)}   [eq:deltalog, dispatch 10],
% which is LARGER -- larger by the power 2 mu/(2 mu + 1) < 1 of a quantity below
% one AND by the logarithmic factor.  SWEPT IN DISPATCH 11: this comment read the
% log-free form, which is the same staleness rem:phirscope(2) carried.
% At r = 2 the frontier reads
%     (2 mu)/(2 mu + 1) * 2/(4 + d_Z) > 1/4,
% i.e. mu > 5/6 at d_Z = 1 but mu > 7/2 at d_Z = 3.  This is consistent with
% the table above, whose mu = 3 row already demands r > 2.1 at d_Z = 3.
% QUOTING THE UNDEFLATED VERSION REPEATS, IN A NEW PLACE, EXACTLY THE MISTAKE
% THE r > d_Z/2 PROHIBITION EXISTS TO PREVENT.  SG-20's framing needs the
% deflation; flagged for whoever owns the grid.

\begin{assumption}[A8: Sieve window]\label{ass:sieve}
The balancing sieve $b_1,\dots,b_{L_x}$ is a \emph{B-spline} basis of \emph{degree exceeding
$s-\tfrac12$} on a quasi-uniform mesh over $\mathcal X$ \emph{with simple interior knots}, with
\[
n^{1/(2s)}\;\ll\;L_x\;\ll\;n^{1/2},
\]
and $\beta_1,\dots,\beta_{L_u}$ is a basis on the residual range with $L_u\le L_x$. \emph{Consequently
the sieve is contained in the Sobolev space the source conditions are written in}:
\begin{equation}\label{eq:sieveinhs}
\mathcal S_n\;=\;\operatorname{span}(b_1,\dots,b_{L_x})\cap\mathcal V\;\subset\;H^{s}(\mathcal X),
\end{equation}
so $\|\zeta\|_{H^{s}}<\infty$ at every $\zeta\in\mathcal S_n$, with $\mathcal S_n$ the sieve of
\cref{def:lfstep} and $H^{s}(\mathcal X)$ the \emph{classical} Sobolev space of H1, 
\emph{the one that carries no boundary condition}, which is what makes the containment available to a
basis that satisfies none.
\end{assumption}

% ---------------------------------------------------------------------------
% THE BASIS CLAUSE IS STATED IN *DEGREE*, AND THE HALF-ORDER IS NOT A ROUNDING.
% Through dispatch 19 this clause read "order exceeding s".  That is half an
% order short of eq:sieveinhs AND it is ambiguous, and the two defects compound.
%
% WHY IT IS SHORT.  eq:jackson is pinned to Schumaker (2007, Thm. 6.27) -- twice,
% and proofs.tex calls that pin "load-bearing" -- so Schumaker's convention
% governs: ORDER m = DEGREE m-1, and a spline of order m on a simple-knot mesh
% lies in C^{m-2}.  Its (m-1)-st derivative is a step function, which lies in
% H^sigma exactly for sigma < 1/2, so
%       spline of order m  \in  H^sigma   <==>   sigma < m - 1/2.
% Hence S_n \subset H^s <==> m > s + 1/2, while "order exceeding s" gives only
% m > s.  The gap {s < m <= s + 1/2} is EMPTY for integer s and for {s} < 1/2,
% and NON-EMPTY otherwise -- and proofs.tex's scope note on eq:jackson states
% outright that "s need not be an integer".
%
% THE WITNESS, AND IT IS ADMISSIBLE.  Take s = 1.7, which clears eq:lxwindow's
% s > 27/16 = 1.6875 at alpha = 3/8, the block's own boundary configuration.
% The old clause admitted m = 2 -- LINEAR splines, C^0 -- and m - 1/2 = 1.5 < 1.7,
% so S_n \not\subset H^{1.7}.  (A second witness, found independently: s = 2.6,
% m = 3.)  The consequence is not a weaker bound but an EMPTY EVENT: every iterate
% lies in S_n by def:lfstep(b),(d) and (P-a.1), and varphi_dagger \in H^s by H1, so
% ||varphi-hat^{(k)} - varphi_dagger||_{H^s} = +infinity at EVERY realisation and
% four printed conditions -- A5(vi), A5(vi'), (P-a.1) and eq:foldrates' second
% display -- held on the empty set rather than with probability tending to one.
% (P-a.1) was the sharpest case: "varphi-hat^{(0)} \in S_n AND
% P(||varphi_dagger - varphi-hat^{(0)}||_{H^s} <= R) -> 1" is two halves of ONE
% clause, jointly unsatisfiable at the witness.
%
% WHY THE REPAIR IS STATED IN DEGREE, AND WHY THAT IS CONVENTION-FREE.  The fix
% deliberately does NOT adjudicate whether the retired "order" meant Schumaker
% order or degree.  DEGREE is the quantity on which the two readings agree:
%       S_n \subset H^s   <==>   degree > s - 1/2,
% which is m > s + 1/2 in Schumaker's order and the SAME inequality on the SAME
% quantity under an order-means-degree reading.  So the amended clause is correct
% under either convention and the ambiguity cannot be re-opened.
%
% WHAT THIS IS *NOT*.  It is NOT a counterexample to prp:phirate or thm:rootn.
% At the witness configuration those results' hypotheses are unsatisfiable, so
% both hold VACUOUSLY; nothing proved in this corpus is false at s = 1.7.  What
% the witness exhibits is an admissible configuration at which a MAINTAINED
% assumption's event is empty -- a defect of the assumption block, not of a
% theorem.  Recorded here so that no later reader converts it into the stronger
% claim.
%
% eq:jackson IS NOW RESTATED IN DEGREE TOO, AND THAT REVERSES THE DISPATCH-20
% DECISION RECORDED IN THIS PARAGRAPH.  The paragraph read: "eq:jackson IS NOT
% TOUCHED ... eq:jackson's own lead-in in proofs.tex still reads 'order exceeding
% s'; that is eq:jackson's HYPOTHESIS and not a quotation of this clause, and it
% is left standing deliberately."  THE PREMISE IS RIGHT AND THE DECISION WAS
% WRONG IN CONSEQUENCE.  Before the amendment the two strings were IDENTICAL, so
% this clause implied eq:jackson's hypothesis trivially under either convention.
% Amending one and not the other put them in DIFFERENT UNITS, and the implication
% then held only under the Schumaker reading -- at s = 2.4 this clause admits
% degree 2 while "order exceeding s" read as a degree needs degree 3.  Two proof
% steps consume it and NEITHER converts (lem:balance Step 4's term (III);
% prp:phirate Step 1, which does not even cite eq:jackson by label), and the
% conversion was printed exactly once, at def:lfstep bullet 4, which does not
% consume it -- and in the looser non-strict form s <= d+1 against eq:jackson's
% strict d+1 > s, so the two printed statements of one hypothesis disagreed at
% s = d+1.  eq:jackson's lead-in now reads "degree at least s-1", the non-strict
% form, which is what the arithmetic below actually gives and what def:lfstep
% bullet 4 already printed; the two statements now agree at s = d+1 as well.
% THE MATHEMATICS WAS FINE THROUGHOUT; THE UNITS WERE NOT.  The arithmetic:
%   APPROXIMATION needs the spline space to have approximation order at least s.
%     A degree-d space reproduces polynomials of degree d, so its approximation
%     order is d + 1 (= Schumaker order m), and Thm. 6.27's rate L_x^{-sigma}
%     is available for sigma <= m.  So eq:jackson's L_x^{-s} needs
%         m >= s,  i.e.  degree >= s - 1.
%   MEMBERSHIP needs  degree > s - 1/2.
% MEMBERSHIP IS THE STRICTLY STRONGER OF THE TWO, by half an order.  So this
% clause is now driven by eq:sieveinhs and eq:jackson rides free with room:
% degree > s - 1/2 > s - 1.  The amendment can only give eq:jackson slack, and
% eq:jackson's exponent L_x^{-s} is consistent with the reading adopted here --
% AND with the other reading, since order-means-degree makes the clause stronger
% still.  NOTHING TO FIX; nothing was changed.
% THE ONE RESIDUAL CHECK, NOW PERFORMED AND RESOLVED.  The above turns on
% Thm. 6.27's smoothness cap being the ORDER m and not the degree m-1.  It is the
% order: the direct theorem runs through the m-th-order modulus of smoothness, so
% the cap is m = d + 1.  Hence APPROXIMATION needs degree >= s - 1 and MEMBERSHIP
% needs degree > s - 1/2, and MEMBERSHIP IS STRICTLY THE STRONGER, by half an
% order.  Had the cap been the degree, eq:jackson would have needed degree >= s
% and would have been the binding clause instead of membership -- it is not.  The
% resolution is printed at eq:jackson in proofs.tex, where the hypothesis is now
% stated in degree, so that no site has to convert between conventions.
%
% WHERE eq:sieveinhs IS CITED.  A5(vi) and A5(vi') above; (P-a.1) and
% eq:foldrates' second display in results.tex; and def:lfstep's fourth bullet,
% whose flag ("the proviso is A8's business") carried the whole burden alone and
% pointed at eq:jackson -- WRONG IN DIRECTION, since eq:jackson says how well
% splines APPROXIMATE H^s functions and what is needed is that splines LIE IN
% H^s.  Approximation power does not give membership.  The pointer is redirected
% to this clause, which is what def:lfstep's own first sentence already said.
% ---------------------------------------------------------------------------

% THE WINDOW IS NON-EMPTY IF AND ONLY IF s > 1, AND THE BASIS IS PART OF THE
% CONDITION.  The lower end comes from the approximation term: the first-order
% bias is O_P(n^{-1/2} delta_n) + O_P(||nu_g-hat|| L_x^{-s}), and
% L_x^{-s} = o(n^{-1/2}) requires L_x >> n^{1/(2s)}.  The upper end is the
% CONDITIONING constraint on the sieve solve and is BASIS-DEPENDENT:
%     splines       L_x << n^{1/2}   ==>  window non-empty iff s > 1;
%     polynomials   L_x << n^{1/3}   ==>  window non-empty iff s > 3/2.
% The memo prices the window without naming the basis; the basis is what
% decides whether the smoothness cost is s > 1 or s > 3/2, and it competes
% with A2, H1 and A5(iii) for the same s budget.  Splines are chosen.
%
% ROUND-7: THE TWO ENDS ARE ROUTE-DEPENDENT, AND THE BLOCK NOW SAYS SO.  N-18,
% delivered at Step 5 and confirmed by proofs.tex:lem:balance, is that A5(iv) and
% A8 are ALTERNATIVE sufficient conditions for the first-order bias, NOT
% cumulative ones:
%     RATE ROUTE      lem:balance(ii): A5(ii) + A5(iv) alone.  NO L_x condition,
%                     NO balancing constraint, NO eq:nuhat, NO A5(vi).  In
%                     particular A8's LOWER end -- and with it the "s > 1 /
%                     s > 3/2" line of the smoothness-budget table -- IS NOT USED.
%     BALANCING ROUTE lem:balance(iii): A8 in full + eq:nuhat in full + A5(vi),
%                     and NO rate for ell-hat.
% A8's UPPER end L_x << n^{1/2} survives on both routes: it is the conditioning
% constraint for the sieve solve and L-13 consumes it.
%
% DISPATCH-8: prp:phirate NOW CITES A8 BY CLAUSE, AND IT HAD TO.  Through
% round 10 prp:phirate listed "A8" in full among its hypotheses while its own
% Step 1 wrote "by A8's lower end and the spline direct theorem" -- two errors
% in one line.  The direct theorem needs the BASIS clause (B-splines of DEGREE
% exceeding s - 1/2 on a quasi-uniform mesh with simple interior knots), not the
% lower end.  DISPATCH-21 DESTALING: this parenthesis read "of order exceeding s"
% -- the pre-amendment wording of the clause, in the PRESENT TENSE, 120 lines
% below the amended clause and inside the same file.  The dispatch-20 sweep
% corrected four such quotations elsewhere and missed this one at home, which is
% the register defect the B1(vi) header already names: a clause amended in one
% place and quoted unamended in another is worse than one nobody amended.  The
% direct theorem's own hypothesis is DEGREE AT LEAST s - 1 and is implied by the
% amended clause with half an order to spare, so nothing in the sentence's
% CONCLUSION changes; only its quotation of the hypothesis does.  And A8's lower end
% together with (P-d) requires s > 1 / (2(1 - 2 alpha)), which at the boundary
% configuration alpha = 3/8 is s > 2 against the s > 1.75 the block recorded --
% AN EMPTY WINDOW.  prp:phirate now takes A8's BASIS and UPPER END and NOT its
% lower end, exactly as this comment already said the rate route should; (P-c)
% is the rate route's own lower end for L_x, calibrated to h_u delta-hat_n
% rather than to n^{-1/2}; and the joint line (P-c) ^ (P-d), eq:lxwindow, is now
% a row of the smoothness-budget table, as is the stronger both-routes line.
% THE FIX IS C2'(i)'s FROM DISPATCH 6, APPLIED TO A8: cite the clause, not the
% assumption.
% CONSEQUENCE FOR THE BUDGET TABLE: the "Sieve window non-empty" row is priced
% ROUTE-CONDITIONALLY there, since s is already contested against A2 (via
% mu = s/2a), H1 and A5(iii), and charging s > 1 unconditionally over-charges the
% one budget the paper's feasibility depends on.
%
% Conservative but honest, and still true on the balancing route: the third term
% is really bounded by ||nu_g-hat|| min(delta_n, L_x^{-s}), so even there the
% s > 1 requirement may relax further.  That refinement is not claimed.

% ---------------------------------------------------------------------------
% PART V.  REGULARISATION OF THE NONLINEAR OPERATOR EQUATION
%
% DISPATCH-7 RESTRUCTURE.  Through round 9 this part fixed the reading of
% A-hat (def:hatA) and then stated two maintained assumptions, A6 and A7,
% ABOUT AN OBJECT -- T-hat -- THAT NO FILE DEFINED.  Part H exhibited it;
% def:hatT below now DEFINES it here, where A7 needs it, so that the maintained
% block no longer names an undefined map.  The integrated kernel def:intker,
% which eq:hatT is written in, moves here with it; results.tex keeps
% lem:intker, which is where its four properties are derived from C3.
% A6 IS RETIRED -- see the retired-assumption block below and prp:tccfails.
% ---------------------------------------------------------------------------

\begin{definition}[The integrated kernel]\label{def:intker}
Let $C_{h_u}$ be the kernel of C3 and, for $I\in\mathfrak I$, define
\begin{equation}\label{eq:intker}
\mathsf C_{h_u}(v;I)\;:=\;\int_{-\infty}^{v/h_u}\Theta_{I}(\varsigma)\,d\varsigma
\;=\;\int_{-\infty}^{v/h_u}C_{h_u}\big(h_u\varsigma,I\big)\,d\varsigma
\;=\;\frac1{h_u}\int_{-\infty}^{v}C_{h_u}(\varsigma,I)\,d\varsigma ,
\qquad v\in\IR ,
\end{equation}
with $\Theta_{I}$ of \eqref{eq:rescaled}, and write $\mathsf C_{h_u}(v):=\mathsf C_{h_u}(v;\IR)$. The same definition with $h_x$ in place of $h_u$
defines $\mathsf C_{h_x}$. \emph{The boundary index is carried only in the $x$-direction}: in the
residual direction the argument $u$ ranges over all of $\IR$ and there is no boundary, so
\eqref{eq:hatT} uses $\mathsf C_{h_u}(\cdot)=\mathsf C_{h_u}(\cdot;\IR)$ throughout --- \emph{in $u$ the
interval is always $\IR$}; in the
$x$-direction $\mathcal X$ is a compact interval by A0 and the index is used.
Its properties are \suppref{Lemma}{intker}.
\end{definition}

\begin{definition}[The sample operator $\hat T$]\label{def:hatT}
Adopt the sample-splitting convention of \suppref{Lemma}{balance}: $\{1,\dots,n\}=\mathcal I_a\cup\mathcal I_m$,
$n_{\mathrm a}:=\#\mathcal I_a$, $n_{\mathrm m}:=\#\mathcal I_m$, the two folds independent by
A0, with $\hat\varrho$ computed on $\mathcal I_a$ as A4(iii) requires.
\emph{The split is \textbf{balanced}}: $n_{\mathrm a}$ and $n_{\mathrm m}$ are non-random with
\begin{equation}\label{eq:foldbal}
n_{\mathrm m}\;\asymp\;n
\qquad\text{(equivalently, }\liminf_n n_{\mathrm m}/n>0\text{)} .
\end{equation}
% EXTRACTED FROM THE STATEMENT OF def:hatT ON 2026-08-03, under the house rule
% (.claude/references/theorem-environment-style.md) that a definition carries
% CONSTRUCTION ONLY.  WHAT eq:foldbal COSTS, and the three sites that consume it:
% rem:hatTread(1).  eq:foldbal ITSELF STAYS, with its \label.
% ARCHAEOLOGY.  Through round 11 eq:foldbal was assumed nowhere.

\emph{Construction convention, and it costs nothing}: $\hat\varrho$ is taken to have values in the
$\varepsilon_\varrho$-interior of the simplex,
\begin{equation}\label{eq:trunset}
\Big\{q\in\IR^{\mathsf L-1}:\ q_l\ge\varepsilon_\varrho\ \text{ for every }l\in\{1,\dots,\mathsf L-1\},\ \ \textstyle\sum_{l=1}^{\mathsf L-1}q_l\le1-\varepsilon_\varrho\Big\},
\end{equation}
with $\varepsilon_\varrho$ the overlap constant of
A4(i); if a first-stage estimator does not, replace it by its metric projection onto that
set, which is well defined because \eqref{eq:trunset} is closed, convex and, since
A4(i) \emph{declares} $\varepsilon_\varrho\le1/\mathsf L$, non-empty --- \emph{the declaration alone and
not the condition}.
% THE TRUNCATION SET CARRIES NO SYMBOL, AND THAT IS THE C78 WRITE-IT-OUT
% PRECEDENT APPLIED FOR THE THIRTEENTH TIME.  cal-Q is a bound variable in
% prp:tccfails (the candidate linearisation) at ~30 sites, cal-D and cal-S are
% claimed, Delta is the differencing operator, and the glossary's own
% calligraphic census leaves only cal-O and cal-Y free -- neither of which is a
% natural name for a subset of R^{sans-L - 1}.  The set is named by its EQUATION
% LABEL, eq:trunset, which every site cites by \eqref; that is exactly how
% eq:hxwindow and eq:foldbal are cited and it costs nothing.
Writing $\widehat{\widetilde W}_{i}:=\big(\un\{W_i=1\}-\hat\varrho_1(\Delta Z_i),\dots,\un\{W_i=\mathsf L-1\}-\hat\varrho_{\mathsf L-1}(\Delta Z_i)\big)'$,
the \emph{sample operator} is
% WHAT THE TRUNCATION BUYS -- (a) it can only decrease the L^2 distance, so A4(iii)
% is preserved; (b) the resulting two-sided bound on |hatWtilde_i| -- and the scope
% note for a smaller known lower bound epsilon: rem:hatTread(2) and (3).
% THE DEFINITION OF hatWtilde_i DID NOT MOVE, AND MUST NOT.  Clause (b) opened by
% DEFINING the symbol and then derived a bound on it; eq:hatT immediately below
% USES the symbol, so cutting "(a) and (b)" wholesale would have left it
% used-but-undefined.  The definitional half stays here; only the derived bound
% epsilon_varrho <= |hatWtilde_i| <= 1 moved.
\begin{equation}\label{eq:hatT}
\hat T(\varphi,\theta)(u,t,l)\;:=\;\frac1{n_{\mathrm m}}\sum_{i\in\mathcal I_m}
\widehat{\widetilde W}_{i,l}\;\mathsf C_{h_u}\big(u-\Delta U_i(\theta,\varphi)\big)\;w\big(t'\Phi(\Delta Z_i)\big),
\qquad(u,t,l)\in\IR\times\mathcal T\times\{1,\dots,\mathsf L-1\} ,
\end{equation}
with $\Delta U_i(\theta,\varphi)$ the trial residual of \eqref{eq:trialres}. Its $x$-smoothed variant
$\hat T_{h_x}$, which is what \cref{def:lfstep} iterates on and what
$\hat{\mathcal A}_\varphi$ differentiates, is \suppeq{hatTx}.
\end{definition}

\noindent\emph{Notes on \cref{def:hatT}.} What the balanced split \eqref{eq:foldbal} costs and where it is
consumed, what the truncation buys, the scope note for a smaller known lower bound, and which assumption
names \eqref{eq:hatT} are \suppref{Remark}{hatTread}.

% ===========================================================================
%  DISPATCH-25 AMENDMENT TO def:hatT.  THREE THINGS MOVE AND ONE OBLIGATION IS
%  HANDED FORWARD.  The label, eq:foldbal, eq:hatT and the fold convention do
%  not move.
%
%  1. THE TRUNCATION SET.  [eps, 1 - eps] becomes eq:trunset.  AT sans-L = 2 THE
%     TWO ARE THE SAME SET: eq:trunset reads q_1 >= eps and q_1 <= 1 - eps,
%     which is the interval, so the amendment is invisible in the binary case.
%     THE CONTRACTION ARGUMENT SURVIVES BECAUSE THE SET IS CONVEX, not because
%     it is an interval -- metric projection onto a closed convex subset of a
%     Hilbert space is 1-Lipschitz, and the interval was the sans-L = 2 instance
%     of that.  A4(i) puts varrho itself in eq:trunset, so the projection can
%     only decrease ||hat-varrho - varrho||_{L^2}: A4(iii) is preserved, exactly
%     as before.  NON-EMPTINESS IS NOT AUTOMATIC AND IS WHY A4(i) BOUNDS
%     eps_varrho BY 1/sans-L: at eps_varrho > 1/sans-L the set is empty, the
%     projection is undefined, and A4(i) is itself unsatisfiable.  The binary
%     range (0, 1/2) was the same requirement and no round had to say so.
%
%  2. THE CONTRAST IS A VECTOR AND hat-T CARRIES A THIRD ARGUMENT.
%     hat-T(varphi,theta)(u,t,l) is the l-th contrast's component.  AT
%     sans-L = 2 there is one component, l = 1, and 1{W_i = 1} = W_i, so
%     eq:hatT is the printed display with a suppressed third argument.
%
%  3. WHY THE TRUNCATION SET IS THE SIMPLEX INTERIOR AND NOT THE BOX
%     [eps, 1 - eps]^{sans-L - 1}.  THE GROUND IS RESTATED IN ROUND 4, BECAUSE
%     ROUNDS 2-3 GAVE A GROUND THAT THE REVERSAL REMOVES.  They said: on the box
%     the implied base cell can be negative, so Omega-hat need not be positive
%     semi-definite and the NORMALISATION Omega-hat^{-1} would be undefined.  The
%     estimator no longer normalises, so that ground is gone.  THE DISCOVERY
%     STANDS ON A DIFFERENT AND SIMPLER GROUND: on the box hat-varrho NEED NOT BE
%     A PROBABILITY VECTOR AT ALL -- sum_l hat-varrho_l may exceed 1, i.e. the
%     implied hat-varrho_0 may be negative -- and a first stage that returns a
%     non-probability vector is not an estimate of the cell probabilities the
%     block is about.  eq:trunset is the set of vectors that ARE probability
%     vectors with every cell, base included, at least eps_varrho.
%     BE HONEST ABOUT WHAT IS AND IS NOT NEEDED: at the raw contrast the
%     componentwise surely-bounds eps_varrho <= |hatWtilde_{i,l}| <= 1 hold on
%     the BOX as well, since they use only hat-varrho_l in [eps, 1 - eps].  SO
%     THE SIMPLEX IS NOT LOAD-BEARING FOR THOSE BOUNDS; it is load-bearing for
%     hat-varrho being a probability vector, hence for hat-varrho_0 meaning
%     anything, hence for Omega-hat being a covariance matrix.  A dispatch that
%     needs only the two componentwise bounds could use the box; a dispatch that
%     reads hat-varrho_0, or Omega-hat, or reports min_l hat-varrho_l as the
%     overlap diagnostic, cannot.
%     WHAT FOLLOWS AND IS RECORDED WITHOUT A CONSUMER: on eq:trunset the
%     derivation of eq:omegafloor applies verbatim to Omega-hat_i, giving
%     lambda_min(Omega-hat_i) >= eps_Omega surely.  NO RESULT IN THE ADOPTED
%     ESTIMATOR CONSUMES THAT, and it is stated as an unconsumed consequence
%     rather than dropped, because it is what a normalised variant would need and
%     because the Section 6 diagnostic reports it.
%
%  4. eq:hatT CARRIES THE RAW CONTRAST, AND SO DOES psi.  ROUNDS 2 AND 3 PUT
%     Omega-hat_i^{-1} HERE; ROUND 4 REMOVES IT, with Omega-hat_i and the sample
%     normalised contrast removed from the printed text along with it, so that no
%     printed symbol is left without a consumer -- the defect round 1 was pulled
%     up for.  The author's ground for the reversal is at the placement block
%     under the notation paragraph.
%     *** CONSEQUENCE FOR THE sans-L = 2 CHECK, AND IT IS THE GOOD ONE:
%     AT sans-L = 2 eq:hatT IS THE PRE-DISPATCH DISPLAY AS AN OBJECT.
%     1{W_i = 1} = W_i, the third argument has one value, and NO FACTOR IS
%     INTRODUCED, so THE BINARY ESTIMATOR IS UNCHANGED and the engine's existing
%     scal convention is undisturbed.  BE PRECISE ABOUT WHAT "UNCHANGED" MEANS:
%     the ESTIMATOR is; the PRINTED DISPLAY is not typographically identical, and
%     "character for character" -- which round 4 wrote here and at the change log
%     -- OVER-CLAIMS.  eq:hatT prints a third argument l and an index set;
%     A0 prints {0,1,...,sans-L - 1} where it printed {0,1}; A4(i) prints a
%     quantifier over l where it printed an interval.  EACH OF THOSE REDUCES TO
%     THE OLD OBJECT AT sans-L = 2 AND NONE OF THEM IS THE OLD GLYPH SEQUENCE.
%     The correct claim, and the one the dispatch makes, is that the sans-L = 2
%     reduction has ZERO EXCEPTIONS AS MATHEMATICS. ***
%     BOUNDS ON THE RAW CONTRAST, SURELY, AT EVERY i AND EVERY l:
%       eps_varrho <= |hatWtilde_{i,l}| <= 1, by rem:hatTread(2)(b);
%       ||hatWtilde_i||_inf <= 1 and ||hatWtilde_i||_2 <= min{sqrt 2,
%           sqrt(sans-L - 1)}, the (P1) bounds, FREE of varrho and UNIFORM in
%           sans-L.
%     THE DEGENERACY GUARD IS THE COMPONENTWISE LOWER BOUND, exactly as it was
%     before this dispatch: eq:gram's row and column vanish iff hatWtilde_i = 0
%     as a vector, which |hatWtilde_{i,l}| >= eps_varrho excludes at every l.
%     No norm subtlety arises, because the raw contrast is nowhere sparse.
%
%  OBLIGATIONS HANDED TO PHASE 2, RECORDED HERE BECAUSE THIS FILE IS WHERE THE
%  THIRD ARGUMENT AND THE NORMALISATION APPEAR FIRST.  UNTIL PHASE 2 RUNS THE
%  CORPUS IS INCONSISTENT ACROSS THE FILE BOUNDARY.  THE FIRST DISPATCH-25 FORM
%  OF THIS NOTE SAID "AT EXACTLY ONE PLACE", WHICH WAS FALSE AND IS THE WORSE
%  KIND OF FALSE: naming two sites and asserting completeness invites the next
%  reader to stop looking.  THE SITES THIS DISPATCH FOUND, ALL IN
%  results_s2_identification.tex, AND THE COUNT IS A LOWER BOUND AND NOT A
%  CENSUS -- results_s3_estimation.tex, results_s4_asymptotics.tex,
%  results_s3_supp.tex, results_supp.tex and proofs.tex were NOT swept:
%    :47        def:model states W in {0,1} in the model display itself;
%    :155-165   def:operator's eq:Mdef and "with Wtilde = W - varrho(Delta Z)" --
%               two-argument, unnormalised, and eq:Tdef's range is L^2(Pi) over
%               R x T;
%    :177       def:admissible's restriction (A), at the scalar contrast;
%    :183-191   def:admissible's factorisation reading, quantified over
%               j in {0,1}, plus "when both W-cells are non-null" -- WHICH IS THE
%               VERY OBJECT THE NEW A1 NOTE CLAIMS AN EQUIVALENCE FOR, and whose
%               overlap proviso the corrected note (b) at A4 shows is not needed;
%    :201-202   eq:Adef's integrand;
%    :226-228   eq:kernel's g_{u,t}, which must become g_{u,t,l};
%    :282-290   the A4(i)-removal note at lem:bierens, whose CONCLUSION survives
%               and whose printed GROUND is binary-only;
%    :298       lem:bierens clause (b);
%    :340-356   cor:udep's c and the nu_{Delta Z + Wtilde c} family;
%    :367       prp:kerchar clause (b);
%    :377-378   rem:a1cwithdrawn's centring sentence;
%    :438-465   def:weight -- MISSED BY THE ROUND-2 CENSUS, INSIDE A FILE THAT
%               CENSUS DECLARED SWEPT, and it needs THREE changes, one of which
%               is already an outright contradiction with this file:
%                 :449  "a Borel b: R^p -> R^p" MUST BECOME R^p -> R^{p x
%                       (sans-L - 1)}.  A2'(iii) in THIS file already says
%                       R^{p x (sans-L - 1)}, so the two files currently
%                       contradict each other on the type of the same symbol;
%                 :439-440  "ell := cal-W ell^Pi ... applied coordinatewise"
%                       must become columnwise in l as well -- cal-W acts on each
%                       contrast column separately;
%                 :451-454  eq:weightdef's type follows: L^b_ell is
%                       R x R^p -> R^{p x (sans-L - 1)}.
%               AND def:weight IS ALSO WHERE THE WEIGHT-CLASS BIJECTION FAILS:
%               its clause (a) is a RANGE condition, ell in R(cal-W), which is
%               what refutes the round-2 redundancy argument.  Phase 2 must not
%               weaken (a) to a bounds condition in the course of retyping it.
%    :481-487   prp:orthog's eq:PsiL -- NOT eq:Psitau, which is a different
%               display at proofs.tex:3197; the round-2 census misnamed it -- and
%               the three perturbation directions, which must become a system in
%               (L_1, ..., L_{sans-L - 1}).  eq:PsiL's second line prints
%               "Wtilde L(...)", which is well formed at sans-L = 2 (both factors
%               scalar-by-vector) and ILL FORMED AT sans-L >= 3: the contrast must
%               stand to the RIGHT of the p x (sans-L - 1) weight.  THE CONTRAST
%               THERE IS AND REMAINS THE RAW Wtilde.
%  def:spaces (:64-137) declares L^2(Pi) over R x T and must gain the third
%  factor with B1(iv.d).
%
%  ONE SITE OUTSIDE results_s2:  results_s3_estimation.tex:216, eq:psidef, WHERE
%    psi IS DEFINED.  It needs the VECTOR contrast and the right-multiplication
%    order, psi = Delta Z Delta U + L(Delta U, Delta Z) Wtilde -- and NOTHING
%    ELSE.  ROUNDS 2-3 ALSO ASKED IT TO CARRY Wtilde^Omega; THAT REQUEST IS
%    WITHDRAWN WITH THE NORMALISATION.  At sans-L = 2 eq:psidef is already
%    correct as printed and needs no change at all.
%    THE NOTATION BLOCK DELIBERATELY DOES NOT RE-DECLARE psi: eq:psidef is a live
%    label in a compiled file, and a second \label{eq:psidef} here would be a
%    DUPLICATE that LaTeX resolves silently to whichever came last -- the same
%    silent-failure species the header records for \ref at an unnumbered
%    assumption.  The block prints an unnumbered reading and cites eq:psidef.
%    THAT FILE HAS NOT OTHERWISE BEEN SWEPT and certainly carries further sites.
%  It is recorded rather than patched here because those are Phase-2 files and a
%  half-migration across a file boundary is worse than a declared one.  eq:hatT's
%  third argument is written now, rather than deferred,
%  because def:hatT is the object A7 names in THIS file and A7 would otherwise
%  name a map of the wrong arity.
% ===========================================================================

\begin{definition}[The smoothed sample linearisation $\hat{\mathcal A}$]\label{def:hatA}
For $\varphi\in\mathcal V$ let $\mathcal A_\varphi:\mathcal V\to L^{2}(\Pi)$ be the Fr\'echet derivative of
$T(\cdot,\theta)$ at $\varphi$, so that $\mathcal A_{\mathrm{cond}}=\mathcal A_{\vdagger}$, with kernel
$\mathsf k(u,t,l;x)=\nu_{g_{u,t,l}}(x)/\pi(x)$ as in \cref{prp:compact}, where
$g_{u,t,l}$ is that proposition's $g_{u,t}$ read at the $l$-th contrast.
Write $\hat{\mathcal A}_\varphi$ for
the \emph{kernel-smoothed sample analogue} of $\mathcal A_\varphi$: the integral operator whose kernel is
obtained from $\mathsf k$ by replacing the indicator $\un\{\Delta U(\theta,\varphi)\le u\}$ inside
$g_{u,t,l}$ by its $C_{h_u}$-smoothed version, and each of the two products
$E[\,\cdot\mid X_j=x]f_{X_j}(x)$ constituting $\nu_{g_{u,t,l}}$ by its $C_{h_x}$ kernel estimate. Write
$\hat{\mathcal A}^\ast_\varphi$ for its adjoint in the $L^{2}(\pi)$--$L^{2}(\Pi)$ pairing, and
\[
\hat{\mathcal A}_{\mathrm{cond}}:=\hat{\mathcal A}_{\vdagger},
\qquad
\hat{\mathcal A}^\ast_{\mathrm{cond}}:=\hat{\mathcal A}^\ast_{\vdagger},
\]
the running Landweber--Fridman iterate being written $\hat{\mathcal A}_{\hat\varphi_m}$ wherever the base
point matters.

\end{definition}

% ===========================================================================
%  ROUND-5 AMENDMENT TO def:hatA: THE KERNEL GAINS THE CONTRAST INDEX.
%  It printed k(u,t;x) and g_{u,t}, which is the WRONG ARITY once B1(iv) in this
%  same file gives Pi a third factor: A_varphi maps into L^2(Pi) over
%  R x T x {1,...,sans-L - 1}, so A_varphi h is a function of (u,t,l) and its
%  kernel carries l.
%  THE POLICY THIS FOLLOWS IS THE FILE'S OWN, STATED AT def:hatT: eq:hatT's third
%  argument was written here rather than deferred to Phase 2 BECAUSE def:hatT IS
%  NAMED BY A7 IN THIS FILE, and a maintained assumption must not name a map of
%  the wrong arity.  def:hatA IS IN THE SAME POSITION -- it is named by the
%  printed retired A6 (ass:tcc, in assumptions_supp.tex), whose displayed event
%  contains A-hat_{varphi-tilde} -- so the same reason applies verbatim and
%  silence would be inconsistent with the policy.
%  WHAT IS NOT DONE HERE.  prp:compact, which supplies k and g_{u,t}, is in
%  results_s2_identification.tex and is Phase 2's; its own statement still prints
%  the two-argument form.  The amendment above is therefore a READING of an
%  object defined elsewhere, and the census under def:hatT already lists
%  results_s2:226-228 (eq:kernel's g_{u,t}) as a site Phase 2 must bring to
%  g_{u,t,l}.  UNTIL IT DOES, THIS DEFINITION AND prp:compact DISAGREE ON ARITY,
%  and that is one more entry in the declared cross-file inconsistency rather
%  than a new one.
%  AT sans-L = 2 NOTHING MOVES: the third argument takes one value, Pi_w is the
%  point mass at 1, and k(u,t,1;x) = k(u,t;x).
% ===========================================================================


\begin{assumption}[A7: Stopping feasibility, at the amended threshold]\label{ass:stopfeas}
The level used by the discrepancy principle dominates the true operator noise level with probability
approaching one: for some $\tau_{\mathrm{DP}}\ge1$,
\begin{equation}\label{eq:a7amend}
P\Big(\big\|\hat T(\vdagger,\theta)\big\|_{L^2(\Pi)}\le\tau_{\mathrm{DP}}\,\hat\delta_n\,\ln n\Big)
\;\longrightarrow\;1 ,
\end{equation}
with $\hat T$ as in \eqref{eq:hatT}. \emph{Any divergent factor serves in place of $\ln n$}; $\ln n$ is
fixed here so that one symbol, not a family, propagates downstream.
\end{assumption}

% A7'S THRESHOLD IS AMENDED IN THIS ROUND, AND THE AMENDMENT IS MANDATORY RATHER
% THAN OPTIONAL.  Rounds 1-9 wrote the threshold tau_DP * delta-hat_n.  Two
% facts force the change.
%  (a) IT IS NOT DERIVABLE AT A FIXED CONSTANT.  prp:stopfeas(i) closes the
%      T-component at a fixed tau_DP; prp:stopfeas(ii) shows the varrho
%      component does NOT, because A4(iii) is an O_P statement and an O_P
%      hypothesis carries no fixed-constant exceedance bound.  The alternative
%      repair -- strengthen A4(iii) to P(||varrho-hat - varrho|| > c r_n) -> 0 --
%      is also sufficient (prp:stopfeas(iii)) and is a strictly stronger
%      hypothesis on the propensity estimator, so the threshold moves instead.
%  (b) def:lfstep ALREADY IMPLEMENTS THE AMENDED LEVEL.  Leaving A7 as written
%      would put the maintained block in contradiction with the estimator.
%
% WHAT THE AMENDMENT COSTS, AND IT IS NOT NOTHING.  THE SUBSTITUTION IS AT THE
% LEVEL OF delta-hat_n, NOT OF delta_n -- SWEPT IN DISPATCH 11, because the old
% wording ("delta_n -> delta_n (ln n)^{2 mu/(2 mu+1)}") reads as a SECOND
% inflation once A5(ii) carries the logarithm in the symbol, which is exactly the
% one-symbol-one-object collision eq:deltalog was written to end.
% prp:stopfeas(iv) shows the substitution delta-hat_n -> delta-hat_n ln n, whose
% effect on the structural rate IS eq:deltalog's
% delta_n = (delta-hat_n ln n)^{2 mu / (2 mu + 1)}, is absorbed by the
% STRICTNESS of eq:frontier, so A5(ii)'s o_P(n^{-1/4}) survives.  IT IS NOT FREE
% AGAINST (P-f).  (P-f) is an EXACT domination, delta_n^2 = O(h_u^2 delta-hat_n
% ln n), and both sides now carry logarithms: the left carries (ln n)^{4 mu /
% (2 mu + 1)} and the right (ln n)^1, with 4 mu / (2 mu + 1) > 1 exactly when
% mu > 1/2 -- i.e. always, in this block.  The consequence is that (P-f)'s
% thresholds are STRICT: alpha (2 mu - 1)/(2 mu + 1) > 1/8 at C5'(i)'s bandwidth,
% and mu > 5/6 at the re-tuned bandwidth of C5'(i) as amended.  At the reference
% configuration (mu, d_Z, r) = (1, 3, 4.5), which sits exactly at the boundary in
% powers of n, (P-f) FAILS by (ln n)^{1/3}.  See rem:h6bandwidth.

% THE DISCREPANCY CONSTANT IS tau_DP, NOT tau -- ROUND-7 RENAME, AND IT IS THE
% CHEAP SIDE OF A COLLISION.  Part F uses tau = (tau_1, tau_2, tau_3) for the
% perturbation parameter of lem:secondord and prp:orthog, and tau-bar for its
% radius; A7's constant is a different object.  The perturbation parameter is
% used at roughly thirty sites and the discrepancy constant at two, so the
% constant moves.  The word subscript follows the eta_TCC precedent, which the
% glossary adopted for exactly this reason: tau_DP never appears without it.

% A7 IS STATED FOR THE RANDOM OPERATOR, and that is the point.
% hanke1995 proves the discrepancy principle for an EXACTLY KNOWN operator with
% NOISY DATA.  Here the noise is IN THE OPERATOR: T-hat is built from smoothed
% indicators and from varrho-hat, so T-hat(varphi_dagger) != 0.  This is the
% Dunker, Florens, Hohage, Johannes & Mammen (2014) noise model, which the
% project ADOPTS AS ITS SETTING even though it rejects IRGN as the regulariser.
%
% Without A7 THE DISCREPANCY PRINCIPLE NEVER STOPS: if the threshold sits below
% the noise floor, the residual never falls under it and the estimator is
% whatever the iteration cap bar-m_n returns, with no rate.  THAT IS EXACTLY THE
% FAILURE THE ROUND-9 THRESHOLD INVITED, and it is why def:lfstep(c) is
% STAGE-ADAPTIVE: the data error of the stage-k LINEARISED equation is not
% delta-hat_n but delta-hat_n ln n + h_u^{-2} eps_{k-1}^2, and a fixed threshold
% sits below it in exactly the regime the recursion must pass through.  A7
% supplies the FIRST of those two terms; the second is supplied by the
% deterministic schedule eq:lfsched, which the estimator computes.
% A7 is testable in simulation (probe CP-13): report the realised stopping index
% BY STAGE and how often the rule fails to stop before the cap, running the
% fixed and the stage-adaptive thresholds side by side.
%
% NOTE what prp:frechet(ii) supplies and what it does not.  It gives
% ||A_{varphi+h} - A_varphi|| <= C ||h||, i.e. the DETERMINISTIC operator's
% derivative map is Lipschitz.  A cone condition on the RANDOM operator T-hat
% does not follow from it -- which is now moot for A6 (retired) and is retained
% because it is the same distinction def:a6b turns on.

% ---------------------------------------------------------------------------
% PART VI.  REGULARITY CONDITIONS CARRIED OVER FROM THE DRAFT APPENDIX
%           (assapp1--assapp5, restated for the panel and for the new operator)
% ---------------------------------------------------------------------------

\begin{assumption}[C1: Densities]\label{ass:app1}~
Write $x=(x_1,x_2)\in\mathcal X^{2}$.
\begin{enumerate}
\item[(i)] The conditional density $f_{\Delta Y,X_1,X_2\mid\Delta Z,W}$ and the joint density
$f_{\Delta Y,X_1,X_2\mid\Delta Z}$ are $p_K$ times continuously differentiable and uniformly bounded
above.
\item[(ii)] $f_{X_1,X_2}$, $f_{X_1}$, $f_{X_2}$ and $\pi$ are uniformly bounded above and below away
from $0$ on $\mathcal X^{2}$ and $\mathcal X$ respectively; consequently
$\bar c_\pi:=\max_{j\in\{1,2\}}\|f_{X_j}/\pi\|_\infty<\infty$.
\end{enumerate}
\end{assumption}

% WHICH CLAUSE IS CONSUMED WHERE -- ROUND-7 NARROWING.  Every result proved
% through Step 5 that cites C1 consumes ONLY C1(ii): prp:frechet, prp:compact,
% lem:nusigma, lem:adjfubini, lem:adjoint, prp:kerchar, cor:a1cprim,
% cor:rangegap, prp:nophifree, lem:solset, lem:secondord, prp:orthog and
% lem:balance all use it for bar-c_pi, for ||1/pi||_inf, for ||pi||_inf, or for
% strict positivity of pi.  Those citations are now written C1(ii).
%   AND THAT SENTENCE WAS FALSE UNTIL 2026-08-23.  prp:frechet, prp:compact and
%   prp:kerchar still block-cited a BARE C1 on their hypothesis lines.  On this
%   block's register rule -- a block citation of the parent tag is a citation of
%   every clause under it -- that charged all three with C1(i); prp:compact is
%   reachable from thm:rootn, so C1(i) was, read by the rule, a hypothesis of a
%   theorem.  TWO LIVE CLAIMS DEPENDED ON THE NARROWING HAVING BEEN DONE: the
%   Part-VI register in assumptions_supp and lem:aerr's opening both state that
%   C1(i) is "a hypothesis of one lemma and, transitively, of no theorem".  The
%   three citations are now actually narrowed.  Each proof was checked first:
%   prp:frechet and prp:compact cite C1(ii) and nothing else, prp:kerchar cites
%   no C1 at all, and lem:nusigma -- which prp:kerchar leans on -- already read
%   C1(ii).  A note written to keep the usage matrix honest was, again, the place
%   the matrix was wrong.
% C1(i): WHERE IT IS ACTUALLY CONSUMED.  THE ROUND-7 FORM OF THIS NOTE WAS
% FALSE AND IS STRUCK.  It read: "C1(i) HAS NO CONSUMER BEFORE STEP 6 ... it
% is consumed by prp:phirate (Step 6) together with C3 and C5'.  It is retained,
% and this note is the forward reference that keeps the usage matrix honest."
% THE FORWARD REFERENCE NEVER LANDED.  prp:phirate's proof cites C1 at no line
% at all, neither clause, and neither lem:unifop nor lem:opnoise carries C1 on
% its hypothesis list.  A note written to keep the usage matrix honest was the
% one place the matrix was wrong, which is the defect species this block keeps
% producing and the reason the correction is printed rather than deleted.
%
% THE TRUE STATEMENT.  C1(i) is cited at exactly two lines in the corpus,
% proofs.tex Step 2 of lem:aerr, where the p_K derivatives of the products
% E[breve-g | X_j = .] f_{X_j} are what make the x-smoothing bias O(h_x^{p_K})
% at order p_K rather than at floor(s - 1/2).  That is its only consumer.
% AND THAT CONSUMER IS ITSELF CONSUMED BY NOTHING: lem:aerr opens "Consumers:
% none".  So C1(i) is a hypothesis of one lemma and, transitively, of no
% theorem.  IT IS NOT RETIRED, because lem:aerr's proof needs it and lem:aerr is
% retained for the two reasons its own reading (iv) gives -- it is the
% re-derivation C5'(ii) asks for, and it is the price-tag on the route Step 0 of
% prp:phirate DECLINES.  A clause whose only consumer is live is not in the
% ass:rank / ass:tcc / A2'(ii) / B1(vi) class and must not be filed there.
% The printed statement of all this is in the Discussion below, under the
% Part VI clause-status list, on the A5(v) model.
%
% CONSEQUENCE FOR prp:phirate'S OWN LIST, closed in the same pass: that list
% read "C1" in full and now reads C1(ii), which it consumes transitively
% through prp:hatAprop and lem:nusigma.  It does not consume C1(i), and the
% struck note above was the reason a reader would have thought it did.
%
% CHANGES FROM THE DRAFT'S assapp1:
%  (a) restated for the PANEL, with x = (x_1,x_2) -- as plan Sec.2.6b directs;
%  (b) conditioned on Delta Z as well as W, which A1 requires;
%  (c) the draft's smoothness order "d" is renamed p_K, the kernel order, both
%      to remove the collision with d_Z and because the two must match for the
%      bias order to be realised;
%  (d) the ratio bound bar-c_pi is made explicit because every boundedness
%      statement about A_cond and about the second-order form uses it, and the
%      memo's L-20 states it on the wrong object.

\begin{assumption}[C2$'$: Smoothness of the residual law]\label{ass:app2}~
\begin{enumerate}
\item[(i)] The conditional density $f_{\Delta Y\mid X_1,X_2,\Delta Z,W}(\cdot\mid\cdot)$ exists, is
$p_K-1$ times continuously differentiable in its first argument, and admits envelopes
$\bar f_k(x_1,x_2,z,\mathsf w):=\sup_{u}\big|\partial_u^{k}f_{\Delta Y\mid X_1,X_2,\Delta Z,W}(u\mid x_1,x_2,z,\mathsf w)\big|$
with $E[\bar f_k]<\infty$ for $k=0,1,\dots,p_K-1$ and $\|\bar f_0\|_\infty+\|\bar f_1\|_\infty<\infty$.

\item[(ii)] For every $t\in\mathcal T$ the map $u\mapsto\mathcal M(u,t)$ is $p_K$ times continuously
differentiable, with $p_K=8$ and
\[
\max_{0\le k\le p_K}\ \sup_{t\in\mathcal T}\ \bar w(t)^{-1}\sup_{u\in\IR}\big|\partial_u^{k}\mathcal M(u,t)\big|\;<\;\infty .
\]
\end{enumerate}
\end{assumption}

% THE DRAFT'S assapp2 IS INSUFFICIENT AND IS REPLACED.  It names only f_U.  Two
% distinct jobs need two distinct objects.
%
%  (a) L-4, the step that produces the Frechet derivative A_cond, needs the
%      density of Delta Y CONDITIONAL ON (X_1, X_2, Delta Z, W) -- i.e.
%      conditional on the ENDOGENOUS regressors -- and needs it C^1 in u with a
%      bounded derivative, one order MORE than the memo's L-4 requests.  The
%      extra order is what turns Gateaux differentiability into Frechet
%      differentiability with a QUADRATIC remainder (prp:frechet(i)) and what
%      makes the derivative map Lipschitz (prp:frechet(ii)).  Clause (i).
%
%  (b) The order-p_K kernel smoothing of the indicator has bias O(h_u^{p_K}),
%      and realising that order needs the SMOOTHED FUNCTION to have p_K
%      derivatives.  ROUND-2 CORRECTION: the smoothed function is NOT the
%      marginal f_{Delta U}.  What T-hat smooths is the CDF-level object
%      u |-> M(u,t), whose derivative is
%        partial_u M(u,t)
%          = E[Wtilde w(t'Phi(Delta Z)) f_{Delta U | X_1,X_2,Delta Z,W}(u)],
%      i.e. the Wtilde-w-WEIGHTED CONDITIONAL density; and the budget must be
%      uniform in t, since the operator norm is taken over the whole index set.
%      Round 1 put the budget on the marginal f_{Delta U}, which is neither
%      necessary nor sufficient.  Clause (ii) states it directly on M.
%      Clause (i), with envelopes to order p_K - 1, is a sufficient primitive
%      for (ii): differentiate under the expectation with dominating functions
%      bar-w(t) bar-f_k.
%
%  This also disposes of the round-1 "one derivative of slack" note, which was
%  computed on the wrong density and is withdrawn.

\begin{assumption}[C3: Kernel, and the defining property of the boundary adjustment]\label{ass:app3}
Fix $p_K=8$. \emph{Write $\mathfrak I$ for the set of closed intervals $I\subseteq[-\infty,+\infty]$
with $\inf I\in[-\infty,0]$, $\sup I\in[0,+\infty]$ and $\sup I-\inf I\ge2$ --- half-lines and $\IR$
included --- and, for a bandwidth $h>0$ and $I\in\mathfrak I$,}
\begin{equation}\label{eq:rescaled}
\Theta_I(\varsigma)\;:=\;C_{h}\big(h\varsigma,\,I\big),
\qquad\varsigma\in\IR,
\end{equation}
\emph{for the kernel read as a function of its scaled argument}. The smoothing kernel used in the residual
direction is a univariate generalized kernel $C_{h_u}$ with:
% EXTRACTED FROM THE STATEMENT OF C3 ON 2026-08-03, under the house rule
% (.claude/references/theorem-environment-style.md) that an assumption carries
% CONDITIONS ONLY.  C3 is very largely math displays, which ARE conditions, so
% only two things left it and the block is otherwise unchanged.
%   -- WHAT eq:rescaled IS FOR -- the object every bound below is stated for and
%      the object eq:intker integrates:                     rem:app3read(1).
%   -- CLAUSE (ii)'s over-long heading, which asserted two readings rather than
%      naming the clause:                                   rem:app3read(2).
% NO \label MOVED: eq:rescaled and eq:truncmom both stayed in the statement.
\begin{enumerate}
\item[(i)] Order $p_K$ at the interior index. Writing $\Theta_1:=\Theta_{\IR}$ --- \emph{the subscript
$1$ is an abbreviation for the untruncated index $I=\IR$, and it is the \textbf{digit} that is not an
interval and so not a member of $\mathfrak I$; the index it abbreviates, $I=\IR$, is a member,} the
notation being retained because it is the interior kernel every bound below names ---
\[
\int_{\IR}\Theta_1(\varsigma)\,d\varsigma\;=\;1,
\qquad
\int_{\IR}\varsigma^{\,j}\,\Theta_1(\varsigma)\,d\varsigma\;=\;0
\quad (1\le j\le p_K-1),
\qquad
\int_{\IR}|\varsigma|^{p_K}\big|\Theta_1(\varsigma)\big|\,d\varsigma\;<\;\infty .
\]
\item[(ii)] \emph{Boundary adjustment.} There is a measurable \textbf{envelope}
$\bar\Theta:\IR\to[0,\infty)$ with $\|\bar\Theta\|_\infty<\infty$, $\int_\IR\bar\Theta<\infty$ and
$\int_\IR|\varsigma|^{p_K}\bar\Theta<\infty$, such that for every $I\in\mathfrak I$ the function
$\Theta_{I}$ vanishes off $I$, satisfies $|\Theta_I|\le\bar\Theta$, and
\begin{equation}\label{eq:truncmom}
\int_{I}\varsigma^{\,j}\,\Theta_{I}(\varsigma)\,d\varsigma\;=\;\delta_{j0}
\qquad (0\le j\le p_K-1),
\end{equation}
with $\delta_{j0}$ the Kronecker symbol. Write
\begin{equation}\label{eq:envconst}
\mathsf c_2\;:=\;\|\bar\Theta\|_\infty ,
\qquad
c_C\;:=\;\frac{\|\bar\Theta\|_{L^{1}}}{2\,\mathsf c_2} ,
\qquad\text{so that}\qquad
\int_\IR\bar\Theta\;=\;2\,c_C\,\mathsf c_2 .
\end{equation}
\item[(iii)] \emph{Regularity at the interior index.} $\Theta_1$ is symmetric and Lipschitz, with
Lipschitz constant $\mathsf c_4$.
\end{enumerate}
The same three conditions with the same $p_K$ apply to the kernel $C_{h_x}$ used in the $x$-direction
inside $\hat{\mathcal A}_{\mathrm{cond}}$ and $\hat{\mathcal A}^\ast_{\mathrm{cond}}$, both as defined in
\cref{def:hatA}. \emph{C3 is not vacuous at $p_K=8$}: \suppref{Lemma}{bdryfam} exhibits a family
satisfying \textup{(i)}--\textup{(iii)} whose interior member is \emph{exactly} the order-$p_K$
Gaussian-based kernel the application computes.
\end{assumption}

% ============================================================================
%  SHORTENED 2026-08-28.  ~500 PRINTED WORDS OF ENVELOPE PROVENANCE MOVED TO
%  assumptions_supp.tex, paragraph "C3 (kernel and boundary adjustment)": the
%  monotonicity of the pair under pointwise domination, the full list of bounds
%  the two constants enter and its non-exhaustiveness, the c_18 non-case, the
%  closed-form majorant's numbers, the third (proved) envelope, and the ruling on
%  which pair a given reading should use.  ALL OF IT IS PROVENANCE OF CONSTANTS.
%  WHAT STAYED: that clause (ii) asks for AN envelope so the pair is not unique,
%  which pair the paper reports and its two numbers, and the two ARTICLE-SIDE
%  displays they feed (eq:normconst and eq:cap, both in results_s3_estimation).
%  An applied reader needs exactly those to read the cap.
% ============================================================================
\noindent\emph{Notes on C3.} What \eqref{eq:rescaled} is for, and how clause (ii)'s boundary index and
defining property are read, are \suppref{Remark}{app3read}. \emph{Clause \textup{(ii)} asks for
\textbf{an} envelope, so \eqref{eq:envconst} returns a different pair for each admissible $\bar\Theta$};
every result below holds at whichever pair the analyst can verify, and none is sharper for a larger one.
\suppref{Lemma}{bdryfam} discharges \textup{(ii)} at the \textbf{smallest} envelope the family admits,
the pointwise supremum $\sup_{I\in\mathfrak I}|\Theta_I|$, and \emph{its} constants at $p_K=8$ are the
ones reported here: $\mathsf c_2=34.0$ and $c_C=0.155$, so that
$2c_C\mathsf c_2=\|\bar\Theta\|_{L^{1}}=10.56$. \emph{Neither constant enters a rate}; they enter the
relaxation parameter \eqref{eq:normconst}, with $\mathsf c_2$ \textbf{squared}, and hence the cap
\eqref{eq:cap}, where the combination is $c_C\mathsf c_2^{4}$. Which further bounds they enter, how a
looser envelope moves each, and which of three admissible envelopes a given reading should use are in
the Discussion of the assumptions, under \emph{C3}.

% ==========================================================================
%  DISPATCH-26 RESTATEMENT OF C3: COMPACT SUPPORT OUT, ENVELOPE IN.  READ THIS
%  BEFORE TOUCHING c_C ANYWHERE.
%
%  THE FINDING.  Round-12 C3(i) required Theta_1 to be SUPPORTED IN [-c_C, c_C].
%  The kernel the project ships is k_order(., 8) of scripts/R/lf_panel.R -- the
%  order-8 Hermite/Gaussian kernel
%      Theta_1(s) = (105 - 105 s^2 + 21 s^4 - s^6)/48 * (2 pi)^{-1/2} e^{-s^2/2},
%  which is Schwartz and NOWHERE ZERO outside a discrete set.  NO FINITE c_C
%  MAKES THE OLD C3(i) TRUE OF IT.  So C3 was false of the estimator the paper
%  runs -- not loose, false -- and every result conditioned on C3 was, read
%  literally, about a kernel the code does not use.
%
%  WHY THE REPAIR IS AN ENVELOPE AND NOT A KERNEL SWAP.  Swapping in a compactly
%  supported order-8 kernel would leave C3 true as written and cost zero theory
%  edits, and it was rejected for two reasons.  (1) It re-runs every number in
%  Secs. 5-6: k_order/K_order/kern8 in lf_panel.R and its dependents, the bw_rot
%  recalibration (a [-1,1]-supported kernel has a different effective scale, so
%  the whole h_x ladder shifts), ~20 tables under paper/tables/, the phi-hat
%  bands, the four-cell localisation, and every claim-source-map entry.  (2) IT
%  BUYS ALMOST NOTHING, because the constructibility gap is not created by the
%  envelope route: C3(ii) ASSERTED a boundary family with a uniform sup bound and
%  a uniform p_K-moment and NO FILE EXHIBITED ONE, for a compact kernel any more
%  than for a Gaussian one.  lem:bdryfam is owed under either route.
%
%  WHAT KEEPS ~28 DOWNSTREAM SITES GLYPH-INVARIANT.  c_C is NO LONGER POSITED --
%  "and a finite c_C > 0" is gone from the opening line -- it is DERIVED at
%  eq:envconst as ||bar-Theta||_{L^1} / (2 c_2).  That is the single most
%  important line of this edit.  Every display of the form 2 c_C c_2 -- and there
%  are ~28 of them, in lem:intker, prp:hatAprop, lem:scaleinv's Schur test,
%  eq:normconst -- previously read "sup times support length" and now reads
%  "sup times (L^1 mass / sup)", i.e. the SAME NUMBER, ||bar-Theta||_{L^1}.  The
%  glyphs do not move and neither do the constants.
%
%  AND THE MAGNITUDE SURVIVES TOO, WHICH IS THE LOAD-BEARING SANITY CHECK -- BUT
%  ONLY AT THE MINIMAL ENVELOPE, AND THAT IS WHY THE MINIMAL ENVELOPE IS THE ONE
%  lem:bdryfam(d) FIXES.  The test is that 2 c_C c_2 -- the object ~28 displays
%  read -- keeps its numerical meaning: the retired half-width was ~2 against an
%  O(1) sup, i.e. 2 c_C c_2 ~ 7.  Minimal envelope: 2 c_C c_2 = 10.56.  PASSES.
%  Analytic majorant: 2 c_C c_2 = 1.64e9, EIGHT ORDERS OUT.  FAILS.  A round-27
%  draft of this block read the test off c_C alone (1.82 vs the half-width's 4)
%  and concluded the opposite; c_C alone is not what any of the 28 displays reads.
%  MEASURED 2026-08-16, AND THE CORPUS'S OWN L1 WAS 8.4% LOW: ||bar-Theta||_L1 is
%  10.5572, not the 9.74 a tanh-grid sweep returned, because the maximiser sits
%  where an interval ENDPOINT CROSSES sigma -- where |Theta_I| jumps -- which a
%  fixed (alpha,beta) grid meets only at isolated sigma.  c_2 = 34.039389 was
%  confirmed exactly, and independently: it is attained at |I| = 2 with ratio
%  1.06373 to the Legendre sum_{j=0}^{7}(2j+1)/2 = 32, above 1 by exactly what
%  exp(-sigma^2/2) <= 1 on I forces.
%
%  THE COST IS NOT "ONE BERNSTEIN EXPONENT" AND THAT UNDERSTATEMENT IS WHAT LET
%  THE ENVELOPE CHOICE LOOK COSMETIC.  c_2 enters eq:normconst SQUARED and eq:cap
%  to the FOURTH POWER; it enters eq:tccrad RECIPROCALLY.  At the
%  majorant, eq:cap is larger by 3.6e29 and prp:tccrepair(ii)'s certified radius
%  is smaller by 4.5e7 -- a SUFFICIENT CONDITION driven to vacuity, and
%  sharpness is its whole content.
%  c_18 IS NOT A SECOND CASE AND AN EARLIER DRAFT OF THIS BLOCK PUT IT HERE AT
%  2.4e16.  c_18 = 3K/(8 (4 c_C c_2 c_11)^2) with K = 4 bar-c_pi (2 c_C c_2)^2
%  ||pi|| ||1/pi|| (eq:convmean); both carry ||bar-Theta||_L1 SQUARED and it
%  cancels exactly, leaving c_18 = (3/8) bar-c_pi ||pi|| ||1/pi|| / c_11^2, free
%  of the envelope.  The 2.4e16 was computed by moving K's numerator and leaving
%  the denominator fixed.  Struck in results_supp on 2026-08-20 and left standing
%  here for one pass -- the same claim in two files, repaired in one.  Splitting c_2^u := ||Theta_1||_inf from
%  c_2^x remains DECLINED (it breaks glyph invariance at ~15 sites); taking the
%  envelope minimal buys the same orders for zero glyph movement.
%
%  THE FLOOR |I| >= 2, AND WHY IT IS AN ABSOLUTE NUMBER.  It keeps {a = b = 0} --
%  where the Gram matrix is 0 -- out of the closure of the compactified parameter
%  set, which is what lem:bdryfam(c) needs of it, AND (c) IS NOT ITS ONLY
%  CONSUMER: (g)'s Steps 4, 6 and 8 read it, Step 4 twice and the second time at
%  lambda_0 = 1, which needs |I| >= 2 outright; and lem:intker(iv)'s hypothesis
%  h_x <= (bar-x - underline-x)/2 IS this floor.  Through 2026-08-22 this block
%  said "used at exactly one place" and, below, "so the floor is free" -- both
%  written before (g) existed, and both false once it did.  It CANNOT be written as
%  a multiple of c_C, which is how lem:intker(iv)'s old hypothesis
%  h_x < (bar-x - underline-x)/(2 c_C) read: c_C is now defined FROM bar-Theta,
%  which is defined from underline-g, which is defined from the floor.  Writing
%  the floor in terms of c_C would close that circle.  The live configuration
%  delivers |I| = |cal-X|/h_x in [10.1, 12.4], clearing the floor by a factor of
%  five, so the floor BINDS NOTHING AT THE LIVE CONFIGURATION.  That is not the
%  same as being free: the VALUE 2 is spent in (g), and loosening it would force
%  eq:envbnds to be re-derived.
%
%  THE WEIGHT IS phi AND NOT Theta_1, AND THAT AMENDMENT IS NOT COSMETIC.
%  The natural construction -- Gram-Schmidt against the WEIGHT Theta_1 -- gives an
%  8x8 signed Hankel matrix whose (j,k) entry is the integral of s^{j+k} Theta_1
%  over I, and whose quadratic form is the integral of P^2 Theta_1.  THAT FORM IS
%  INDEFINITE: Theta_1 < 0 on 1.156 < |s| < 2.366 and on |s| > 3.75, and even over
%  the whole line the matrix carries mu_14 = -2,702,700 against mu_0 = 1.  So
%  nonsingularity is a codimension-one condition on a two-parameter noncompact
%  family; NO FLOOR ON |I| REPAIRS IT, and it must not be asserted anywhere.  (It
%  is also not REFUTED here; the phi weight makes the question moot, which is why
%  no symbol is coined for that matrix.)  Weighting by phi instead makes the Gram
%  matrix a genuine positive-definite Gram matrix on every nondegenerate interval
%  -- lem:bdryfam(a) is then two lines -- and, decisively, ITS INTERIOR MEMBER IS
%  EXACTLY THE SHIPPED KERNEL, by uniqueness rather than by convention.  That is
%  what makes the swap free: it delivers the C3(i)/C3(iii) compatibility for
%  nothing.
%
%  WHAT THE ROUND-12 BLOCK BELOW STILL GETS RIGHT, AND WHAT IT NO LONGER DOES.
%  Its (a), (b) and (d) stand verbatim: Theta is declared in the assumption, the
%  normalisation is printed, and (ii) has a defining property.  ITS (c) IS
%  SUPERSEDED.  (c) resolved a C3(ii)/C3(iii) clash by NORMALISING a scalar index
%  s_0 in [0,1] so that s_0 = 1 is the interior index and the support is
%  [-c_C s_0, c_C].  The interval indexing dissolves the clash rather than
%  normalising it away: I RECORDS WHICH ENDPOINT IS TRUNCATED AND BY HOW MUCH,
%  so the reflection convention of lem:intker(iv) becomes a THEOREM
%  (lem:bdryfam(f)) and then becomes unnecessary, and the BOTH-ENDPOINTS-CLOSE
%  case -- which a scalar s_0 could not express at all, and which lem:intker(iv)
%  excluded by hypothesis -- is covered by the same clause.  The old (iii)-clash
%  paragraph in proofs.tex (P9) is DELETED with it: the clash was an artefact of
%  the scalar index.
% ==========================================================================

% ROUND-12 RESTATEMENT OF C3.  THREE THINGS CHANGE AND EACH CLOSES A NAMED
% FINDING; NOTHING PROVED UNDER THE OLD C3 IS LOST, BECAUSE EVERY BOUND
% CONSUMED WAS |Theta_{s_0}| <= c_2 ON A SET OF LENGTH <= 2 c_C PLUS THE ORDER
% CONDITIONS, AND BOTH SURVIVE VERBATIM.
%
%  (a) Theta_{s_0} IS NOW DEFINED IN THE ASSUMPTION, eq:rescaled.  It was
%      previously introduced as a "standing convention (H2c)" in
%      archive/proofs_partH.tex -- stamped DO NOT COMPILE -- and used at six live
%      sites in proofs.tex with NO DECLARATION ANYWHERE IN A LIVE FILE.  That is
%      the assumption-auditor's U-1 and the regularity-checker's A8, and the
%      repair is to state the object where it belongs.  proofs.tex now carries
%      (P9) as a pointer to eq:rescaled rather than as an independent
%      declaration.
%
%  (b) int Theta_{s_0} = 1 IS PRINTED.  The old C3 said only "of order p_K" and
%      the normalisation is what "of order" means for a generalized kernel; it is
%      consumed at lem:intker(i), (iii) and (iv), so leaving it implicit charged
%      three clauses to a phrase.  The old (H2c) note (a) already said it should
%      be printed.
%
%  (c) THE BOUNDARY INDEX NOW HAS A DEFINING PROPERTY, AND THE C3(ii)/C3(iii)
%      CLASH IS DISSOLVED RATHER THAN RE-FILED.  Three incompatible meanings of
%      s_0 were in play: the OLD C3(ii)'s abstract index, whose support clause
%      [(s_0-1)/h, s_0/h] cap [-c_C, c_C] is VACUOUS for fixed s_0 in (0,1) and
%      small h and ONE-SIDED at s_0 = 1, contradicting (iii)'s SYMMETRY; the
%      literature's boundary distance IN BANDWIDTHS; and lem:intker(iv)'s proof,
%      which chose s_0 = (x* - underline-x)/(bar-x - underline-x), a FRACTION OF
%      THE INTERVAL and dimensionally wrong.  The new (ii) is the second of the
%      three, NORMALISED so that s_0 = 1 remains the interior index: support
%      [-c_C s_0, c_C], i.e. left truncation at the boundary distance measured in
%      bandwidths and divided by c_C.  At s_0 = 1 the support is exactly
%      [-c_C, c_C], TWO-SIDED, so (iii)'s symmetry is consistent and the clash
%      DISAPPEARS -- it was an artefact of the old interval, not a real ambiguity.
%      THE OLD VERDICT "nothing depends on which is meant" WAS FALSE, and it was
%      false at exactly one place: lem:intker(iv)'s proof uses the two-sided
%      reading at interior points and the one-sided reading at boundary points.
%      Under the new (ii) both are the SAME clause at different s_0.
%
%  (d) WHAT (ii) NOW SUPPLIES AND THE OLD C3 DID NOT: eq:truncmom, the TRUNCATED
%      moment conditions.  This is the defining property of a boundary-adjusted
%      generalized kernel and it is what lem:intker(iv) needs; the old C3 pinned
%      nothing of the kind, and (iv)'s proof appealed to "the rescaling implicit
%      in C3(ii)", which was not a statement.  LOAD-BEARING:
%      prp:hatAprop(iii) -- that A-hat* maps into V -- rests ENTIRELY on (iv),
%      and its own proof says that without it the identity fails by an O(1)
%      amount and A-hat* alpha leaves the parameter space, which (L2) inherits.
%
% The draft's assapp3a is otherwise retained, with three renamings: its order
% symbol "l" becomes p_K (l is one glyph from ell, the debiasing weight); its
% boundary-adjustment parameter "t" becomes s_0 (t is the Bierens index); and its
% support interval, written C_0 in round 2, is WRITTEN OUT as [-c_C, c_C].
% ROUND-7: THE SUPPORT INTERVAL LOSES ITS SYMBOL.  Calligraphic C_0 collided with
% cor:udep(c)'s admissible weight class -- a set of functions on R^p -- and the
% glossary carried both rows without noticing.  The kernel's support interval is
% a purely local object mentioned twice in one assumption, so it is written out,
% exactly as C12 wrote out the compact cube [-c_Phi, c_Phi]^p rather than naming
% it.  c_C joins c_Phi and c_Pi in the constant family.  The symbol C_0 is
% therefore free and belongs to cor:udep alone.
%
% THE DRAFT'S assapp3b IS RETIRED.  It introduces a discrete product kernel for
% smoothing over W, with bandwidth h_w, and a Dirac-delta representation of
% f_W.  Under the architecture of plan Sec.17, W enters ONLY through Wtilde
% (a multiplication) and, inside A-hat, through exact cell-splitting on the
% finitely supported W, which A4(i) makes feasible.  So h_w = 0 identically,
% assapp3b reduces to the indicator, and its symbol L_{h_w} -- which collides
% with the debiasing weight L -- disappears with it.  The Dirac device at draft
% `:257' goes too, removing a third meaning of Delta.
%   DISPATCH-25 SCOPE CORRECTION, AND IT WIDENS THE GROUND RATHER THAN NARROWING
%   IT.  The sentence read "exact cell-splitting on the BINARY W" through
%   dispatch 24.  The retirement never depended on there being two cells: what
%   it depends on is that W has FINITELY MANY cells, each of them conditionally
%   non-null, so that conditioning on {W = l} is exact averaging within a cell
%   and no smoothing in the W direction is defined, let alone needed.  A4(i)
%   delivers non-nullity at every l and at every sans-L, so h_w = 0 identically
%   at every sans-L and assapp3b stays retired for the same reason it always
%   was.  C5'(iv)'s "no bandwidth in W" is the same statement and needs no
%   amendment.  A DISPATCH THAT REACHES FOR A DISCRETE KERNEL BECAUSE THERE ARE
%   NOW FOUR CELLS HAS MISREAD THIS: discrete smoothing over W would be relevant
%   only if sans-L GREW WITH n, which A0 excludes.


% STATUS OF C4' -- PRINTED, NOT LEFT HERE.  ell_n and kappa_n occur nowhere in
% this project outside the clause above, and nothing is truncated to
% [-ell_n, ell_n].  The clause is INERT and is NOT RETIRED; the printed statement
% of that standing is in the Discussion, under "Five clauses of Part VI that no
% theorem rests on", which is the A5(v) model applied uniformly to C1(i), C4',
% C5'(ii), C5'(iii) and C5'(iv).
%
% THREE CORRECTIONS TO THE DRAFT'S assapp4.
%  (a) The draft requires kappa_n -> 0 AS PART OF THE ASSUMPTION.  That is not
%      an assumption anyone would want: it is a statement that the truncation
%      degenerates.  It also contradicts the same assumption's own "either
%      bounded or diverges slowly", since with ell_n bounded kappa_n is a
%      positive constant.  The requirement is struck.
%  (b) ROUND 2: kappa_n > 0 is a CONCLUSION and needs a hypothesis.  Continuity
%      alone gives an infimum over a compact set that is attained; strict
%      POSITIVITY of f_{Delta U} on [-ell_n, ell_n] is what makes it nonzero,
%      and it is now stated.  C2'(ii) gives smoothness of M, not positivity of
%      the marginal, so it cannot supply this.
%  (c) The tail caveat.  With ell_n bounded there is NO rate penalty.  If ell_n
%      diverges the penalty is SUB-POLYNOMIAL, not logarithmic: for Gaussian
%      f_{Delta U} and ell_n ~ sqrt(2 c log n), kappa_n^{-1} ~ n^{c}.  The plan
%      Sec.2.6 echo asserting a kappa_n^{-1} rate penalty is struck (memo Sec.9
%      item 11).

\begin{assumption}[C5$'$: Smoothing parameters]\label{ass:app5}~
\begin{enumerate}
\item[(i)] \emph{Residual direction, operator component.} $p_K=8$ and
\[
h_u\;\asymp\;\mathrm{sd}(\Delta U)\,\hat\delta_n^{1/p_K}\;=\;\mathrm{sd}(\Delta U)\,n^{-\alpha/8},
\qquad \alpha:=\frac{r}{2r+d_Z}\in\big(0,\tfrac12\big),
\]
% WHY h_u -> 0 AND h_u^{p_K} asymp hatdelta_n FOLLOW from the display above, and
% why no lower bound on h_u is required for this component: rem:app5read(1).
\item[(ii)] \emph{$x$-direction, derivative and adjoint components.} $h_x\le(\bar x-\underline x)/2$ at
every $n$, $h_x\to0$, and
\[
\big(n\,h_x\,h_u\big)^{-1}\ln n\;\longrightarrow\;0 .
\]
% THE FIRST CLAUSE IS "AT EVERY n" AND THAT IS THE WHOLE POINT OF ADDING IT.
% It is lem:intker(iv)'s hypothesis, which is what puts I(x*) of eq:ivalx in
% frak-I.  eq:smdiff, def:hatTexp and def:lfstep's measurability paragraph are
% PER-SAMPLE-SIZE statements: they need I(X_ij) in frak-I at the n they are read
% at, not eventually.  h_x -> 0 is an eventual statement about an object defined
% at each n, so it delivers the hypothesis only for n large, and the three sites
% inherited it silently.  IT BINDS NOWHERE: the live configuration gives
% |I| = |cal-X|/h_x in [10.1, 12.4], clearing the floor |I| >= 2 by a factor of
% five -- the same margin the comment at results_supp eq:ivalx already records,
% which is where this condition has been living as a comment rather than as a
% condition.
% EXTRACTED FROM THE STATEMENT OF C5' ON 2026-08-03.  Fourteen lines of WORKED
% ARITHMETIC stood here -- which lemma consumes the displayed form, the stronger
% inherited form and where it may be imposed, the (P-g) upper bound it must meet,
% and the resulting window on s.  A derivation is not a condition, so it is now
% rem:app5read(2).
% THE DISPLAY eq:hxwindow MOVED WITH IT, INTACT AND WITH ITS OWN \label.  It is
% cited at 12 live sites -- assumptions_supp (3), results_supp (5), proofs (1),
% and the undivided assumptions.tex/results.tex copies of those -- every one of
% them an \eqref, which resolves across environments and across files, and every
% one of them in the SUPPLEMENT or the corpus.  NOTHING in main.qmd or in the four
% shared files references it, which is what makes the move safe: the label leaves
% the article, and no article-side \eqref is left pointing at nothing.
\item[(iii)] \emph{Propensity.} $\hat\varrho$ satisfies A4(iii) with $r$ obeying the
frontier \eqref{eq:frontier}.
\item[(iv)] \emph{No further smoothing.} No bandwidth is used in $t$, in $\Delta Z$, or in $W$.
\end{enumerate}
\end{assumption}

\noindent\emph{Notes on C5$'$.} What clause (i)'s display buys, and clause (ii)'s stronger inherited form
together with the window on $s$ it leaves against (P-g), are \suppref{Remark}{app5read}.

% STATUS OF CLAUSES (ii), (iii), (iv) -- PRINTED, NOT LEFT HERE.
%
% THE HEADER OF THIS NOTE WAS FALSE THROUGH DISPATCH 20 AND IS CORRECTED HERE.
% It read "No hypothesis list in this corpus carries any of the three."  THREE
% LISTS CARRIED ALL THREE, by block citation of the parent tag C5': lem:aerr,
% prp:phirate and rem:rootnhyp(3) -- the last being thm:rootn's (T-5) inheritance.
% Dispatch 19's own register rule states that a block citation of the parent tag
% IS a citation of every clause under it; the same dispatch narrowed B1 and C1 on
% those lines by clause and left C5' un-narrowed.  All three now read C5'(i), and
% the header is true as of this dispatch and not before.
%
% (ii) is consumed ARITHMETICALLY, as the lower end of the admissible h_x window,
% at rem:tccwhy and lem:scaleinv, and lem:aerr reading (i) shows the operator-error
% bound needs only (n h_x h_u)^{-1} ln n -> 0, which is now what the clause states.
% (iii) is a POINTER to A4(iii) and eq:frontier, both cited directly wherever the
% propensity rate is needed.  (iv) is a SCOPE DECLARATION discharged by
% def:hatTexp and def:hatAcon.  NONE OF THE THREE IS RETIRED.  The printed
% statement is in the Discussion, under "Five clauses of Part VI that no theorem
% rests on" -- the A5(v) model, applied uniformly to C1(i), C4', C5'(ii),
% C5'(iii) and C5'(iv).
%
% assapp5 WAS WRITTEN FOR A DIFFERENT OPERATOR AND HAD TO BE REBUILT.  The
% draft's condition is  h_u, h_w, h_x -> 0  and  (n h_x^p h_u^3)^{-1} ln n -> 0.
% Three things are wrong with it here.
%
%  (a) THE EXPONENT p IS NOT dim(Z).  It was dim(x) in the cross-sectional
%      companion.  In the panel X is SCALAR and x = (x_1,x_2), so the exponent
%      is 2.  Plan Sec.1's reserved-symbol table lists draft `:302' under
%      "p = dim(Z)", which is a mis-attribution.
%
%  (b) THE OPERATOR NEEDS NO LOWER BOUND ON h_u.  T(varphi,theta)(u,t) is an
%      UNCONDITIONAL expectation of a bounded summand: there is no
%      local-averaging denominator, so the stochastic error is O_P(n^{-1/2})
%      pointwise and O_P(n^{-1/2} sqrt(ln n)) uniformly over (u,t), IRRESPECTIVE
%      of h_u.  The smoothing bias is O(h_u^{p_K}) under C2'(ii).  At p_K = 8
%      and h_u ~ n^{-1/16} the two balance exactly at n^{-1/2}.  This is where
%      clause (i) comes from, and it RECONCILES plan Sec.16.1's inherited
%      bandwidth choice with the theory rather than treating it as a MATLAB
%      artefact.
%
%      GUARD, REWRITTEN IN DISPATCH 7 BECAUSE THE OLD GUARD PROTECTED A MATCHING
%      TO THE WRONG TARGET.  ROUND-9 TEXT, STRUCK: "h_u ~ n^{-1/16} is simply the
%      largest bandwidth whose bias h_u^8 sits at the FLOOR n^{-1/2}."  Two parts
%      of that survive and one does not.
%       SURVIVES.  The exponent comes from matching bias to a level, NOT from a
%       bias-variance trade-off.  Because there is no local-averaging
%       denominator, the stochastic term does not inflate as h_u decreases, so
%       the standard balance h_u^8 against (n h_u)^{-1/2} -- which would give
%       h_u ~ n^{-1/17} -- DOES NOT APPLY HERE.  A later editor must not
%       "correct" the exponent by that route.
%       DOES NOT SURVIVE.  THE LEVEL TO MATCH IS delta-hat_n, NOT n^{-1/2}.  The
%       stochastic floor n^{-1/2} is what eq:opnoise's FIRST term sits at; the
%       operator noise the whole of Step 6 is denominated in is
%       delta-hat_n = max{n^{-1/2} sqrt(ln n), r_n}, and A5(i)'s FLOOR records
%       delta-hat_n >= n^{-1/2} sqrt(ln n) >> n^{-1/2} at every n and in every
%       configuration -- not merely where the propensity branch binds.  Matching
%       the bias to n^{-1/2} therefore chooses a bandwidth SMALLER than necessary
%       by a factor n^{(1 - 2 alpha)/16}, and pays for it: (P-f) of prp:phirate
%       is delta_n^2 = O(h_u^2 delta-hat_n ln n), which is HARDER the smaller
%       h_u is.  At h_u ~ n^{-1/16} it reads alpha (2 mu - 1)/(2 mu + 1) > 1/8 --
%       a JOINT restriction on (mu, alpha) that is IMPOSSIBLE at mu = 3/4 (it
%       would need alpha > 5/8 > 1/2) and FAILS at the BOUNDARY configuration
%       (mu, d_Z, r) = (1, 3, 4.5) once logarithms are counted.  At the amended
%       h_u ~ delta-hat_n^{1/8} it reads mu > 5/6, free of alpha.
%      WHAT THE AMENDMENT COSTS.  (1) The bandwidth now depends on alpha, hence
%      on r and d_Z, which the analyst must supply -- the same input A4(ii)-(iii)
%      already require, so this is not a new infeasibility, but it IS a new
%      dependence and is stated rather than hidden.  (2) The smoothing bias is no
%      longer negligible against the stochastic floor: it is asymp delta-hat_n,
%      which is all eq:opnoise needs (lem:opnoise(ii)) but which means
%      prp:stopfeas(i) closes at delta-hat_n and not at the sharper
%      delta-hat_n^{(T)}.  See the A5(i) comment, clause (c).
%      (3) mu > 5/6 STRICT joins the smoothness budget table.
%
%  (c) h_w DISAPPEARS (see the note under C3), and no bandwidth in t or in
%      Delta Z ever existed: the Bierens device replaces conditioning on
%      Delta Z by the scalar index t'Phi(Delta Z), and the propensity is a
%      SCALAR regression, priced separately in (iii).
%      DISPATCH-25 CORRECTION TO ONE WORD.  The propensity is no longer a
%      SCALAR regression: A4(iii) prices a MULTINOMIAL fit returning sans-L - 1
%      functions of Delta Z.  Nothing else in the clause moves -- the fit still
%      carries no bandwidth priced by C5', still enters only through eq:trunset
%      and eq:hatT, and is still priced by one rate symbol r_n, since A4(iii) is
%      stated in the vector norm.  The word is corrected rather than the clause,
%      because "scalar" was describing the ESTIMAND of the first stage and not
%      any smoothing parameter.
%
% Clause (ii) is CONSERVATIVE and is flagged as such.  The exponent 3 on h_u is
% inherited from the companion paper's density-derivative estimator; how many
% u-derivatives the Landweber step actually needs is prp:phirate's business
% (Step 6), and the exponent should be re-derived there.  A stronger condition
% is safe in the meantime.
%
% STATUS OF CLAUSE (ii): THE RE-DERIVATION WAS DONE IN DISPATCH 7, THE RELAXATION
% WAS DECLINED THERE, AND IT IS TAKEN NOW.  lem:aerr consumes
% (n h_x h_u)^{-1} ln n -> 0 -- ONE power of h_x and ONE of h_u -- and eq:gamman
% shows why: the exponents 2 and 3 pair with an estimator of a DENSITY DERIVATIVE
% in x and a SECOND derivative in u, which is the cross-sectional companion's
% object; A-hat_varphi estimates a density-weighted CONDITIONAL MEAN.
%
% THE DISPATCH-7 GROUND FOR KEEPING THE STRONGER FORM WAS AN INVALID INFERENCE
% AND IS STRUCK.  It read: "Clause (ii) as written IMPLIES the relaxed condition
% whenever h_x, h_u <= 1, so nothing proved is lost by keeping it."  The premise
% is true and the conclusion does not follow.  A LOWER bound that is stronger
% than needed is free only in isolation; it is NOT free once it must intersect an
% UPPER bound imported from elsewhere, because the strengthening moves the floor
% up towards a ceiling it does not control.  The ceiling here is prp:phirate's
% (P-g), R h_x^{s-1/2} = O(h_u delta-hat_n), and the intersection is EMPTY at
% admissible parameters.
%
% THE WHOLE COMPUTATION BELOW WAS RE-RUN 2026-08-10 (Phase 6a) WHEN (P-g) ACQUIRED
% ITS FACTOR h_u, WHICH IS THE TRIGGER THE LAST LINE OF THIS BLOCK ASKS FOR.  With
% h_u ~ n^{-alpha/8} the right-hand side is h_u delta-hat_n = n^{-9 alpha/8}, so
% every (P-g) ceiling below is h_x = O(n^{-9 alpha/(8(s - 1/2))}) and every
% threshold has its alpha-term multiplied by 9/8.  NOTHING BELOW CHANGED SIGN: the
% witness is still empty (by a wider gap) and both spot-checks are still non-empty.
%
% THE WITNESS, AND IT IS ADMISSIBLE.  (mu, s, a, r, d_Z) = (2, 1.2, 0.3, 8/3, 3),
% so alpha = r/(2r+d_Z) = 8/25 = 0.32 and kappa = 2 mu/(2 mu + 1) = 4/5.
%   FRONTIER  eq:frontier: kappa alpha = 0.256 > 1/4.               CLEARS.
%   eq:lxwindow: s > 9 alpha / (8 (1 - 2 alpha)) = 1.  s = 1.2.      CLEARS.
%   STRONG FORM (n h_x^2 h_u^3)^{-1} ln n -> 0 at h_u ~ n^{-alpha/8} = n^{-0.04}:
%       h_x >> n^{-0.44}.
%   (P-g) at s - 1/2 = 0.7:  h_x = O(n^{-9(0.32)/(8 x 0.7)}) = O(n^{-2.88/5.6})
%       = O(n^{-0.51429}).                     [was O(n^{-0.45714})]
%   WINDOW EMPTY, by the exponent gap 0.07429. [was 0.01714]
% The threshold is eq:hxwindow, s > 1/2 + 18 alpha/(8 - 3 alpha), which at
% alpha = 0.32 is s > 1.31818 and is FAILED by s = 1.2.  Emptiness holds on a band
% of s that is NON-EMPTY at every alpha in (1/4, ~0.359) -- and eq:frontier forces
% alpha > 1/4 always, so the band is never vacuous at the low end.  The band
% closes at alpha ~ 0.359 and is empty at alpha = 3/8, which is why the block's
% two spot-checks below missed it: both are run at s = 3 and s just above 2.
% (The band's upper root solves 9 alpha/(8(1 - 2 alpha)) = 1/2 + 18 alpha/(8-3 alpha)
% at ~0.35914; it was ~0.34940 under the old (P-g).  Both are printed rounded DOWN
% so that "non-empty at every alpha in the interval" stays true as printed.)
%
% WHAT THIS IS *NOT*.  It is NOT a counterexample to lem:aerr, prp:phirate or
% thm:rootn.  At the witness the block-cited hypothesis sets of those results were
% UNSATISFIABLE, so each held VACUOUSLY; nothing proved in this corpus is false at
% (s, alpha) = (1.2, 0.32).  What the witness exhibits is an admissible
% configuration at which a MAINTAINED clause and an imported proposition
% hypothesis have no common solution -- a defect of the conditions, not of a
% theorem.  Recorded so that no later reader converts it into the stronger claim.
%
% THE TWO REPAIRS, BOTH TAKEN.  (a) The three hypothesis lists that block-cited
% C5' now cite C5'(i), so no hypothesis set in the corpus contains the conflict.
% (b) The clause itself now states the weaker form, with the stronger one retained
% as an option, gated on eq:hxwindow -- so the MAINTAINED BLOCK read as a whole is
% satisfiable at every admissible (s, alpha) as well, which (a) alone does not
% deliver.  Both are needed and neither substitutes for the other.
%
% WHAT THE RETENTION OF THE STRONGER FORM IS STILL FOR.  It is the only place in
% the block where h_x carries a LOWER bound, and two printed computations read the
% h_x window at that end: rem:tccwhy(3) and lem:scaleinv's cap inflation.  Both are
% read at s = 3, where eq:hxwindow holds with room, so both stand unchanged.
% RECORDED, NOT DONE, AND STILL NOT DONE: promoting (P-g) into C5' as a clause.
% Whoever performs it must carry eq:hxwindow with it, since promotion puts the
% floor and the ceiling in the same maintained block.
%
% DISPATCH-8: THE WINDOW BELOW IS RECOMPUTED AT THE AMENDED BANDWIDTH.  The
% round-10 version of this comment read the window at h_u ~ n^{-1/16}, which
% C5'(i) RETIRED in the same pass, and at s > 1.75, which the joint (P-c) ^
% (P-d) ^ A8 line in the budget table now rules out as a DESIGN (it clears
% eq:lxwindow but not A8 in full).  Both numbers are replaced.
%
%   AT THE FRONTIER-CLEARING CONFIGURATION (mu, s, a, alpha) = (2, 3, 3/4, 0.32),
%   which is the admissible one:
%     h_u ~ n^{-alpha/8} = n^{-0.04}.
%     Clause (ii), (n h_x^2 h_u^3)^{-1} ln n -> 0, is
%         h_x >> n^{-1/2} h_u^{-3/2} (ln n)^{1/2} = n^{-0.44} (ln n)^{1/2}.
%     (P-g), R h_x^{s - 1/2} = O(h_u delta-hat_n), is
%         h_x = O(n^{-9 alpha/(8(s - 1/2))}) = O(n^{-2.88/20}) = O(n^{-0.144}).
%     WINDOW: (n^{-0.44}, n^{-0.144}) -- non-empty with room.  [ceiling was
%         n^{-0.128}; the window narrows at the cheap end and stays non-empty,
%         and lem:scaleinv's cap inflation is recomputed against it.]
%     The RELAXED lower bound lem:aerr consumes, (n h_x h_u)^{-1} ln n -> 0, is
%         h_x >> n^{-0.96} ln n, which is comfortable, so the relaxation buys a
%         much wider window and is available whenever it is wanted.
%
%   AT THE BOUNDARY CONFIGURATION (mu, d_Z, r) = (1, 3, 4.5), alpha = 3/8, kept
%   only as the clause-(0) illustration and NOT a design:
%     h_u ~ n^{-3/64}.  Clause (ii) gives h_x >> n^{-55/128} = n^{-0.4297};
%     (P-g) at s just above 2 gives h_x = O(n^{-9 alpha/(8(s - 1/2))})
%         = O(n^{-3.375/12}) = O(n^{-0.28125}).           [was O(n^{-0.25})]
%     WINDOW: (n^{-0.4297}, n^{-0.28125}) -- non-empty.  Under the RETIRED
%     bandwidth h_u ~ n^{-1/16} the lower end was n^{-13/32} = n^{-0.40625} and
%     the upper end, at the then-quoted s > 1.75, was n^{-0.3}; the amended
%     bandwidth WIDENS the window at the lower end and the corrected s
%     requirement NARROWS it at the upper end.
%
% Whoever promotes (P-g) into C5'(ii) must re-run these two lines, because both
% ends move with h_u and with s.


% ============================================================================
%  ROUND-2 CHANGE LOG (against proofs_review.md, 56/100)
%
%  C-1  B1(ii) restricted to w = exp; the cos/sin claim deleted; the general-w
%       route recorded as rem:swintercept in results.tex with the correct
%       anchor (Stinchcombe & White 1998 Thm 3.1, intercept-indexed).
%  C-2  A4 note (c) rewritten: the identity holds at sigma(Delta Z) ONLY; the
%       "same range, same kernel" gloss WITHDRAWN, with the exact extra
%       condition (p_X = varrho) named and rejected.
%  M-1  A2, A2' and A3's referent moved from nu to nu/pi, with the Picard
%       orthogonality test recorded as the check that settles it.
%  M-3  A1b(iv) added: nonsingularity of E[Delta Z Delta Z'] and the second
%       moments.  Identification of theta_dagger no longer routed through A3.
%  M-4  Discussion: A2's necessity described as conjectural in this nonlinear
%       problem, per the contract; the "no estimator is root-n" claim removed.
%  M-5  Domain of A_cond fixed to V (see def:spaces); H1's note added.
%  M-6  B1(iv) weakened to full topological support (Pi_u) and non-mutual
%       singularity (Pi_t) -- exactly what Step 5 consumes.
%  M-7  C2'(ii) restated on u |-> M(u,t) uniformly in t, not on f_{Delta U};
%       C2'(i) reduced to p_K - 1 derivatives, which is what it is used for.
%  M-9  C6 retagged A1c and added to the stable-tag list.
%  M-10 A4(i) removed from prp:wellposed / prp:frechet / prp:compact; A0's
%       non-use in this dispatch stated explicitly for the assumption-auditor.
%  m-4  assapp3's boundary parameter renamed s_0; varsigma and s_0 declared
%       reserved local bound variables in the glossary.
%  m-5  Tag-vs-number bridge stated in the header.
%  m-6  B1(v) now carries the explicit dominating function for differentiating
%       L under the t-integral.
%  m-7  C4' now assumes positivity of f_{Delta U}, which kappa_n > 0 requires.
%  F-1  comment corrected: |Wtilde| <= 1, A4(i) idle; cross-reference F-3 -> F-1.
% ============================================================================
%
% ============================================================================
%  DISPATCH-2 CHANGE LOG.  FOUR CHANGES; ONE IS A STRENGTHENING, ONE IS A
%  DELETION, ONE IS AN ADDITION, ONE IS A TYPE CORRECTION.  TAGS ARE UNCHANGED.
%
%  1. B1(iv.a) STRENGTHENED.  Pi_u must now be absolutely continuous with a
%     bounded, strictly positive density varpi_u.  This IMPLIES the round-2 full
%     topological support clause, so lem:bierens Step 5 -- (iv.a)'s only
%     consumer before this dispatch -- is unaffected and no existing proof
%     changes.  It is a condition on a CHOSEN reference measure, not on P.  It
%     is what closes escalation N-10 by construction: see lem:adjoint(ii)-(iii)
%     and convention (P7) in proofs.tex.
%
%  2. A2'(0) DELETED.  It is now a theorem, lem:adjoint(ii).  Clauses (i) and
%     (ii) KEEP THEIR NUMBERS, so every downstream citation is unaffected.  A2'
%     now names the relation ell_star = W ell^Pi_star in its own statement
%     rather than in a comment.
%
%  3. A1d ADDED, NOT MAINTAINED.  Instrument completeness given the residual; a
%     primitive sufficient condition for A1c, cited by cor:a1cprim alone.  The
%     analogue of Feve & Florens (2014, Prop. 2.2)'s "strongly identified by W".
%     It is added rather than left inline so that cor:a1cprim can cite by tag
%     and so that the assumption-auditor sees a single declared consumer.
%
%  4. TYPE CORRECTION (theory_memo.md finding F-15), affecting no assumption
%     statement.  L, ell, ell_star and cor:udep's c are R^p-VALUED, one
%     coordinate per coordinate of theta.  A2's own wording -- "each coordinate
%     of nu/pi" -- already implies it; the glossary typed them as scalars, which
%     made psi and A3's J dimensionally ill-formed.  Found by
%     scripts/symbolic/specs/jacobian-dims-numeric.
%
%     WHICH TAGS ARE COORDINATEWISE, CORRECTED IN ROUND 7.  A2'(i), A2'(iii) and
%     N1 are read COORDINATEWISE -- each is p copies of a scalar statement.
%     A2'(ii) is written with the Euclidean norm, which is the coordinatewise
%     statement summed.  A3 IS NOT COORDINATEWISE AND MUST NOT BE READ THAT WAY:
%     J = E[Delta Z Delta Z'] + E[Wtilde ell_star Delta Z'] is p x p precisely
%     because ell_star is a genuine p x 1 VECTOR, so that
%     E[Wtilde ell_star Delta Z'] is (p x 1)(1 x p).  A coordinatewise reading
%     returns p scalars and J is ill-formed.  The round-2 change log listed A3
%     with the coordinatewise tags and the glossary Sec.8.4 item 3 said the
%     opposite; the glossary was right.


%
%  ALSO RECORDED, in the A2 comment: A2 acquires a primitive reading through
%  lem:adjoint(iii), and MAY BE STRICTLY STRONGER than what the estimator needs
%  (F-20, recorded not claimed); and, in rem:Tcond, the first half of B1(v) is
%  IMPLIED by B1(iv.c) AT gamma = Pi_t (F-19).
%
%  ROUND-7 SCOPE CORRECTION TO F-19, AND IT RUNS THE OTHER WAY FROM THE ROUND-6
%  SUMMARY.  It is NOT the case that "the second half of B1(v) is now the only
%  live half".  F-19 retires the first half ONLY AT gamma = Pi_t, which is the
%  POPULATION-side identification of the correction measure made by
%  lem:adjoint(ii).  lem:density(ii) quantifies over ALL finite signed gamma
%  satisfying B1(v) and its proof CONSTRUCTS AN ATOMIC gamma != Pi_t, verifying
%  the first half against that gamma directly -- so the FIRST half is live on the
%  estimation side.  The SECOND half is live in rem:Tcond only, since cor:elldiff
%  declines it twice and says why (it dominates partial_u lambda, one
%  u-derivative short of what eq:Wderiv consumes).  BOTH HALVES ARE RETAINED, and
%  each for its own reason.
% ============================================================================
%
% ============================================================================
%  DISPATCH-3 ROUND-7 CHANGE LOG (fix-only; against regularity-audit
%  dispatch2-partF, notation-table.md, usage-matrix.md).  TAGS ARE STABLE:
%  N1' and A5(vi) are ADDITIONS, A2'(iii) is an APPENDED CLAUSE, and no existing
%  tag or clause number moves.
%
%  A0     "compact" -> "COMPACT INTERVAL" (F-9).  Free; four results consume
%         connectedness.  The lem:adjfubini appeal to A0 is STRUCK as
%         over-attribution: P is sigma-finite by def:model.
%  A1b(i) reduced to its content, varphi_dagger in V; the generic-varphi half is
%         def:spaces and not an assumption.
%  A2'    "admits a VERSION" added (m-3); (ii) written with the Euclidean norm;
%         NEW CLAUSE (iii), the anchor of L_star, which closes CRITICAL F-2 --
%         L_star was consumed by eq:Psitau, lem:secondord, prp:orthog(ii)-(iii),
%         A5(iv) and the glossary and was DEFINED NOWHERE.  def:weight in
%         results.tex supplies the definition; (iii) prices its anchor and
%         nothing else, and costs nothing at either anchor of interest.
%  A4(d)  variance referent corrected from E[Wtilde^2 ell_star^2] to
%         E[Wtilde^2 ||L_star||^2].
%  N1     restated as a PROPERTY of a weight, not a construction (auditor
%         Finding 1: three incompatible statuses across five sites).  Struck from
%         lem:secondord, lem:solset and lem:balance, where it was unused.
%  N1'    NEW TAG for the direction condition E[D | Delta Z] = 0, which round 6
%         cited as "N1 on the direction" -- a third object under one tag.
%  A5(vi) NEW CLAUSE tagging ||h||_{H^s} = O_P(1), the last untagged substantive
%         condition in Part F; single consumer lem:balance(iii).
%  A7     tau -> tau_DP, to free tau for Part F's perturbation parameter.
%  A8     window priced ROUTE-CONDITIONALLY, per N-18; budget table gains a
%         "Route" column.
%  C1     five citations narrowed to C1(ii); C1(i)'s forward reference to Step 6
%         recorded.  THAT FORWARD REFERENCE WAS FALSE AND IS STRUCK IN DISPATCH
%         18: Step 6 (prp:phirate) cites C1 at no line and its list is narrowed
%         to C1(ii); C1(i)'s one consumer is lem:aerr Step 2.  See the corrected
%         note at C1 and the Part VI clause-status list in the Discussion.
%  C3     support interval C_0 written out as [-c_C, c_C], freeing C_0 for
%         cor:udep(c) (notation collision T2-7).
%  A1c    head comment DESTALED: prp:kerchar proves the passage the head called
%         unavailable (m-11).
%  A3     removed from the coordinatewise list; it is a genuine vector statement.
% ============================================================================

% ============================================================================
%  ROUND-8 CHANGE LOG (fix-only; against notation-table_round8.md 0/100 and
%  usage-matrix_round8.md 7/100).  TAGS REMAIN STABLE: A2b and def:hatA are
%  ADDITIONS; no existing tag or clause number moves.
%
%  A1b(i)  MODE SENTENCE ADDED, and it is the fix for auditor Finding 1 (-30).
%          A1b(i) is consumed for WELL-POSEDNESS, never for a bound; a result
%          about a generic h in V consumes def:spaces.  Four idle citations
%          struck or redirected: lem:nusigma(ii) (proofs.tex + the statement
%          line), lem:levels (hypothesis line -- it appeared NOWHERE in the
%          proof), prp:nophifree (hypothesis line + roadmap sentence), and
%          cor:a1cprim (which inherited it).  The four genuine consumers are
%          NAMED in the clause.
%  A2b     NEW TAG, NOT MAINTAINED: "admissible representer", the PROPERTY half
%          of A2'(i).  Auditor Finding 4, which GATED STEP 4.  A2'(i) gains an
%          explicit status sentence, on the N1 model.  cor:elldiff is restated at
%          a generic admissible representer with ell^Pi_star as the
%          SPECIALISATION, so that L-21 knows which object it discharges.
%  A2'     CITED BY CLAUSE at the four Part F results (Finding 8): A2, A2'(i),
%          A2'(iii).  A2'(ii) is consumed by none of them and is a Step-8 clause;
%          this matches the C1 -> C1(ii) and B1(iv) -> B1(iv.a) narrowings the
%          block performed in round 7 for exactly this reason.
%  A5(iv)  RESTATED (Finding 2, -15).  The clause now prices
%          ||ell-hat - ell_star||_{L^2(P)} ONLY.  The first equality of the
%          round-7 display, ||partial_u(L-hat - L_star)|| = ||ell-hat -
%          ell_star||, IS AN ASSERTION AND NOT NOTATION: it presupposes
%          L-hat = L^{b-hat}_{ell-hat} per def:weight(b), which is precisely the
%          obligation rem:dirmember records as OWED.  A forward obligation had
%          migrated into the block unstated.  It is now stated as a CONDITIONAL
%          consequence, with the two consumers separated by which form each uses.
%  A5(iv)  SCOPE NOTE NARROWED (Finding 7).  The round-7 claim "no bound in the
%   note   block is denominated in a norm this assumption does not price" is
%          block-wide and FALSE as printed -- proofs.tex refutes it in italics
%          two hundred lines later.  Scoped to the D-direction of clause (iv).
%  A4(iii) is now the tag cited for the r-direction rate at lem:secondord(iv)'s
%          instantiation (Finding 5), replacing a mis-citation to A5(iv); and
%          the bound itself is re-denominated in ||r||_{L^2(P)}, the norm A4(iii)
%          prices.  NOTHING prices ||varrho-hat - varrho||_infty, and nothing
%          now needs it.
%  A6/C3   def:hatA ADDED (warden V-8).  The three symbols A-hat_{phi-tilde},
%          A-hat_cond, A-hat*_cond appeared in a MAINTAINED assumption and in C3
%          with no defining display in any file.  The hat convention alone is
%          insufficient because A_phi is a Frechet derivative, so "sample
%          analogue" is ambiguous between "the derivative of T-hat" and "a
%          separately smoothed operator" -- and C3, by attaching a bandwidth,
%          implies the second.  def:hatA FIXES THE SECOND READING and records the
%          construction as a Step-6 obligation.
%  B1(ii)  comment: Phi(+-1) = +-a written out as +-Phi(1) (warden V-6); a is the
%          degree of ill-posedness and Sec.1.6 of the glossary built the sans-serif
%          constant slot expressly to keep constants off that letter.
% ============================================================================

% ============================================================================
%  DISPATCH-14 CHANGE LOG (fix-only, against four narrow audits arriving
%  together).  TWO CLAUSES ARE RETIRED IN PLACE AND NO TAG OR CLAUSE NUMBER
%  MOVES.  No new maintained assumption.  No conclusion of any result changes.
%
%  A2'(ii)  RETIRED IN PLACE, on the ass:rank precedent -- retired because
%           PROVED, not because false and not because too strong.  Under A2'(i)
%           IN FULL with B1(i), (ii), (iv.a), cor:elldiff gives
%           ||ell_star||_inf < infinity and (P1) gives |Wtilde| <= 1, so
%           E[Wtilde^2 ||ell_star||^2] <= ||ell_star||_inf^2 < infinity.  The
%           clause had NO CONSUMER -- results.tex says so three times and this
%           file's own ROUND-8 CHANGE LOG says so -- and A4's comment (d)
%           justified retaining it on a ground that contradicts that log.  The
%           forward Step-8 need under tier T3, where A2'(i) is relaxed and the
%           clause becomes a genuine Khan-Tamer irregularity boundary, is a
%           FORWARD OBLIGATION and is recorded as one at the retirement note.
%           CLAUSE NUMBER KEPT as a struck placeholder, so a stale citation lands
%           somewhere.  Two DESCRIPTIVE mentions amended: rem:adjointreconcile(i)
%           in results.tex and its rendering in paper/sections/theorems/.
%
%  B1(vi)   RETIRED IN PLACE, and for the THIRD of the three retirement reasons:
%           neither "it is a theorem" (A3) nor "it is false" (A6) but "it is
%           about a different estimator".  T_n occurred NOWHERE outside its own
%           sentence; def:hatT integrates against Pi_t over ALL of T.  Its
%           heading claimed identification of varphi_dagger, which in this block
%           is A1b(iii), assumed outright.  The counterexample-hunter's
%           finite-grid witness is recorded WITH the retirement, so a later
%           dispatch reinstating the growing-grid design finds both together.
%
%  A1b(i)   CONSUMER LIST REPAIRED.  lem:biasexp was missing from a list the
%           clause declares BINDING ON EVERY CITATION, while lem:biasexp lists
%           and cites A1b(i) CORRECTLY -- it is what puts h_k in V and makes the
%           balancing pairing the object lem:nusigma(iii) computes.  A reader
%           following the list would have struck a correct citation: the
%           deletion-inviting direction.  THE LIST WAS WRONG AND THE CITATION WAS
%           RIGHT.  Found independently by the assumption-auditor and the
%           regularity-checker, the latter as "a citation gap and not a
%           mathematical one".
%           Three further sites (lem:prelimtheta Step 4, lem:epterm Step 3,
%           lem:sandwich Steps 3 and 6') were NOT added, because lem:nusigma(i)
%           is weakened this round from "for every h in V" to "for every h in
%           L^2(X,pi)" -- all its proof establishes -- after which none of the
%           three needs the tag at all.
%
%  A4(d)    TWO FALSE STATEMENTS STRUCK, one legitimate ground kept.  (1) "A2'(ii)
%           is the clause the second-order and bias terms consume" -- false, and
%           contradicted by this file's own round-8 log.  (2) The bare "(i) +
%           (iii)", read inside a comment attached to A4, attributes finiteness
%           of E[Wtilde^2 ||L_star||^2] to OVERLAP; A4(i) plays NO role in that
%           bound, whose tags are A2'(i), A2'(iii), A1b(iv), B1(i), (ii), (iv.a)
%           and (P1).  What survives, and is now the only ground stated, is that
%           overlap governs the SIZE of the variance: with the conditional laws
%           given W held fixed and varrho == eps, ranges are invariant (so A2 is
%           untouched) while E[Wtilde^2 ell_star^2] asymp eps^{-1}.
%
%  A4(ii)   d_Z >= 1 STATED, and it was assumed silently at four sites.  At
%           d_Z = 0 the exponent alpha = r/(2r+d_Z) is exactly 1/2, and (P-d)'s
%           sieve-entropy condition L_x ln n = o(n delta-hat_n^2) then reads
%           L_x = o(1) -- no integer L_x >= 1 satisfies it, eq:lxwindow is empty
%           at every s, and (T-5) is unsatisfiable.  Three further sites print
%           "alpha < 1/2" as a REASON and are false as printed although their
%           conclusions survive at the boundary.  d_Z = 0 is an EASIER problem
%           (a constant propensity) excluded for bookkeeping, not a pathology,
%           and the note says so.
%
%  H1       THE IDENTITY IS SUBSCRIPTED AND SCOPED: mu_{H1} := s/(2a) is the
%           source index of A2, at the POPULATION operator on ALL of V at
%           theta_dagger.  BARE mu is (P-a.3)'s PRIMITIVE index at the
%           finite-rank sieve compression of the REALISED operator at theta_o.
%           They are equal only under the interpolation proviso rem:pasrc(4)
%           DECLINES.  The two budget-table lines that compute mu_{H1} carry a
%           dagger note; the sentence deducing mu > 5/6 from s > a + 1/2 carries
%           the proviso; the flagship (mu,s,a,alpha) = (2,3,3/4,0.32) is read as
%           FOUR numbers of which two happen to satisfy s/(2a) = mu.
%           rem:pasrc(5) carries the binding reading; glossary Sec.8.8 documents
%           the hazard.
% ============================================================================

% ============================================================================
%  ROUND-9 CHANGE LOG (memo Sec.10 STEP 4).  ONE TAG IS ADDED AND NO EXISTING
%  TAG, CLAUSE NUMBER OR STATEMENT MOVES.
%
%  A2c     NEW TAG, NOT MAINTAINED: the source condition of index mu'' > 3/2 on a
%          representer, in the u-Sobolev scale on R with respect to LEBESGUE
%          measure, integrated against bar-w d Pi_t.  Modelled on A1d: a primitive
%          sufficient condition with exactly one consumer (lem:reprsmooth = L-21),
%          never asserted.  It DISCHARGES obligation (O1) of rem:obligations by
%          supplying a route to A2'(i.a)-(i.b) IN SUP-INSIDE FORM ON ell^Pi_star.
%          A2'(i) REMAINS MAINTAINED; what it loses is the status of an
%          unverifiable primitive.
%  A2'     comment blocks DESTALED at two places, both saying L-21 was unproved.
%          The (alpha)-versus-(beta) note and item 2 of "FOUR THINGS" now record
%          that L-21 is proved, that the threshold is mu'' > 3/2, and that the
%          sup-inside target is MET rather than approximated -- because A2c is
%          stated in the same bar-w L^1(Pi_t) geometry as the clauses it implies,
%          so no supremum is ever moved across the t-integral.
%  A2b     comment block CORRECTED at three lines that said "R1" where they meant
%          A2b -- the tag this block itself records as REJECTED because rem:R1
%          occupies it.  Found by sweeping the retired string across the file
%          rather than by any finding.  Glossary Sec.1.5b's rule applied to a
%          comment block: the round-8 rename reached every statement and citation
%          and not the prose that explains the rename.
%  BUDGET  the "Representer" line of the smoothness table now points at A2c and
%   TABLE  lem:reprsmooth instead of A2'(i), pins mu'' > 3/2 rather than leaving
%          the index free, and says in the following paragraph that the line
%          records a SUFFICIENT CONDITION and not an assumption.
%
%  WHAT IS DELIBERATELY NOT DONE.
%   - NO CLAUSE OF A2' IS WEAKENED OR REMOVED.  An assumption with a proved
%     sufficient condition is still an assumption; deleting (i) and keeping A2c
%     would make A2c maintained, which is exactly the A1c/A1d relation the block
%     already declines to collapse.
%   - NO NEW CITATION.  The three classical facts consumed at Step 4 -- the
%     one-dimensional Sobolev embedding, the numerical-radius identity for a
%     self-adjoint operator, and the essential-spectrum invariance named as Weyl's
%     theorem at lem:solset(iii) -- are each PROVED INLINE in proofs.tex, at the
%     level of detail the step consumes.  The reason is recorded rather than
%     silent: no spectral-theory or Sobolev-space entry resolves in this project's
%     bibliography (conway2000 was DEPINNED in round 6), and adding an
%     unresolvable \cite breaks the build.  The intended pins are named in a
%     comment at each site so that the bibliography pass has a target.
%     **THIS CLAUSE IS A RECORD OF THIS DISPATCH'S SCOPE AND IS NOW STALE AS A
%     DESCRIPTION OF proofs.tex.  AMENDED 2026-08-02 AND AGAIN 2026-08-03; the
%     amendment is here rather than by deletion, because the clause is what
%     licensed the choice at the time.**  What is now true, fact by fact:
%       * ~~no spectral-theory or Sobolev-space entry resolves~~ -- brezis2011,
%         halmos1982, conway1990 and kato1995 all resolve, every field and locator
%         read off an object on disk.
%       * ~~The intended pins are named in a comment~~ -- they are no longer
%         "intended".  Three are attached: conway1990 Thm II.1.5 at
%         lem:formstruct(i), which is why a FOURTH fact is no longer proved inline
%         at all; three integer-order sub-facts of lem:reprsmooth Step 1' to
%         brezis2011; and kato1995 Thm IV.5.35 at BOTH of lem:solset(iii)'s
%         previously bare "Weyl's theorem" invocations.
%       * The three facts listed above ARE still proved inline, and now for stated
%         reasons rather than for want of a book: the Sobolev embedding because
%         Brezis has no fractional space and no usable constant (M-3); the
%         numerical-radius identity because Halmos is complex-only and the
%         statement is false over R for a normal operator (M-5); and the
%         essential-spectrum CRITERION because Kato is complex-only by a global
%         convention and his nearest statement prices approximate nullity with a
%         normalized rather than an orthonormal sequence (M-4d).  The INVARIANCE,
%         by contrast, is pinned.
%     The register at the head of proofs.tex carries all of it, item by item.
% ============================================================================

% ============================================================================
%  DISPATCH-25 CHANGE LOG (2026-08-08).  PHASE 1 OF THE MULTINOMIAL-INSTRUMENT
%  PLAN, quality_reports/plans/2026-08-08_multinomial-instrument.md.
%
%  SCOPE.  This dispatch edits assumptions_main.tex and assumptions_supp.tex and
%  NOTHING ELSE.  results_s2_identification.tex, results_s3_estimation.tex,
%  results_s4_asymptotics.tex, results_s3_supp.tex, results_supp.tex and
%  proofs.tex are Phase 2 and 3 and are UNTOUCHED; the retired mirrors
%  results.tex, assumptions.tex and results_main.tex are out of scope by the
%  quality.md binary-gate repointing of 2026-08-05 and were not opened.
%
%  NO TAG MOVES.  NO CLAUSE NUMBER MOVES.
%  ~~NO \label IS ADDED TO OR REMOVED FROM ANY THEOREM-LIKE ENVIRONMENT, so the
%  shared counter is untouched and no number shifts in any of the eight compiled
%  files.~~  FALSE FROM ROUND 4 ONWARD AND CORRECTED IN ROUND 5.  ONE numbered
%  environment is added -- rem:htread, at the foot of assumptions_supp.tex -- so
%  every theorem-like number in results_s3_supp.tex, results_supp.tex and
%  proofs.tex ADVANCES BY ONE.  And FOUR numbered displays are added in this file
%  and THREE in the supplement, so equation numbers shift by four from
%  eq:contrast onward IN THIS FILE and, cumulatively, in every file after it:
%  results_s2, results_s3, results_s4 shift by four; assumptions_supp by four
%  until its own first insertion and by more thereafter; results_s3_supp,
%  results_supp and proofs by SEVEN.  ROUND 4 SAID ONLY "shifts every subsequent
%  equation number IN IT by four", which understated the reach: the equation
%  counter is global to the document, not per file.  NOTHING BREAKS -- every
%  cross-reference is an \eqref, a \suppeq or a \cref and resolves by label --
%  but a reader comparing printed numbers against an older PDF will find them all
%  moved, and that is this dispatch and not an error.
%
%  ONE CLAUSE IS APPENDED: B1(iv.d), the third factor of the range reference
%  measure.  Appending is the A2'(iii) / A5(vi) precedent.  Any list
%  block-citing B1(iv) acquires it; it is forced at sans-L = 2 and costs nothing
%  about P at any sans-L.
%
%  FOUR NUMBERED DISPLAYS ARE ADDED IN THIS FILE, WHICH SHIFTS EVERY SUBSEQUENT
%  EQUATION NUMBER IN IT BY FOUR AND BREAKS NOTHING, since every cross-reference
%  in this corpus is an \eqref or a \suppeq and resolves by label:
%     eq:contrast    the propensity vector and the contrast vector
%     eq:omegadef    the conditional contrast covariance
%     eq:omegafloor  A4(i)'s consequence, lambda_min(Omega) >= eps_Omega
%     eq:trunset     def:hatT's truncation set
%  THE COUNT IS FOUR AND THE LIST IS FOUR LONG.  ROUNDS 2-3 ADDED eq:omeganorm
%  AND eq:hatomeganorm HERE AND LEFT THE HEADING AT "FOUR" OVER A SIX-ITEM LIST.
%  BOTH ARE REMOVED IN ROUND 4: eq:hatomeganorm with the normalisation, and
%  eq:omeganorm because it MOVES to assumptions_supp.tex, where the reading it
%  belongs to now lives.
%  THREE DISPLAYS ARE ADDED IN assumptions_supp.tex -- eq:omegaprefac,
%  eq:omegainv and eq:omeganorm -- SO SEVEN IN TOTAL ACROSS THE TWO FILES.  THE
%  SUPPLEMENT'S OWN CHANGE LOG CARRIES THE SAME TWO NUMBERS; if the two logs ever
%  disagree again, that disagreement is itself the defect, not a rounding.
%  A build that reports a changed equation number is reporting this and not an
%  error.
%
%  ONE NUMBERED ENVIRONMENT IS ADDED, AND IT IS THE ONLY COUNTER CHANGE IN THE
%  DISPATCH: rem:htread, at the foot of assumptions_supp.tex beside the four
%  remarks already there.  That file is position 5 of 8, so every theorem-like
%  number in results_s3_supp.tex, results_supp.tex and proofs.tex ADVANCES BY
%  ONE.  No reference breaks -- all are \cref -- but printed numbers move, and
%  this is stated because the file's own contract requires it.
%
%  STATEMENTS AMENDED, IN FILE ORDER.
%   A0        support {0,1} -> {0,1,...,sans-L - 1}; sans-L declared finite,
%             fixed in n and known.  No consumer of A0 reads the W factor.
%   A1        UNCHANGED.  A printed unnumbered note is added after it recording
%             that the statement is cardinality-free and that this dispatch
%             therefore touches representation only.
%   B1(iv)    Pi = Pi_u (x) Pi_t -> Pi_u (x) Pi_t (x) Pi_w.
%   B1(iv.d)  NEW.  Pi_w a probability measure on {1,...,sans-L - 1} with
%             strictly positive masses.
%   A2'(i.a)  ell^Pi_star gains the contrast argument l; the budget is the
%             MAXIMUM over l.  Choice flagged in the comment beneath A2'.
%   A2'(i.b)  same, and the differentiability is quantified over l.
%   A2'(iii)  b_star's CODOMAIN R^p -> R^{p x (sans-L - 1)}; the norm is
%             Frobenius.  Its DOMAIN is untouched.
%   A2b       the same two budgets at a generic representer.
%   A2c       "p coordinates" -> "p(sans-L - 1) coordinates"; the slice and the
%             budget carry l; eq:lambdamu aggregates by the maximum.
%   A4(i)     one-sided floor at every cell, eps_varrho in (0, 1/sans-L), plus
%             the CONSEQUENCE eq:omegafloor.  At sans-L = 2 the conjunction over
%             l in {0,1} IS the printed two-sided interval.
%   A4(ii)    per component, common Holder order r.
%   A4(iii)   hat-varrho is the vector fit; the L^2 norm is spelled out as the
%             vector one.  Same single rate symbol r_n.
%   def:hatT  truncation set eq:trunset; hatWtilde_i is a vector; eq:hatT gains
%             the contrast argument l AND the normalised contrast
%             ~~eq:hatomeganorm, which is where Omega^{-1} is placed (ROUND 2)~~
%             -- REMOVED IN ROUND 4 with the normalisation.  eq:hatT carries the
%             RAW contrast and at sans-L = 2 is the pre-dispatch display.
%
%  ROUND-2 ADDITIONS (against theorist-critic 61/100).
%   C1  Omega^{-1} placed on the operator.  ~~and psi left at Wtilde on a
%       redundancy argument~~ -- THE SECOND HALF IS WITHDRAWN IN ROUND 3.
%   C2  THE DEGENERATE REGION WAS WRONG AT sans-L >= 3 AND IS CORRECTED.  Round 1
%       wrote D = {det Omega = 0} = {min_l varrho_l = 0}; the correct set is
%       {Omega = 0} = {max_l varrho_l = 1}, and the equivalence (a) <=> (b) needs
%       NO overlap condition at any sans-L.  Two under-claims that followed from
%       the error -- the A1 note here and the A4 paragraph in the supplement,
%       both of which listed A4(i) as a hypothesis of the equivalence -- are
%       corrected with it.
%   M4  the overidentification motivation is restated on the joint-consistency
%       route (supplement).
%   m4  A2'(i.a)-(i.b), A2b and A2c now carry the coordinate index j explicitly.
%
%  ROUND-3 CHANGES (against theorist-critic 45/100).  EVERY ITEM IS TAGGED WITH
%  THE FINDING IT CLOSES so the next reviewer can audit them one by one; the
%  round-2 log tagged only four of its own and that is repaired here.
%  *** NOTE ON TAG NAMESPACES, ADDED IN ROUND 4 (finding m-d).  TAGS IN SQUARE
%  BRACKETS -- [C-2], [M-1], [m-4] -- ARE CRITIC FINDINGS.  Tags in round
%  parentheses -- (C-1)...(C-5) -- WERE COST ITEMS OF THE ROUND-3 PLACEMENT AND
%  ARE A DIFFERENT NAMESPACE WITH DIFFERENT REFERENTS; [C-2] the sparsity error
%  and (C-2) the lost sup-norm bound are not the same object.  THE COST
%  NAMESPACE IS RETIRED WITH THE PLACEMENT IN ROUND 4, so only the bracketed
%  critic tags remain live. ***
%
%   ~~** THE sans-L = 2 GUARANTEE, RESTATED AFTER ROUND 3. **  Every symbol,
%   every assumption clause and every display in this dispatch reduces verbatim
%   at sans-L = 2 WITH TWO EXCEPTIONS, both of them the normalisation: eq:hatT,
%   and the reading eq:psidef is brought to.~~  STRUCK IN ROUND 4.  Round 2
%   claimed one exception and round 3 claimed two; BOTH COUNTS WERE ARTEFACTS OF
%   THE NORMALISATION.  With the raw contrast restored the count is ZERO.  See
%   the round-4 entry below.
%
%   ~~[AUTHOR DECISION, superseding round 2's placement]  BOTH SIDES NORMALISE.
%       psi is read at Wtilde^Omega and so is the operator.~~  SUPERSEDED IN
%       ROUND 4: NEITHER SIDE NORMALISES.  What survives from this entry is the
%       refutation of round 2's bijection -- def:weight's clause (a) is a RANGE
%       condition, ell in R(cal-W), and right-multiplication by a merely Hoelder
%       Omega(z) leaves that range because cal-W's z-sections are real-analytic
%       -- which is kept because it is the same analytic-versus-Hoelder mechanism
%       that makes the author's ground for the reversal bite.  The four
%       bound-checks round 2 ran are each correct and are kept on record as the
%       lesson that bounds are not membership.
%   [C-2]  THE SPARSITY OF Wtilde^Omega WAS WRONG.  "At most two entries nonzero"
%       is false: at W = 0 ALL sans-L - 1 entries equal -1/varrho_0.  The pattern
%       is "one entry, or all of them".  Consequently the Euclidean constant is
%       sqrt(sans-L - 1) eps_varrho^{-1}, SHARP -- not
%       min{sqrt 2, sqrt(sans-L - 1)} eps_varrho^{-1}, which is right at
%       sans-L = 2, 3 and UNDERSTATES AT sans-L = 4, i.e. at Section 6's own
%       design.  Corrected at both sites (this file's def:hatT comment and the
%       A4 note; the supplement's rem:hatTread(2), which was PRINTED).  The
%       min{sqrt 2, .} form for the RAW contrast is correct and stays.
%   [M-1]  THE GROUND FOR THE POINTWISE NORMALISATION WAS A MISQUOTATION.  Round 2
%       cited lf_panel.R:308-313 as saying the alternative "solves an equation
%       whose population zero is not varphi-dagger"; that passage is about
%       CENTRING, not normalisation, and all candidates already centre
%       conditionally.  It also contradicted this file's own statement that a
%       global positive scalar changes no root set.  Struck; the ground is now
%       the three interpretability identities and nothing else.
%   [M-2]  "EFFICIENCY AND CONDITIONING" WITHDRAWN.  No efficiency calculation
%       exists here, and under A1 the efficient conditional weight for the
%       indicator moment is {Omega(z) F(u|z)}^{-1}, which differs from Omega^{-1}
%       by a scalar depending on u and z.  The conditioning claim may run
%       backwards near the overlap boundary, which is where Section 6's n = 66
%       cell sits.  The stated ground is INTERPRETABILITY.
%   [M-3]  THE COST LIST IS FIVE ITEMS AND IS NOT ASSERTED COMPLETE.  (C-1) A4(i)
%       becomes a hypothesis of the operator, reversing M-10 at prp:wellposed,
%       prp:frechet, prp:compact; (C-2) the free bound |Wtilde| <= 1 is lost;
%       (C-3) lem:bierens's clauses naming M or T acquire A4(i) as
%       WELL-FORMEDNESS while (a) <=> (b) does not -- the collision with the C2
%       hand-off is resolved in both files rather than left to Phase 2; (C-4)
%       T-hat's dependence on varrho-hat becomes a RATIO, giving sensitivity
%       eps_varrho^{-2}, curvature eps_varrho^{-3} and a genuine O_P(r_n^2)
%       remainder that was exactly zero before -- a PROOF OBLIGATION at
%       lem:opnoise, not a rate break, since r_n^2 = o(r_n); (C-5) the binary
%       estimator changes and tau_DP must be re-derived.  The five unswept files
%       are named with the list.
%   [M-4]  def:weight ADDED TO THE CENSUS, inside the file round 2 declared
%       swept, with its three required changes -- b's codomain (which ALREADY
%       contradicts A2'(iii) here), cal-W applied columnwise, eq:weightdef's
%       type -- and with the note that its clause (a) is the range condition on
%       which the bijection argument founders.
%   [m-1]  the spanning claim is scoped: dimension exactly sans-L - 1 at a z with
%       every cell non-null, and one less than the number of non-null cells
%       otherwise.
%   [m-2]  norm conventions named where both occur: Frobenius for L, ell, b,
%       ell^Pi; spectral elsewhere; the A4 paragraph says so at its head.
%   [m-3]  type and ordering slips in the A4(d) archaeology: the contrast stands
%       to the RIGHT, "Wtilde^2" names nothing for a vector, ell_star = ell
%       Omega^{-1} and not Omega^{-1} ell.  The rule is stated once at the head
%       of that paragraph rather than re-typing an archaeological record.
%   [m-3b] the range-invariance sentence is SCOPED TO CONSTANT varrho, where it
%       is true; the unqualified "at every sans-L" is withdrawn, on the same
%       ground as the weight-class refutation.
%   [m-4]  the header no longer says Wtilde^Omega "enters eq:hatT and nothing
%       else": psi carries it, and eq:hatT carries the HATTED version.
%   [m-5]  the supplement no longer redefines Z_0 in an unlabelled display; it
%       cites eq:degenreg and fixes the reading.  Amending eq:degenreg is Phase 3.
%   [m-6]  "strictly smaller than {det Omega = 0}" is scoped to the simplex, with
%       the note that for a given DGP the z-sets may coincide.
%   [m-7]  the census cites eq:PsiL (results_s2:481) and not eq:Psitau, which is
%       a different display at proofs.tex:3197.
%   [m-8]  this log.
%   [NEW, found while fixing m-7]  eq:psidef ALREADY EXISTS at
%       results_s3_estimation.tex:216.  Round 3's first draft declared a second
%       \label{eq:psidef} in this file, which would have been a DUPLICATE LABEL
%       resolving silently to whichever came last.  The notation block now prints
%       an UNNUMBERED reading and cites the existing label; eq:psidef is added to
%       the Phase-2 census as the site where psi must be amended.
%
% ---------------------------------------------------------------------------
%  ROUND-4 CHANGES.  THE ESCALATION WAS RESOLVED BY THE AUTHOR AND THE
%  NORMALISATION IS REVERSED.  THE ESTIMATOR CARRIES THE RAW CONTRAST ON BOTH
%  SIDES.  Tags below are critic findings from the round-3 report.
%
%   ** THE sans-L = 2 GUARANTEE, FINAL FORM: **ZERO EXCEPTIONS AS MATHEMATICS**.
%   Every symbol, every assumption clause and every display in this dispatch
%   reduces AT sans-L = 2 TO THE OBJECT IT REPLACED, and the binary estimator is
%   unchanged.  IT IS NOT A CLAIM ABOUT GLYPHS: A0 prints {0,...,sans-L - 1},
%   A4(i) prints a quantifier over l, eq:hatT prints a third argument, and none
%   of those is the old character sequence.  Rounds 4-5 wrote "character for
%   character" at two sites and that OVER-CLAIMS; the claim is about objects.
%   CHECKED ITEM BY ITEM:
%     A0's support {0,1,...,sans-L - 1} is {0,1};
%     eq:contrast's Wtilde is 1{W=1} - varrho = W - varrho, the printed contrast;
%     eq:omegadef's Omega is varrho(1 - varrho), the printed prefactor;
%     A4(i)'s per-cell floor over l in {0,1} IS the printed two-sided interval,
%       and eps_varrho in (0, 1/sans-L) IS (0, 1/2);
%     eq:omegafloor asserts varrho(1-varrho) >= eps/2, true and weaker than the
%       sharp eps(1-eps) -- used only as a positive floor, never at its value;
%     eq:trunset IS the interval [eps, 1 - eps], its two faces being q_1 >= eps
%       and q_1 <= 1 - eps;
%     A4(ii), A4(iii) are one component and one rate;
%     A2'(i.a)-(i.b), A2b, A2c: max over a one-element index set;
%     A2'(iii): R^{p x 1} is R^p, Frobenius is Euclidean;
%     B1(iv.d): Pi_w = delta_1 is forced, so c_Pi does not move;
%     eq:hatT: 1{W_i = 1} = W_i, one contrast component, NO FACTOR -- the
%       pre-dispatch display;
%     eq:psidef: L Wtilde = Wtilde L at a scalar contrast, the printed form.
%   THE ONLY THING A sans-L = 2 READER MEETS THAT IS NEW IS rem:htread, which is
%   an interpretive remark and changes nothing.
%
%   [AUTHOR DECISION]  NORMALISATION REVERSED.  Ground: normalising makes A_cond
%       a DIFFERENT operator, and A2's Picard/range condition, A2', A1c and H1
%       (through sigma_j(A_cond) asymp j^{-a}) are ALL stated at A_cond in this
%       file, with mu = s/(2a), eq:frontier, the feasibility table and the CFF
%       ceiling comparison downstream of a.  This file contains its own proof
%       that the move is not free: A4's note (c) says the two operators must not
%       be assumed to share range and kernel, and round 3 itself scoped
%       range-invariance to CONSTANT varrho.  The gain was interpretational.
%       STRUCK WITH IT: the five-item cost list; "A4(i) becomes a hypothesis of
%       the operator" and the reversal of M-10; the affine-to-ratio second-order
%       remainder at lem:opnoise; "the binary estimator changes"; the
%       two-exception sans-L = 2 guarantee; eq:hatomeganorm, Omega-hat_i and the
%       sample normalised contrast, removed from def:hatT's printed text so that
%       no printed symbol is left without a consumer.
%       RETAINED: the sans-L-category instrument, the contrast VECTOR, Omega as
%       covariance and prefactor, A4(i), eq:trunset, eq:omegafloor, and every
%       sans-L >= 3 correction of rounds 2 and 3.  eq:omeganorm and eq:omegainv
%       survive AS INTERPRETATION at rem:htread, explicitly not as a metric.
%   [CRITICAL, surviving form]  THE TWO MOMENTS OF THE NORMALISED CONTRAST ARE
%       DIFFERENT OBJECTS AND ROUND 3 COMPUTED ONE WHILE CITING THE OTHER:
%         E[Wtilde' Omega^{-1} Wtilde | Delta Z] = sans-L - 1        BOUNDED
%         E[||Omega^{-1} Wtilde||^2 | Delta Z]   = tr(Omega^{-1})    DIVERGENT
%       and tr(Omega^{-1}) = sum_{l>=1} 1/varrho_l + (sans-L - 1)/varrho_0.  At
%       sans-L = 4 with a thin NON-BASE cell at 0.05 and the other three at
%       0.31667 these are 3.00 and 35.79, against eps^{-1} = 20 -- so the second
%       exceeds the level bound by a factor of order sans-L and is NOT controlled
%       by it.  ALSO CORRECTED: E[|Wtilde^Omega_l| | Delta Z] = 2 EXACTLY, FREE
%       of varrho, which is the standard IPW fact -- the mean is untouched and
%       the variance diverges.  Round 3 had this backwards at both ends.  All
%       three are now stated at rem:htread.
%   [CRITICAL, surviving form]  "Wtilde^Omega IS NOT INTEGRABLE WITHOUT A4(i)"
%       IS WRONG and is struck wherever it appeared.  Where a cell vanishes
%       Wtilde^Omega is UNDEFINED, Omega having no inverse; where it is defined
%       it is integrable with a varrho-FREE bound, E|Wtilde^Omega_l| = 2.
%       UNDEFINEDNESS AND NON-INTEGRABILITY ARE DIFFERENT FAILURES.
%   [M-B]  the C-4 Taylor constants were stated sans-L-free and are not:
%       Wtilde-hat depends on ALL sans-L - 1 coordinates through
%       varrho-hat_0 = 1 - 1'varrho-hat, so the gradient carries
%       sqrt(sans-L - 1) eps^{-2} and the curvature (sans-L - 1) eps^{-3}.  MOOT
%       under the reversal -- the map is affine again -- and recorded so that any
%       future ratio-in-varrho-hat statement gets the cardinality right.  The
%       census entry calling eps^{-1} "THE LARGEST OF THE FIVE" while omitting
%       eps^{-2} and eps^{-3} is DELETED with the eps entries.
%   [M-C]  stale counts fixed: the display list here is four items under a
%       heading of FOUR; the two change logs now carry the SAME total (seven);
%       and the notation-decision heading reads FIVE over five items.
%   [m-a]  the supplement's spanning derivation now carries the degenerate-z
%       statement main promises, so the two files agree that spanning does not
%       need A4(i).
%   [m-b]  "Wtilde L_star is ill-formed" scoped to sans-L >= 3; at sans-L = 2 it
%       is the printed eq:psidef and is well formed.
%   [m-c]  the Pi_w constant is max_l varpi_w(l)^{-1/2} = (min_l varpi_w(l))^{-1/2},
%       not "min_l varpi_w(l)^{-1/2}".
%   [m-d]  the two tag namespaces are separated at the head of the round-3 log.
%   ~~[m-f]  "unnumbered" -> "unnormalised" at the census entry for eq:psidef.~~
%       MIS-LOGGED IN ROUND 4 AND CORRECTED IN ROUND 5.  The word "unnumbered" at
%       the eq:psidef census entry is CORRECT -- the notation block does print an
%       unnumbered display -- and the round-4 rewrite of that entry removed the
%       "unnormalised scalar form" wording the finding was actually about, so the
%       fix landed and the LOG entry named the wrong site.  A log that points at
%       the wrong line is worse than no log, because it is read as an audit.
%   [tau_DP]  RECORDED, AND ROUND 5 REPLACES THE FRAMING ENTIRELY.  ~~A7 requires
%       tau_DP >= 1; the Landweber rate requires tau_DP > 2; the
%       engine runs 0.7 against a PLUG-IN noise level roughly 3x the theory's
%       delta-hat_n.  That is a discrepancy between two NOISE-LEVEL conventions,
%       NOT a tau_DP below the theory's floor.  The relationship Phase 6 must
%       satisfy is written out at the rejection block.  The ratio c is not
%       established here.~~
%       ALL OF THAT IS STRUCK.  THE ENGINE RUNS tau_DP = 2.1 > 2, SO THE THEORY'S
%       FLOOR IS SATISFIED AND NOTHING DEVIATES FROM A7.  There are no "two
%       noise-level conventions" and no ratio c: the theory's delta-hat_n is a
%       DETERMINISTIC RATE (A5(i), constant 1) and the engine's res$noise is a
%       DATA-DEPENDENT PLUG-IN, different species, which is why no constant ever
%       reconciled them.  THE RECONCILIATION IS eq:lfsched's c_14, a constant the
%       paper already declares must be CHOSEN because it is not identified by the
%       data; res$noise is an ESTIMATOR OF c_14 delta-hat_n ln n, and c_14 ~ 0.018
%       on the application's numbers.  THE RESIDUAL GAP IS RECORDED AND NOT
%       CLOSED: prp:phirate Step 2 needs c_14 >= c_16 and c_16 is not known to the
%       analyst, so a SMALL c_14 is the more exposed choice.  The calibration is
%       admissible but unverifiable, and that is the THEORY's declared gap rather
%       than one the implementation introduces.  Full statement, with the
%       arithmetic, at the rejection block.
%
%  ROUND-5 CHANGES, TAGGED.
%   [C-1]  FOUR SITES STILL ASSERTED THE REVERSED DECISION AND ARE FIXED: this
%       file's notation-decision item 4 (which stated "BOTH THE OPERATOR AND psi
%       ARE WRITTEN AT Wtilde^Omega ... NOT REOPENED HERE" in the present tense
%       and re-asserted the two-exception guarantee, giving the file two
%       contradictory [AUTHOR DECISION] entries); item 5 beside it; the
%       Phase-2 TYPES handoff, which said "psi = ... L Wtilde^Omega" and so
%       handed Phase 2 the reversed object in the one clause meant to fix its
%       types; the A4 NOTE heading "NOW HALF STALE"; and, in the supplement, the
%       round-3 log entry [C-3].  ROUND 4 EDITED THE ITEM COUNT OF THE DECISION
%       BLOCK AND NOT ITS CONTENT.
%   [M-1]  supp: "what makes the normalised operator of the next paragraph well
%       defined" pointed, in PRINTED text, at an object round 4 had removed.
%   [M-2]  Omega-hat had four printed consumers in the supplement and no printed
%       declaration anywhere, round 4 having removed the declaration with the
%       normalisation and kept the consumers.  Declared at its first consumer.
%   [M-3]  def:hatA's kernel gains the contrast index; it printed k(u,t;x), the
%       wrong arity under B1(iv).  Amended on this file's own stated policy, the
%       one that put eq:hatT's third argument here rather than deferring it.
%   [M-6]  `a` added to the sans-L census, with the two ways of moving A_cond
%       distinguished -- normalising (rejected) against enlarging the range
%       (unavoidable, taken) -- and with the two uncomfortable facts stated: the
%       L^2(Pi) norm AVERAGES over contrasts so more cells do not monotonically
%       improve the singular values, and `a` now depends on Pi_w.
%   [AUTHOR DECISION 2]  "converts A1 from untestable" and the sans-L - 2 count
%       are STRUCK.  A1 is testable at sans-L = 2.  The motivation is
%       LOCALISATION; joint consistency is the apparatus.
%   [m-1] worst case vs the example's 1.79.  [m-2] "character for character" ->
%       a claim about objects.  [m-3] the "Omega^{-1} appears ONLY in rem:htread"
%       claim corrected.  [m-4] the equation-shift understatement corrected.
%       [m-5] this entry.  [m-6] the uniqueness step no longer cites
%       eq:omegafloor.
% ---------------------------------------------------------------------------
%
%  WHAT IS DELIBERATELY NOT DONE, AND EACH IS A HANDOFF AND NOT AN OMISSION.
%   - THE CROSS-FILE INCONSISTENCY IS NOT ONE SITE.  The full list this dispatch
%     found -- twelve sites in results_s2_identification.tex plus def:spaces --
%     is at the comment under def:hatT, WITH the statement that the count is a
%     LOWER BOUND: results_s3_estimation.tex, results_s4_asymptotics.tex,
%     results_s3_supp.tex, results_supp.tex and proofs.tex were not swept.  The
%     round-1 form of this entry said "THAT ONE PLACE", which named two sites and
%     asserted completeness -- worse than naming none, because it stops the next
%     reader looking.
%   - prp:orthog's system in (L_1, ..., L_{sans-L - 1}),
%     lem:adjoint's sum over l, and the NEW localisation proposition are
%     Phase 2.  This file fixes their TYPES -- L, ell, b matrix-valued, ell^Pi
%     four-argument (j, u, t, l), psi = Delta Z Delta U + L(Delta U, Delta Z)
%     Wtilde WITH THE RAW CONTRAST -- in the notation block, and nothing more.
%     *** THIS CLAUSE READ "L Wtilde^Omega" THROUGH ROUND 4 AND WAS THEREFORE
%     HANDING PHASE 2 THE REVERSED OBJECT, in a clause whose whole purpose is to
%     fix the types Phase 2 consumes, and in contradiction with this file's own
%     notation block.  A contract clause naming the wrong object is a first-order
%     defect and it is corrected here. ***
%   - WHAT THIS DISPATCH COULD NOT ESTABLISH, STATED RATHER THAN LEFT IMPLICIT.
%     ROUND 4 REWRITES THIS LIST: items (1), (2) and (5) were about the rejected
%     normalisation and are moot; they are kept in struck form because they are
%     the reasons a re-inventor would offer.
%     ~~(1) that Omega^{-1} is efficient~~ -- IT IS NOT, and no efficiency
%         calculation for this estimator exists in the corpus.  MOOT, and the
%         disproof is retained at the rejection block.
%     ~~(2) that the normalisation improves conditioning~~ -- unestablished and
%         plausibly false near the overlap boundary.  MOOT.
%     (3) THE CROSS-FILE CENSUS IS NOT COMPLETE.  Five files are unswept:
%         results_s3_estimation.tex, results_s4_asymptotics.tex,
%         results_s3_supp.tex, results_supp.tex, proofs.tex.  THIS ONE STANDS AND
%         IS THE LIVE OMISSION OF THE DISPATCH.
%     ~~(4) whether the weight-class bijection can be repaired under a weaker
%         class~~ -- MOOT: nothing needs it once neither side normalises.
%     ~~(5) the sharp constant in lem:opnoise's new second-order term~~ -- MOOT:
%         with the raw contrast the map is affine and there is no such term.
%     (6) THE RATIO c BETWEEN THE CODE'S PLUG-IN NOISE LEVEL AND THE THEORY'S
%         delta-hat_n IS NOT ESTABLISHED HERE.  The tau_DP note above states the
%         relationship Phase 6 must check; it does not verify c ~ 3.
%   - NO CONSTANT IS RE-DERIVED.  Those that carry sans-L, at the sites where
%     each is flagged.  ROUND 4 REMOVES THE eps_varrho ENTRIES, which existed
%     only under the normalisation:
%        sqrt(sans-L - 1)          A4(iii)'s norm equivalence; the Frobenius
%                                  passage from the coordinatewise budgets;
%        min{sqrt 2, sqrt(sans-L - 1)}  the Euclidean bound on Wtilde, = 1 at
%                                  sans-L = 2 and sqrt 2 at sans-L >= 3.  FREE of
%                                  varrho and UNIFORM in sans-L, which is what
%                                  keeps A4(i) idle in prp:wellposed, prp:frechet
%                                  and prp:compact;
%        sans-L                    eps_Omega = eps_varrho / sans-L, eq:omegafloor;
%        a  (THE DEGREE OF ILL-POSEDNESS)  -- ADDED IN ROUND 5, AND ITS ABSENCE
%                                  WAS THE MOST CONSEQUENTIAL OMISSION IN THIS
%                                  CENSUS, because the reversal's own ground is
%                                  that A_cond must not move.  *** THIS DISPATCH
%                                  MOVES A_cond TOO, BY A DIFFERENT ROUTE, AND
%                                  THE TWO ROUTES MUST BE DISTINGUISHED. ***
%                                  (i) BY NORMALISING -- REJECTED.  That
%                                  multiplies the integrand by a z-varying
%                                  Omega^{-1}, and the analytic-versus-Hoelder
%                                  mechanism means range and kernel need not
%                                  transfer at all; A4's note (c) says so.
%                                  NOTHING about A2, A2', A1c or H1 would carry
%                                  over, and that is why it was declined.
%                                  (ii) BY ENLARGING THE RANGE -- UNAVOIDABLE AND
%                                  TAKEN.  B1(iv) makes Pi = Pi_u (x) Pi_t (x)
%                                  Pi_w, so A_cond maps into a LARGER space and
%                                  its kernel carries l (def:hatA, as amended).
%                                  At sans-L >= 3 the singular values sigma_j are
%                                  those of a DIFFERENT operator from the binary
%                                  one, so `a` MAY be a different number -- not
%                                  IS, which was asserted where only "may" is
%                                  supported, and which the Pi_w-invariance
%                                  above further narrows.  Phase 2 is to settle
%                                  whether the per-cell operators share a decay
%                                  exponent: if they do, `a` is unchanged; if
%                                  they do not, a = min_l a_l and the exponent
%                                  can only weakly IMPROVE, making the budget
%                                  table's reading conservative at sans-L = 4
%                                  rather than merely unverified.  And
%                                  mu_H1 = s/(2a), eq:frontier's readings and the
%                                  budget-table configurations are read at it.
%                                  NOTHING IN THE FILE IS FALSE ON THIS ACCOUNT
%                                  -- every statement is made at "A_cond",
%                                  whatever sans-L is -- BUT THE NUMBERS IN THE
%                                  BUDGET TABLE WERE CALIBRATED AT THE BINARY
%                                  OPERATOR AND ARE NOT AUTOMATICALLY THE RIGHT
%                                  NUMBERS AT sans-L = 4.
%                                  TWO FURTHER FACTS, BOTH UNCOMFORTABLE AND BOTH
%                                  STATED RATHER THAN LEFT: (1) because Pi_w is a
%                                  PROBABILITY measure, ||A_cond h||^2 =
%                                  sum_l varpi_w(l) ||A^{(l)} h||^2 is an AVERAGE
%                                  over contrasts, NOT a sum, so adding a cell
%                                  with little signal can DECREASE it -- more
%                                  cells do not monotonically improve the
%                                  singular values and no sentence may say they
%                                  do; (2) ~~`a` therefore DEPENDS ON Pi_w, an
%                                  analyst's choice, which is a dependence the
%                                  binary block never had.~~ -- STRUCK, AND IT
%                                  DOES NOT FOLLOW.  Under B1(iv.d) every
%                                  admissible Pi_w has strictly positive mass on
%                                  a FINITE index set, so for two admissible
%                                  masses the range norms are EQUIVALENT: with
%                                  c = min_l pi'(l)/pi(l) > 0 and C = max_l the
%                                  same ratio, c||Ah||^2_pi <= ||Ah||^2_pi' <=
%                                  C||Ah||^2_pi for every h.  Singular values of
%                                  a compact operator are its approximation
%                                  numbers a_j(T) = inf{||T - F|| : rank F < j},
%                                  so sqrt(c) sigma_j^pi <= sigma_j^pi' <=
%                                  sqrt(C) sigma_j^pi with J-INDEPENDENT
%                                  constants, and H1 defines `a` only through
%                                  sigma_j ~ j^{-a}.  So `a` IS INVARIANT to
%                                  Pi_w.  What Pi_w moves is the CONSTANT in the
%                                  ~, which the census entry
%                                  max_l varpi_w(l)^{-1/2} already prices.  At sans-L = 2,
%                                  Pi_w = delta_1 and A_cond IS the binary
%                                  operator, so `a` is the binary `a`.
%                                  RE-READING THE BUDGET TABLE AT sans-L >= 3 IS
%                                  NOT DONE HERE AND IS FLAGGED FOR PHASE 2.
%        max_l varpi_w(l)^{-1/2}   the passage between the L^2(Pi) norm and the
%                                  per-contrast L^2(Pi_u (x) Pi_t) norms, since
%                                  ||T||^2_{L^2(Pi)} = sum_l varpi_w(l) ||T_l||^2
%                                  and a bound stated per contrast costs that
%                                  factor when read in the product norm.  IT IS
%                                  max_l varpi_w(l)^{-1/2} = (min_l varpi_w(l))^{-1/2}
%                                  AND ROUNDS 2-3 WROTE "min_l varpi_w(l)^{-1/2}",
%                                  WHICH IS THE RECIPROCAL OF THE LARGEST MASS AND
%                                  THEREFORE THE WRONG END.  Invisible at a uniform
%                                  Pi_w, where every mass is equal and the two
%                                  agree at sqrt(sans-L - 1); WRONG at the
%                                  non-uniform Pi_w that B1(iv.d) permits, which is
%                                  exactly the case a reader would consult it for.
%                                  B1(iv.d)'s strict positivity is what keeps it
%                                  finite.
%     THE PLAN'S "Rates" ITEM ASKS WHETHER ANY CONSTANT IS CLAIMED sans-L-FREE
%     WHEN IT IS NOT; this file flags the ones it can see and the sweep of
%     proofs.tex is Phase 3's.
%
%  FIVE NOTATION DECISIONS THE AUTHOR MAY OVERTURN, ALL RECORDED AT THEIR
%  SITES AND NONE OF THEM SILENT.
%   1. sans-L, NOT ITALIC L, for the category count.  The plan writes italic L;
%      italic L is the debiasing weight, L_x and L_u the sieve dimensions and
%      L^2 the spaces, so "L - 1" would be ambiguous three ways.  One collision
%      survives, sans-L_Gamma in lem:bierenscore's proof (Phase-3 file).
%   2. Omega FOR THE CONTRAST COVARIANCE, AS THE PLAN DIRECTS.  IT WAS a
%      declared case pair against the intertwining family
%      Omega(varphi, varphi-tilde) of prp:norangeinv, separated BY TYPE ALONE --
%      the round-1 claim of arity separation was false, the family having been
%      bound to bare Omega with a bare Omega^{-1} printed.  *** CLOSED 2026-08-09:
%      the family is renamed cal-O and declared inside prp:norangeinv; the nine
%      printed lines that recorded the pair and gave the reader a disambiguation
%      rule are DELETED from the type-declaration block above.  Omega now names
%      one object.  The glossary census certified cal-O and cal-Y as the only
%      free calligraphic letters. ***
%   3. THE BUDGETS AGGREGATE EUCLIDEAN OVER j AND BY THE MAXIMUM OVER l.
%      *** THIS ENTRY READ "over j AND l by the maximum" THROUGH ROUND 1 AND WAS
%      SUPERSEDED BY THE ROUND-2 REPAIR at A2'(i.a)-(i.b), A2b and eq:lambdamu,
%      which carry the p coordinates of theta inside a EUCLIDEAN norm; a
%      coordinatewise maximum in j would leave the downstream constants short. ***
%      Over l, sum and Pi_w-weighted sum were the alternatives; the maximum is
%      conservative and free of Pi_w.
%   4. NEITHER THE OPERATOR NOR psi IS NORMALISED.  BOTH CARRY THE RAW CONTRAST
%      Wtilde.  *** THIS ENTRY IS REWRITTEN IN ROUND 5 AND ITS PREVIOUS CONTENT
%      WAS THE REVERSED DECISION, LEFT STANDING IN THE PRESENT TENSE. ***  It read
%      "BOTH THE OPERATOR AND psi ARE WRITTEN AT Wtilde^Omega ... THIS IS THE
%      AUTHOR'S DECISION AND IS NOT REOPENED HERE", and it re-asserted the
%      two-exception sans-L = 2 guarantee by saying eq:hatT "reproduces a division
%      by varrho(1 - varrho), pointwise in i, which the pre-dispatch display did
%      not carry".  THAT IS THE EXACT NEGATION OF THE HEADLINE CLAIM OF ROUND 4,
%      AND IT SAT IN THE FILE'S DECISION REGISTER, so the file carried two
%      contradictory [AUTHOR DECISION] entries.  Round 4 edited this block's item
%      COUNT and not its CONTENT -- the failure mode this corpus calls a rename
%      reaching the statements and not the prose.
%      WHAT IS TRUE: the raw contrast on both sides; eq:hatT at sans-L = 2 IS the
%      pre-dispatch display; the sans-L = 2 reduction has ZERO exceptions.  The
%      author's ground for declining the normalisation is that it moves A_cond,
%      at which A2, A2', A1c and H1 are all stated.  See the rejection block under
%      the notation paragraph.
%   5. WHAT IS NOT CLAIMED ABOUT THE NORMALISATION, KEPT BECAUSE THE CLAIMS WERE
%      MADE AND WOULD BE MADE AGAIN.  It is NOT efficient, NOT shown to improve
%      conditioning, and changes NO root set.  Since round 4 it is also not USED:
%      these are properties of a device that lives only in rem:htread.  What that
%      remark rests on is interpretability -- the conditional forward residual is
%      the cell difference, the contrast is the Horvitz-Thompson one, the
%      quadratic form is base-cell-free -- and that is all the file says for it.
% ============================================================================
